\section{Trajectories, marked fields and ultraviolet convergence}\label{sec:partII-E}

This section develops the trajectories, renormalized insertions, volume-uniform bounds and full fixed-volume limits of Theorem~\ref{prop:MT}, Proposition~\ref{prop:UB} and Theorem~\ref{prop:UV}. The proofs of the bounds and limits still require the higher-mode attachment and stopping estimates, the actual marked Gaussian comparison, and the fixed-level cutoff comparison identified at their points of use below. The section also develops the cutoff equation identities, coefficient corrections and finite-rank comparisons from which Section~\ref{sec:partII-F} builds the quotient fields. All these results concern the Wilson measures of Definition~\ref{def:D02} and the insertions of Definition~\ref{def:D08}.

Three different tasks are involved, and each has its own mechanism. A continuous shooting argument chooses the bare coupling so that the trajectory reaches one prescribed coarse coupling (Section~\ref{sec:partII-E.1}). Weighted estimates for the specified recursion give uniqueness and analytic dependence near that endpoint. Finally, at fixed volume, convergence of the effective coordinates at each fixed level is combined with a bound on the error made by stopping the insertion transfer there (Section~\ref{sec:partII-E.3}). The additive scalar normalization cancels from normalized expectations, so its convergence is never needed.

The main difficulty in the bounds of Section~\ref{sec:partII-E.2} is differentiation in position. Each lattice difference introduces an inverse spacing, so a bound for undifferentiated insertions does not suffice. Local comparison replaces a differentiated insertion by one of higher degree, with a negligible chessboard error. Exact equation identities remove its equation component. A finite matrix of free Maxwell correlation tests bounds the remaining coefficient. This argument takes place at finite cutoff, before any continuum field space has been constructed.

Throughout, the group, metric, penalty, scales and admissible classes are those of Section~\ref{sec:renormalization} and Section~\ref{sec:constants}. The trajectories are built from the path actions of Lemmas~\ref{lem:C1}--\ref{lem:C2}, which fix the classical and one-loop terms, the Wilson initialization of Lemma~\ref{lem:D3}, and the inductive step of Theorem~\ref{prop:RG}, which carries an admissible representation one level down. Additional dependence of constants is indicated. Write $M_{N,p}=\latmoment_{N,p}$ and $\langle F\rangle_{N,p}=\int F\,d\mu_{N,p}$.

For an insertion label $v=(d,c)$ put $O_{N,p}(v,y)=O^{(d)}_{N,p}(c,y)$. Physical points $x$ in a moment are read as lattice sites $x/a_N$, modulo $L_pD_N$. We use the torus separation
\begin{equation*}
\operatorname{sep}_{\rm tor}(y)=\min_{i\ne j}\min_{z\in L_p\Z^4}
|a_N(y_i-y_j)+z|_\infty,
\end{equation*}
with value $+\infty$ for at most one slot. Multilinear interpolation on the physical grid means the affine interpolation in each coordinate of each grid cube, using its vertex values.

\subsection{Trajectories and tuning}\label{sec:partII-E.1}

The entry levels allow the ratios and stage constants to increase sufficiently slowly. Besides the conditions of Section~\ref{sec:constants}, we use the following quantitative choices. Let $c_{{\rm dir},j}$ be the direct clipping coefficient at stage~$j$, and let $\nu_G^{\rm mat}$ be the real dimension of the ambient matrix space. Put
\begin{equation}\label{eq:E-vacuum-majorant}
Q_j(b)=2(2+4A_j^{\rm path})+1+6(c_{{\rm dir},j}+(A_j^{\rm path})^2)
 +4\nu_G^{\rm mat}(\log16+\tfrac12\log b),\qquad q_m=Q_{J_m}(\widehat\beta_m).
\end{equation}
Here $A_j^{\rm path}$ is the path bound of Lemma~\ref{lem:C1}. For an admissible unmarked representation $\rho_m$ on $V$ sites write $\Phi_m=v_{{\rm sc},m}V-\pathaction_{N,m}-\sum_i\mathbf1_{X_i\,\mathrm{healthy}}E_i^{\rm reg}$ and $K_{m,Y}$ for its forced activities. The next estimate bounds, per site, the ratio obtained by replacing the positive density by absolute values in a normalized integral.
\begin{lemma}[Normalization cost]\label{lem:E-normalization-cost}
For every sufficiently late level $m$ and every admissible unmarked representation $\rho_m$ on $V$ sites,
\begin{equation}\label{eq:E-vacuum-ratio}
 \frac{\int e^{\Phi_m(U)}\sum_{\mathcal A}\prod_{Y\in\mathcal A}|K_{m,Y}(U)|\,dU}
 {\int\density\rho_m\,dU}\le e^{Vq_m}.
\end{equation}
\end{lemma}
\begin{proof}
Put $\widehat g_m=\widehat\beta_m^{-1/2}$. The scalar-free unmarked exponent is at most $VE_m$, and its absolute forced sum is at most $e^{V\epsilon_m^{\rm for}}$ with $\epsilon_m^{\rm for}\le1$; at every sufficiently late level the path bounds give $E_m\le2+4A_{J_m}^{\rm path}$. Restrict the denominator to links of chord at most $\widehat g_m/4$. There are no flags there, the forced sum is one, and its scalar-free exponent is at least $-V(E_m+6(c_{{\rm dir},J_m}+(A_{J_m}^{\rm path})^2))$. The elementary Haar-ball bound gives this set measure at least $(\widehat g_m/16)^{4V\nu_G^{\rm mat}}$. The common scalar normalization in numerator and denominator cancels, so the ratio is at most
\begin{equation*}
 \exp\!\left\{V\left[2E_m+\epsilon_m^{\rm for}
 +6(c_{{\rm dir},J_m}+(A_{J_m}^{\rm path})^2)
 +4\nu_G^{\rm mat}\log(16/\widehat g_m)\right]\right\}\le e^{Vq_m}.
\qedhere
\end{equation*}
\end{proof}
The per-site exponent $q_m=O(1+\log\widehat\beta_m)$ comes from the Haar-ball loss $\widehat g_m^{-\nu_G^{\rm mat}}$ on every link. Lemma~\ref{lem:E-source-integral} improves it to a constant depending only on the stage: in an axial gauge the denominator contributes Haar-ball factors on $3V+1$ links and the numerator contributes one-link integrals on $3V-3$ links, so only a loss of order $\log\widehat\beta_k$ for the whole torus remains. The source comparison needs that improvement, because its gain per cell is only a power of $\cutoffsize_t$. Enlarge $Q_j$ at the finitely many initial levels, so that $q_m\ge1$ and~\eqref{eq:E-vacuum-ratio} hold at every level. The choices include
\begin{equation}\label{eq:E-marker-scales}
 B_{i+1}^2\le2B_i^2,\qquad (\ell_iB_i^2)^4q_{i+1}=o(s_{i+1}),\qquad
 \varpi_{J_i}=o(s_i).
\end{equation}
Define $k(m)$ to be the first $k\ge m$ with $B_k^2\le a_m/a_k$.

The short-distance comparison uses the following comparison windows of levels. For $b>0$ and a level~$t$ in stage~$j$, write $L=\ell_j^{\rm st}$ and
\begin{equation}\label{eq:E-window}
 h_t(b)=\left\lceil\frac{(b+2)\log\cutoffsize_t}{\log L}\right\rceil,
 \quad v_t=\left\lceil\frac{10\log B_t}{\log L}\right\rceil,
 \quad k_t=t-h_t(b)-v_t.
\end{equation}

The numerical requirements of the marked proofs serve three purposes: contraction of one-insertion remainders (items~\ref{it:E-columns}--\ref{it:E-match}), control of replicated sources and comparison with the free field (items~\ref{it:E-log}--\ref{it:E-entries}), and the initialization of insertions (item~\ref{it:E-birth}). At each stage the finite-dimensional and support constants are fixed first, then the block ratio, and finally the coupling entry. The constants named here are introduced with their estimates; we list them so that the order of choice is visible in one place.

For each stage~$j$, let $K_j\ge1$ be a nondecreasing bound for the leading coefficient maps and their inverses through that stage's finitely many degrees and symmetry averages. These bounds are fixed before the column bounds~$C^{\rm col}_j$. Let $A_d^{\rm sym}$ bound the finite basis recombination in degree~$d$, and let $C_j^{\rm loc}$ bound the local marked error of the one-step estimate before the ratio is chosen (proof of Lemma~\ref{lem:E-marked-step}). Each step below has $j=J_{m+1}$ and $L=\ell_m$; in items~\ref{it:E-graph}--\ref{it:E-match} and in the proofs of Lemmas~\ref{lem:E-marked-step}, \ref{lem:E-local-comparison} and~\ref{lem:E-coarse-bounded} we write $K=K_j$, $S=C^{\rm col}_j$ and $\varepsilon_L=L^{-3/8}$. For the retained incoming regular term, and for the unsubtracted incoming contribution in Section~\ref{sec:partII-J}, let $R_j^{\rm supp}\in\N$ bound the enlargement of its projected incoming read support in units of $B_m$, over the finitely many adjacent stage pairs with fine stage~$j$. Put
\begin{equation}\label{eq:E-support-conversion}
 \mathcal A_m=\exp\!\left\{(\sigma_{J_m}+2)(2R_{J_{m+1}}^{\rm supp}+13)^4\right\}.
\end{equation}
This factor converts an incoming support weight into the outgoing weight. Indeed, for fine and coarse covering numbers $s_f,s_c$, the connected-cover bound gives
\begin{align*}
 Ls_c&\le(L+10s_f)(2R_j^{\rm supp}+13)^4,\\
 (\sigma_{J_m}+2)s_c
 &\le\log\mathcal A_m+\tfrac12\sigma_{J_{m+1}}s_f,
\end{align*}
provided $20(\sigma_{J_m}+2)(2R_j^{\rm supp}+13)^4\le L\sigma_{J_{m+1}}$. Thus half the incoming exponent remains for the other support factors. The constant $\mathcal C_j\ge2$ includes these bounds and the finite source comparison constants. In item~\ref{it:E-birth}, $C_d$ is the section bound in~\eqref{eq:A-global-bound}.
For the source estimates, choose an integer $q^{\rm ht}_d$ bounding the powers of the inverse health threshold in degree~$d$, and put
\begin{equation*}
 \zeta_{d,i}=\widehat\beta_i^{-(q^{\rm ht}_d+4)}.
\end{equation*}
This bound controls attachments of other marked factors to a distinguished mark. Let $A_{j,d}$ bound the finite section, derivative and support-conversion constants in these attachments at stage~$j$, and let $\eta^{\rm att}_i>0$ be the attachment threshold in Lemma~\ref{lem:D-connected} for the step at level~$i$. Decrease it so that $\eta^{\rm att}_i\le1/4$. All these constants are fixed before the coupling entry. Lemma~\ref{lem:E-source-attachments} uses
\begin{equation}\label{eq:E-attachment-allowance}
 2A_{j,d}(2B_i+1)^4(1+B_i)^4
 \widehat\beta_i^{q^{\rm ht}_d}\zeta_{d,i}\le\eta^{\rm att}_i/4,
 \qquad j=J_{i+1},\quad d\le J_i.
\end{equation}
The factor $(1+B_i)^4$ bounds the spatial sum from marked anchors to a site; it is derived in the proof of that lemma. For transport to finer levels we also require
\begin{equation}\label{eq:E-source-allowance-transport}
 \widehat\beta_{i+1}^{-1}\ell_i^d C^{\rm col}_{J_{i+1}}\zeta_{d,i}
 \le\zeta_{d,i+1}.
\end{equation}
\begin{enumerate}
\item\label{it:E-columns} \emph{Coefficient and support bounds.} Choose nondecreasing columns with $C^{\rm col}_j\ge64K_{j+1}^2$, above the derivative factorials and finite basis-recombination constants at stage~$j$. The support exponent is $\sigma_j=100$; we also retain the estimates at every exponent in $[98,100]$.
\item\label{it:E-graph} \emph{One-insertion contraction} (Lemma~\ref{lem:E-marked-step}). Choose $L^{3/5}>2C_j^{\rm loc}$, $\mathcal A_m\le L$ and $20(\sigma_{J_m}+2)(2R_j^{\rm supp}+13)^4\le L\sigma_{J_{m+1}}$. The local marked error $\epsilon^{\rm mk}_m$ satisfies $\epsilon^{\rm mk}_m\le L^{-2/5}$, and
\begin{equation}\label{eq:E-graph-reserve}
 \epsilon^{\rm mk}_m(1+S)A_d^{\rm sym}\le\varepsilon_L,\quad
 2K\varepsilon_L(1+S)\le1,\quad4S\varepsilon_L\le1,\quad
 2K\le S,\quad4K^2\le C^{\rm col}_{J_m}.
\end{equation}
\item\label{it:E-match} \emph{Backward matching} (Lemmas~\ref{lem:E-local-comparison} and~\ref{lem:E-coarse-bounded}). Require
\begin{equation}\label{eq:E-matching-choices}
 4K\,32(1+S)^2\varepsilon_L\le1,\qquad
 (4K^2+1)L^{1/8}\varepsilon_L\le1.
\end{equation}
For each auxiliary section through the stage's degrees, require also $2(1+C^{\rm col}_j)H_{j,d}A_d^{\rm sym}L^{-1/8}\le1/4$, where $H_{j,d}$ bounds its normalized healthy derivatives.
\item\label{it:E-log} \emph{Accumulated transport losses.} Choose ratios to satisfy~\eqref{eq:E-log-cost}, including every constant in~$\mathcal C_j$.
\item\label{it:E-entries} \emph{Gaussian comparison and source bounds} (Lemmas~\ref{lem:E-source-polynomial}--\ref{lem:E-source-integral}). On a comparison window, for an arity bound $n\ge2$, put $\mathsf w_t=(\cutoffsize_t)^{1/(8n)}$. Let $\mathcal E_t^{\rm G}$ be the sum of the covariance, rounding, Taylor and localization errors in the fixed finite source comparison. Require
\begin{equation}\label{eq:E-Gaussian-entry}
 \mathsf w_t^n\mathcal C_j^{h_t(b)+v_t+2}\mathcal E_t^{\rm G}\le1,
 \qquad Q_j^{\rm cmp}\le\log\cutoffsize_t/(16n),
 \qquad \widehat\beta_i\in[\widehat\beta_t/2,\widehat\beta_t]\quad(k_t\le i\le t).
\end{equation}
Here $Q_j^{\rm cmp}$ is the stage normalization constant in~\eqref{eq:E-comparison-cost}, enlarged to cover the source modes and amplitudes of the finitely many arity bounds $n\ge2$ used at stage~$j$, and the second bound is required for each of these~$n$. These constants are fixed first; the entry is then raised until $\log\cutoffsize_t\ge16n\,Q_j^{\rm cmp}$. The first bound is imposed with the stronger smallness used in~\eqref{eq:E-source-iteration}. The entry also meets the local marked bound in item~\ref{it:E-graph}, the forced-source bounds at every exponent between $98$ and $100$, and~\eqref{eq:E-attachment-allowance}--\eqref{eq:E-source-allowance-transport}. The unmarked attachment sum in Lemma~\ref{lem:D-connected} is at most $\eta^{\rm att}_i/2$, and $64\log\widehat\beta_{i+1}/B_i\le1/4$.
\item\label{it:E-birth} \emph{Initialization of insertions} (Lemma~\ref{lem:E2.1}). At every level~$m$ and $d\le J_m$, require $B_m$ at least the radius of~$\Sigma_d$ and $C_dt_m^{-d}e^{\sigma_{J_m}}(e^{|\Sigma_d|\epsilon_m^{\rm for}}-1)\le(\epsilon_m^{\rm for})^{1/2}$.
\item\label{it:E-ratio-log} \emph{Ratio below the logarithmic coupling} (Lemma~\ref{lem:I-radii}). At every level~$m$, $\ell_m\le\log\widehat\beta_m$.
\end{enumerate}
The finite comparisons in item~\ref{it:E-entries} exhaust bounded configurations with positive point separation, fixed arities and finite derivative orders as~$j$ increases. Their errors and constants are calculated in Lemma~\ref{lem:E-Gaussian-sources}; none depends on the period.

The following choice gives both the sparse insertion grids and the windows used for free comparison.
\begin{lemma}[Scale and period choices]\label{lem:E1.1}
All requirements~\ref{it:E-columns}--\ref{it:E-ratio-log} and~\eqref{eq:E-marker-scales} can be met with
\begin{gather*}
\ell_m\ge3\text{ odd},\qquad \log\ell_m\to\infty,
\qquad \log\ell_m/s_m\to0,\qquad4\mid L_p,\\
(a_m/a_{k(m)})^4q_{k(m)}=o(s_{k(m)}).
\end{gather*}
The sides $L_p$ increase strictly; every positive integer divides infinitely many. For every finite set of exponents $b\ge1$ and finite comparison requests, arbitrarily late terminal levels~$t$ meet the comparison requirements simultaneously on their windows~\eqref{eq:E-window}, with $k_t\to\infty$ and $J_i=J_t$ throughout each window.
\end{lemma}
\begin{proof}
The choices follow the order of Table~\ref{tab:constant-choices}: stage constants, then ratios, then entries. First fix the entire sequence $K_j$ and the finite section bounds. Choose $C^{\rm col}_j\ge64K_{j+1}^2$, and then fix $\mathcal C_j$ to include these columns and all source constants. Since $J_{m+1}\le J_m+1$, this gives $4K_{J_{m+1}}^2\le C^{\rm col}_{J_m}$ even at a stage transition. Next choose nondecreasing odd stage ratios $\ell_j^{\rm st}\ge\max(3,j)$ large enough that
\begin{equation}\label{eq:E-log-cost}
\frac{\omega_j^{\rm RG}}{\log\ell_j^{\rm st}}
\le\min(\slope/4,j^{-2}),\qquad
\frac{\log\mathcal C_j}{\log\ell_j^{\rm st}}\le j^{-2}.
\end{equation}
Each denominator can be increased after its numerator is fixed. The ratio requirements in items~\ref{it:E-graph} and~\ref{it:E-match} are lower bounds on this same later choice: $L^{-2/5}/L^{-3/8}=L^{-1/40}\to0$ and $L^{1/8}L^{-3/8}=L^{-1/4}\to0$. Choose also $L^{3/5}>2C_j^{\rm loc}$, so $C_j^{\rm loc}/L<\tfrac12L^{-2/5}$. For this fixed $L$, raise the coupling entry until the absolute value of the $o_{\widehat\beta_m}(1)$ term in the one-step estimate is at most $\tfrac12L^{-2/5}$. The local error $C_j^{\rm loc}/L+o_{\widehat\beta_m}(1)$ is then at most $L^{-2/5}$. The column and ratio choices also give the fixed auxiliary-section inequality. The support radii and exponents in~\eqref{eq:E-support-conversion} are fixed before~$L$, so its compression condition and $\mathcal A_m\le L$ are further fixed lower bounds on the ratio. Thus no column is chosen in response to a ratio or to a later-stage column.

Choose consecutive constant-stage intervals, postponing each new stage until all its fixed coupling requirements hold. On a fixed stage $s_i$ grows linearly, while $B_i=O((\log(2+s_i))^2)$ and $q_i=O(1+\log(2+s_i))$. Thus
\begin{equation*}
 (\ell_iB_i^2)^4q_{i+1}/s_{i+1}\longrightarrow0
 \quad\text{on that stage as its entry is postponed}.
\end{equation*}
For the initial stage, choose $\bar\beta$ after its fixed ratio~$L_0$. For $B(b)=\lceil(\log b)^2\rceil$, the quotient $B(b+\slope\log L_0)^2/B(b)^2$ tends to one as $b\to\infty$. Choose $\bar\beta$ beyond a threshold where it is at most two for every $b\ge\bar\beta$. This covers every step of the first stage, including $i=0$. At a transition to a later ratio, impose the same bound at its entry and thereafter, with that ratio fixed. Impose the second condition of~\eqref{eq:E-marker-scales} with bound $1/j$ after entry to stage~$j$. Impose also $\varpi_j\le s_i/j$ and $\log\ell_j^{\rm st}\le s_i/j$ at its entry. These are lower bounds on that entry. By item~9 of Table~\ref{tab:constant-choices}, $t_m^{-1}\le\widehat\beta_m$ and $(\epsilon_m^{\rm for})^{1/2}\widehat\beta_m^{d+1}\to0$, so item~\ref{it:E-birth}, including the finite stencil-radius requirement, is another entry bound. Since $\log\widehat\beta_m\to\infty$ at a fixed ratio, item~\ref{it:E-ratio-log} is also an entry bound, imposed for both ratios that occur at the levels of a stage. At each fixed ratio, the same threshold formula gives $t_m/t_{m+1}\to1$ as $\widehat\beta_m\to\infty$; impose $\tfrac12\le t_m/t_{m+1}\le2$ at entry and thereafter, with $t_m\le1$ supplied by item~7 of Table~\ref{tab:constant-choices}. The preceding entry restrictions give both little-oh assertions and $\log\ell_i/s_i\to0$; $s_i\ge i\log3$ at every level.

For each stage, the length $h_t(b)+v_t$ in~\eqref{eq:E-window} is $O_j((\log\log(2+s_t))^2+\log\log(2+s_t))=o(t)$. Lengthen that constant-stage interval to contain a terminal point with the whole window~\eqref{eq:E-window} in the same stage and with its left endpoint beyond all earlier requested levels. Do this with the longest window $b=j$, all arities $n\le j$ and the finitely many comparison orders at stage~$j$, on the compact sets of test configurations with point separations at least $1/j$ and coordinates of absolute value at most~$j$. Every $1\le b\le j$ uses a subinterval of this window. The length and error bounds are uniform for $b$ in any fixed compact subinterval of $[1,j]$; take the maximum of the finite comparison constants before postponing the common terminal level. Their finitely many derivative bounds are uniform on each such compact set. The estimates~\eqref{eq:E-source-loss}--\eqref{eq:E-source-iteration} show that the first requirement in~\eqref{eq:E-Gaussian-entry} holds once that entry is late enough. The second is a lower bound on $\log\cutoffsize_t$ after $Q_j^{\rm cmp}$ is fixed. The total window length is $o(t)$, so its coupling variation is $o(\widehat\beta_t)$ and gives the third requirement. These are lower bounds on a later entry or interval length. Each fixed request is then met at arbitrarily large stages. The source requirements are also entry bounds. On a fixed stage, the left side of~\eqref{eq:E-attachment-allowance} is a fixed constant times $(2B_i+1)^4(1+B_i)^4\widehat\beta_i^{-4}$ and tends to zero. In~\eqref{eq:E-source-allowance-transport}, division by $\zeta_{d,i+1}$ leaves
\begin{equation*}
 \ell_i^d C^{\rm col}_{J_{i+1}}\widehat\beta_{i+1}^{-1}
 (\widehat\beta_{i+1}/\widehat\beta_i)^{q^{\rm ht}_d+4},
\end{equation*}
which tends to zero at fixed ratio. Include both adjacent stages at a transition. There are finitely many degrees at each entry. The unmarked attachment sum tends to zero by the factor rates in the proof of Lemma~\ref{lem:D-connected}, so its half-threshold requirement is another entry bound. These choices do not alter earlier intervals. Taking $\ell_i=\ell^{\rm st}_{J_{i+1}}$ gives odd ratios, and $\ell_j^{\rm st}\ge j$ gives $\log\ell_i\to\infty$.

The level $k(m)$ exists: $a_iB_i^2$ decreases by at least the factor $2/3$ at each step. If $k(m)=u+1>m$, minimality gives
\begin{equation}\label{eq:E-first-open}
 a_m/a_{k(m)}\le\ell_uB_u^2.
\end{equation}
Consequently~\eqref{eq:E-marker-scales} gives $(a_m/a_{k(m)})^4q_{k(m)}=o(s_{k(m)})$. The case $k(m)=m$ satisfies the same conclusion directly from $q_m=o(s_m)$.

Choose $L_p=4(M_{\rm per}+p+1)$. For $j=J_{m+1}$, the required period inequalities, including item~7 of Table~\ref{tab:constant-choices}, are
\begin{gather*}
 L_p e^{s_m}\ge100(B_m+1),\qquad
 L_p e^{s_{m+1}}\ge100((j+3)\ell_m+1),\\
 L_p e^{s_m}\ge v_j^{\rm per}(1+B_m+B_{m+1})^{h_j^{\rm per}}.
\end{gather*}
The first follows by choosing $M_{\rm per}$ large, since $(B_m+1)e^{-s_m}$ is bounded. After canceling $\ell_m$ in the second, postpone each stage until $e^{s_m}\ge100(j+4)$. For the third, $v_j^{\rm per}$ and $h_j^{\rm per}$ are fixed before entry to stage~$j$, and $v_j^{\rm per}(1+B_m+B_{m+1})^{h_j^{\rm per}}e^{-s_m}\to0$ on that stage. Postpone each later stage until this expression is at most one thereafter, and enlarge $M_{\rm per}$ for the finite initial interval. This proves all three inequalities. The third supplies the period bound for the five-torus localization in the proof of Lemma~\ref{lem:D-localization}. The periods are strictly increasing, divisible by four, and every integer divides infinitely many of them.
\end{proof}

The next identities retain the normalization of the original Wilson integral through every integration.
\begin{lemma}[Exact trajectory identities]\label{lem:E1.2}
For the trajectories constructed below with endpoint $\bar\beta$ and every $m<N$,
\begin{gather*}
\rho_{N,p,N}=\rho^{\rm W}_{\beta_{N,p}},\qquad
|\beta_{N,p,m+1}-\beta_{N,p,m}-\slope\log\ell_m|
\le\omega_{J_{m+1}}^{\rm RG},\\
\int\density\rho_{N,p,m}=Z_{L_pD_N,\beta_{N,p}}.
\end{gather*}
All representations are admissible, and the terminal integral has the same value.
\end{lemma}
The following proof also proves Theorem~\ref{prop:MT}.
\begin{proof}[Proof of Lemma~\ref{lem:E1.2} and Theorem~\ref{prop:MT}]
The trajectory is obtained by bracketing, level by level, along the continuous Wilson curve of Lemma~\ref{lem:D3}, carried down by the continuous step of the specified recursion in Lemma~\ref{lem:D2}; the drift margin comes from~\eqref{eq:E-log-cost}. Set
\begin{equation*}
b_m^- =\bar\beta+\tfrac12\slope s_m,\qquad
b_m^+ =\bar\beta+2\slope s_m.
\end{equation*}
These intervals lie in the coupling bands $\mathfrak B_m$. Suppose the fine endpoint couplings are $b_{m+1}^-$ and $b_{m+1}^+$. The drift estimate and~\eqref{eq:E-log-cost} give coarse endpoint bounds
\begin{equation*}
\beta_m^-\le b_m^- -\tfrac14\slope\log\ell_m<b_m^- ,
\qquad
\beta_m^+\ge b_m^+ +\tfrac34\slope\log\ell_m>b_m^+ .
\end{equation*}
For a continuous function whose endpoints bracket $[b_m^-,b_m^+]$, restrict between its last lower hit before its first upper hit. Its values then remain in this interval. Restrict the parameter interval accordingly. Every preceding representation is restricted along with the new one, so all already obtained identities persist.

Start at level~$N$ with the Wilson curve restricted to $[b_N^-,b_N^+]$ and repeat this operation down to zero. Since $b_0^-=b_0^+=\bar\beta$, it produces a trajectory with the required endpoint. The case $N=0$ is the Wilson representation at $\bar\beta$. All retained representations and adjacent pairs belong to the chosen classes and relation because Theorem~\ref{prop:RG} returns admissible representations whose adjacent pairs lie in the chosen relation.

The exact density identity at each step is the one in Theorem~\ref{prop:RG}, now applied only on the retained interval. By Lemma~\ref{lem:B1}, integration of the block kernel of Definition~\ref{def:D04} in its coarse argument is one. Fubini therefore preserves the integral, and the Wilson start gives its displayed value at every level. The path arrays were fixed before~$p$ in Lemma~\ref{lem:C1}; restricting the parameter interval does not change them.

Summing the drift from zero gives
\begin{equation}\label{eq:E-running}
|\beta_{N,p,m}-\bar\beta-\slope s_m|
\le\sum_{i<m}\omega_{J_{i+1}}^{\rm RG}.
\end{equation}
For a fixed splitting level~$k$, the tail is at most
$\sup_{i\ge k}J_{i+1}^{-2}s_m$ by~\eqref{eq:E-log-cost}. The initial sum is fixed. Divide by $s_m$, let $m\to\infty$ and then $k\to\infty$. Since $\bar\beta/s_m\to0$, this proves the uniform running law in Theorem~\ref{prop:MT}; neither bound depends on $N\ge m$ or~$p$.
\end{proof}

The uniqueness estimate concerns the specified recursion, rather than arbitrary representations of a density. Lemma~\ref{lem:D-E1} supplies its weighted analytic coordinates and their real power-series extension. Section~\ref{sec:partII-E.3} uses the fixed-level convergence of Lemma~\ref{lem:D-fixed-level}.
\begin{lemma}[Analytic endpoint tuning]\label{lem:E1.3}
There is $w_{\rm tun}>0$, depending only on the construction, such that, for every $N,p$ and $s\in[-1,1]$, exactly one bare coupling $\beta_{N,p}(s)$ starts a trajectory of the specified recursion ending at $\bar\beta+w_{\rm tun}s$. Moreover,
\begin{equation*}
\beta_{N,p}(0)=\beta_{N,p},\qquad
\beta_{N,p}(s)=b_{N,p}(s)\quad(-1<s<1),
\end{equation*}
where $b_{N,p}$ is holomorphic on $|z|<1$.
\end{lemma}
\begin{proof}
Uniqueness and the analytic extension both come from the weighted contraction of Lemma~\ref{lem:D-E1}, applied to histories of the specified recursion (Definitions~\ref{def:D06} and~\ref{def:D25a}). Let $\mathfrak d$ be the history distance of two histories in the domain of that lemma, including the coupling coordinate, and $\delta$ their coarse endpoint difference. With the constants $b_{\min},\lambda_*>0$ fixed before that lemma, independent of $N,p$, it gives
\begin{equation}\label{eq:E-analytic-input}
\mathfrak d\le |\delta|/b_{\min}+8g_*\mathfrak d,\qquad g_*=1/128.
\end{equation}
For equal endpoints, $\mathfrak d\le \mathfrak d/16$ implies equality of the histories and hence of their bare couplings.

Take $w_{\rm tun}=\lambda_*b_{\min}/4$, where $\lambda_*=1/16$ is fixed before Lemma~\ref{lem:D-E1}. The local action is affine in the varying coupling; the masks and sector domains are fixed in its analytic coordinates. Its power-series map preserves the radius-$\lambda_*$ ball and contracts by $1/2$ for endpoint series $w_{\rm tun}z$. Successive substitution converges in the norm $\sum_j\|u_j\|$ of absolutely summable series. The map preserves real coefficients, so evaluation gives a real trajectory in that domain for every $s\in[-1,1]$, pointwise in $s$. This supplies existence on the whole stated interval, including its endpoints. The selected coefficient at $s=0$ is the trajectory of Lemma~\ref{lem:E1.2}, by uniqueness. The bare-coupling series is holomorphic in $|z|<1$ and has the claimed real values. All clipping and health conditions are covered by the analytic neighbourhood in Lemma~\ref{lem:D-E1}; they are not differentiated across their boundaries.
\end{proof}

\subsection{Marked fields and volume-uniform estimates}\label{sec:partII-E.2}

Proposition~\ref{prop:UB} needs bounds on differentiated insertions, uniform in the cutoff and the volume. The argument has four parts. (i)~One marked insertion is carried to coarser levels as a polynomial coordinate plus remainders, with a two-sided bound on the coordinate that loses no factor $\ell_m^d$ in the inverse direction (Section~\ref{sec:partII-E.2.1}). (ii)~Placing one insertion in every reflected cell of a tiling of side $a_m$ and integrating down to a level~$k$ with $(a_m/a_k)^4$ sites per cell gives the chessboard bound $Aa_m^{-d-\varepsilon}$ (Section~\ref{sec:partII-E.2.2}). (iii)~A lattice difference of a degree-$d$ insertion is a degree-$(d+1)$ insertion up to a negligible error; its coefficient splits into an equation part, which vanishes in separated correlations, and a part in a space of dimension at most $\dim_\C\maxwellspace{d+1}$ (Sections~\ref{sec:partII-E.2.3}--\ref{sec:partII-E.2.4}). (iv)~That part is bounded through a nondegenerate matrix of free Maxwell tests (Section~\ref{sec:partII-E.2.5}). Lemma~\ref{lem:E2.4} combines (ii)--(iv) and proves Proposition~\ref{prop:UB} (Section~\ref{sec:partII-E.2.6}). The bounded approximations of Section~\ref{sec:partII-E.2.7} serve the phase argument of Section~\ref{sec:partII-H}.

\subsubsection{One marked insertion}\label{sec:partII-E.2.1}

Fix a degree~$d$ and suppress $N,p,d$ while describing one transfer. We use the marked densities of~\eqref{eq:I12-marked-evolution}, the local section $P_d=\localsection{d}$ of Lemma~\ref{lem:A1}, and the factor-count scaling $t^{\factorcount}$ of Definition~\ref{def:D08}. Write $\rho_m$ for the unmarked representation, $v_{{\rm sc},m}$ for its scalar normalization and $\Lambda_{n_{p,m}}$ for its site torus. An anchor is the lattice site at which the insertion is placed. Denote the marked regular supports by $Z_i$, to distinguish them from the unmarked supports~$X_i$. For a bare polynomial coordinate $\mathbf v\in V_d$,
\begin{equation}\label{eq:E-marked-evolution}
 H_N(\mathbf v)=\density\rho_N\,P_d[t_N^{-\factorcount}\mathbf v],\qquad
 H_m(\mathbf v)=\ell_m^d\blockmap{m}_{\beta_{N,p,m+1}}H_{m+1}(\mathbf v),
\end{equation}
so that $H_m(\mathbf v)=H^{\rm ins}_{N,p,m;d}(\mathbf v)$. These densities keep track of one insertion under exact integration. The equality is exact; the estimates below control the coordinates in which it is written.

\begin{definition}[Marked density]\label{def:D26}
A representation of $H_m(\mathbf v)$ at level~$m$ consists of a coordinate $\mathbf x=x_0\mathbf1+\mathbf x^\circ\in V_d$, complex regular remainders $R_i$ on supports~$Z_i$, and forced remainders $Q_i$ on forced addresses of~$\rho_m$, such that the exact decomposition~\eqref{eq:I12-marked-representation} holds with $\mathbf x_m=\mathbf x$, $R^{\rm reg}_{m,i}=R_i$ and $Q^{\rm for}_{m,i}=Q_i$. The section and the stencil~$\Sigma_d$ are translated to the anchor.
\end{definition}
The remainders restore the exact density; eligibility is that of Definition~\ref{def:D06}. The marked regular norm below is anchored at the insertion site, taken to be zero; it is not anchored separately at each support. For selected majorants $e_i\ge\|R_i\|_m$, $k_i\ge\sup|Q_i|$, set
\begin{equation}\label{eq:E-marked-mass}
\begin{aligned}
\mathcal E^{\rm mk}_m&=\sum_i e^{\sigma_{J_m}\cover_{B_m}(Z_i)}e_i,
\qquad \mathcal K^{\rm mk}_m=\sum_i e^{\sigma_{J_m}\cover_{B_m}(Y_i)}k_i,
\\
\|g_m\|_{\rm rem}&=\mathcal E^{\rm mk}_m+(\epsilon_m^{\rm for})^{-1/2}\mathcal K^{\rm mk}_m.
\end{aligned}
\end{equation}
The derivative seminorm is Definition~\ref{def:D05a}. A coefficient graph at level~$m$ is a linear choice $\mathbf x\mapsto g_m(\mathbf x)=(R(\mathbf x),Q(\mathbf x))$ of such remainders, with the supports fixed and $\|g_m(\mathbf x)\|_{\rm rem}\le C^{\rm col}_{J_m}|\mathbf x^\circ|_1$. As in the graph transform for invariant manifolds, the marked step maps coefficient graphs at level $m+1$ to graphs at level~$m$ and halves the distance between two of them (Lemma~\ref{lem:E-marked-step}); the current polynomial coordinate thus bounds the whole inserted density.

The following step retains the complete inserted density while reexpressing its remainders in the new polynomial coordinate. The inverse estimate is essential: a bound in the original bare coordinate would accumulate at every step.

To leading order a step multiplies a homogeneous coordinate of engineering weight $w$ by $\ell_m^{d-w}$: the factor $\ell_m^d$ comes from the normalization in~\eqref{eq:E-marked-evolution} and $\ell_m^{-w}$ from blocking. These factors form the diagonal map
\begin{equation*}
 (\mathsf S_L\mathbf x)_i=L^{d-w_i}x_i,
 \qquad 0\le w_i\le d,
\end{equation*}
where $(w_i,k_i)$ are the engineering weight and factor count of the basis polynomial~$b_{d,i}$; the unit has $w_0=0$. What remains is a change of basis $T_m$, bounded with its inverse uniformly in $N$, $p$ and $m$ within a stage, obtained by reexpressing the fine section, evaluated on the free interpolation of a coarse jet, in the coarse basis. Precisely, for a coarse radial jet~$\mathsf j$, let $U^{W_0}(\mathsf j)$ be the fine stencil configuration obtained by exponentiating the free interpolation~$W_0$ of Lemma~\ref{lem:B3} of the coarse potential of~$\mathsf j$, the degree-$d$ Taylor polynomial, in the jet coordinates, of the coarse-link logarithms. Define $\mathsf R^{\rm lead}_m$ by
\begin{equation*}
 \left.\frac{d^q}{dh_{\rm lat}^q}(\mathsf R^{\rm lead}_mb)(\operatorname{sc}_{h_{\rm lat}}\mathsf j)\right|_{h_{\rm lat}=0}
 =\left.\frac{d^q}{dh_{\rm lat}^q}P_d[b]\bigl(0;U^{W_0}(\operatorname{sc}_{h_{\rm lat}}\mathsf j)\bigr)\right|_{h_{\rm lat}=0}
 \qquad(0\le q\le d),
\end{equation*}
put $(\mathsf D^{\rm eng}_L\mathbf x)_i=L^{w_i}x_i$, and set
\begin{equation*}
 T_m=t_m^{\factorcount}\mathsf R^{\rm lead}_m\mathsf D^{\rm eng}_{\ell_m}t_{m+1}^{-\factorcount}.
\end{equation*}
Taylor matching gives $(\mathsf R^{\rm lead}_m)_{bi}=\delta_{bi}\ell_m^{-w_i}$ for $w_b\le w_i$ and $=0$ for $k_b<k_i$. To bound the other entries, use the kernel form of Lemma~\ref{lem:B3}: at averaging order $j\ge d+3$,
\begin{equation*}
 |\nabla_f^rW_0(\ell_my+a,y')|
 \le C_{j,r}\ell_m^{-1-r}e^{-c_j|y-y'|},\qquad 0\le r\le d.
\end{equation*}
The choice of the reference level ensures this order. A jet potential grows polynomially, so the weighted bound~\eqref{eq:B-weighted-interp} with $c'=0$ and $h=d$ applies to it.
Thus its kernel sum and the required differences converge absolutely. Applying the displayed bound to each factor of a homogeneous jet monomial gives $|(\mathsf R^{\rm lead}_m)_{bi}|\le C\ell_m^{-w_i}$, with $C$ depending only on $d$ and the stage. Periodization gives the same bound on a torus by summing over its translates. The threshold formula in item~9 of Table~\ref{tab:constant-choices} and the entry choices of Lemma~\ref{lem:E1.1} give $0<t_m\le1$ and $\tfrac12\le t_m/t_{m+1}\le2$. Hence $T_m$ sends each basis vector to its own and higher weights, with nonzero diagonal $(T_m)_{ii}=(t_m/t_{m+1})^{k_i}$, and back substitution in increasing weight bounds $T_m^{\pm1}$ uniformly in $N$, $p$ and $m$ within a stage.

The marked step constructs, from a coefficient graph at level $m+1$, the coarse coordinate and a coefficient graph at level~$m$ for which the marked identity holds again. Its statement is in the coordinate $\mathbf u=\mathsf M_m\mathbf x$ extracted by the step.
\begin{lemma}[One marked step]\label{lem:E-marked-step}
For $m\ge\reflevel{d}$ and each coefficient graph $g_{m+1}$ at level~$m+1$, there are an invertible linear $\mathsf M_m$ on $V_d$ and a coefficient graph $g_m$ at level~$m$ such that every representation of $H_{m+1}(\mathbf v)$ with coordinate~$\mathbf x$ and remainders $g_{m+1}(\mathbf x)$ yields one of $H_m(\mathbf v)$ with coordinate $\mathbf u=\mathsf M_m\mathbf x$ and remainders $g_m(\mathbf u)$. Moreover:
\begin{enumerate}
\item[(i)] $|\mathbf u|_1\le C^{\rm col}_{J_{m+1}}\ell_m^d|\mathbf x|_1$ and $|\mathbf x|_1\le C^{\rm col}_{J_{m+1}}|\mathbf u|_1$;
\item[(ii)] $\mathsf M_m$ and $g_m$ respect conjugation and signed permutations, and $\mathsf M_m\mathbf1=\ell_m^d\mathbf1$ with zero remainders;
\item[(iii)] fine graphs differing by at most $\tau C^{\rm col}_{J_{m+1}}|\mathbf x|_1$ give coarse graphs differing by at most $\tfrac\tau2C^{\rm col}_{J_m}|\mathbf u|_1$.
\end{enumerate}
\end{lemma}
The proof gives the construction and the bounds, including the contraction between two choices of incoming remainders.
\begin{proof}
The new coordinate is read off by Taylor extraction after one exact block integration (Definition~\ref{def:D04}) in the local sections of Definition~\ref{def:D07}. The connected expansion of Lemma~\ref{lem:D-connected}, in a nilpotent algebra, controls the marked terms; the pre-reset estimates of Lemma~\ref{lem:D4}, the subtraction of Lemma~\ref{lem:D-local-subtraction} and the density identity of Lemma~\ref{lem:D5} provide the unmarked part of the step on which the marked terms ride. Put $L=\ell_m$, $S=C^{\rm col}_{J_{m+1}}$, $K=K_{J_{m+1}}$ and $\kappa_m=(\epsilon_m^{\rm for})^{1/2}$. All norms in this proof are coefficient sums or the anchored norm~\eqref{eq:E-marked-mass}.

\emph{Step 1: Marked terms in the local integration.}
Introduce $z^2=0$. Replace the regular exponential factor by its product with
$1+z(x_0+\mathbf1_{\Sigma_d\,\mathrm{healthy}}P_d[t_{m+1}^{-\factorcount}\mathbf x^\circ]+\sum_i\mathbf1_{Z_i\,\mathrm{healthy}}R_i)$ and replace a forced factor $K_i$ by $K_i+zQ_i$. The coefficient of $z$ in the eligible covering sum is exactly the representation of Definition~\ref{def:D26}. Apply the coefficientwise rooted clause of Lemma~\ref{lem:D-connected} in this nilpotent algebra. Each connected term now contains one distinguished factor. For a term with ordinary factors $a_1,\ldots,a_h$, its marked coefficient is
\begin{equation*}
 [z]\prod_{i=1}^h(a_i+zb_i)
 =\sum_{i=1}^h b_i\prod_{j\ne i}a_j.
\end{equation*}
Root its connected support at this distinguished factor. The sums over all other factors are bounded by the rooted clause of Lemma~\ref{lem:D-connected}; there is no sum over possible translates of the marked anchor. Absolute convergence justifies the displayed extraction and integration term by term. A regular mark has its derivative majorant $e_i$; a forced mark has its amplitude majorant $k_i$. The latter retains its original eligibility event, so the large-field estimate used in the forced expansion still applies.

The three contributions to a basis input are bounded as follows. The map $T_m$ fixes the unit and preserves the nonconstant subspace. It is triangular in the homogeneous weights: $(T_m)_{ij}\ne0$ and $i\ne j$ imply $w_j<w_i$. Its diagonal entries are nonzero, so inversion proceeds in increasing weight, with the finite stage bounds included in~$K$. Let $C_{\rm loc}\ge1$ bound its two operator norms and the following local remainder estimates before symmetry averaging; the leading stage bound was chosen with $K\ge2A_d^{\rm sym}C_{\rm loc}$. Write $\epsilon=\epsilon^{\rm mk}_m$ for the local marked error of item~\ref{it:E-graph}. Taylor subtraction of a primary mark leaves coefficient error at most $\epsilon|\mathbf x^\circ|_1$ and remainder at most $C_{\rm loc}|\mathsf S_L\mathbf x^\circ|_1$. A regular remainder of mass $e^{\sigma\cover(Z)}e_i$ contributes at most $\epsilon e^{\sigma\cover(Z)}e_i$ to each of these two bounds. A forced remainder contributes at most $\epsilon e^{\sigma\cover(Y)}k_i/\kappa_{m+1}$.

For the primary mark, the local polynomial is subtracted through degree~$d$ before rescaling. Each remaining term has at least one additional scaled jet weight, and the block rescaling gives the inverse block factor in its norm. The finite connected sums and the ordered derivative product rule introduce only stage constants.

An incoming regular mark is never moved to the forced family. A forced remainder must vanish off an eligibility event with a nonempty inventory of flagged sites (Definition~\ref{def:D25}), whereas the support of a regular mark may wrap around the torus while every plaquette in it is healthy. The integrated contribution of a regular mark consists of regular coarse remainders and, from those terms of its connected expansion that carry a flagged inventory at level~$m$, of forced coarse remainders on their own eligibility events. Wrapping plays no role in this assignment; it only decides whether the radial cancellation can be used. Fix such a mark $R_i$ on $Z_i$, with derivative majorant $e_i$, let $\widehat Z_i$ be the coarse support of its contribution, and put
\begin{equation*}
 d_i=\cover_{B_{m+1}}(Z_i),\qquad
 \widehat Z_i\subset N_{R_j^{\rm supp}B_m}(\operatorname{parent}(Z_i)).
\end{equation*}
Call $Z_i$ short if $d_i^3<\ell_m$. Three classes of supports are estimated separately.

(i) If $\cover_{B_m}(\widehat Z_i)>B_m$, the contribution is bounded at a coarse support exponent at least $\sigma_{J_m}+3$, one more than the exponent $\sigma_{J_m}+2$ in~\eqref{eq:E-support-conversion}. Half of this unit excess absorbs the derivative losses, which are polynomial in $\cover_{B_m}(\widehat Z_i)$. The other half gives a factor $e^{-B_m/2}$, since $\cover_{B_m}(\widehat Z_i)>B_m$. As for the large supports in the proof of Lemma~\ref{lem:D-inherited}, the resulting factor tends to zero as the coupling entry is raised at fixed stage and ratio; it is part of the term $o_{\widehat\beta_m}(1)$ below.

(ii) If $\cover_{B_m}(\widehat Z_i)\le B_m$ and $Z_i$ is short, then $Z_i$ is nonwrapping about the anchor. Indeed, as in the proof of Lemma~\ref{lem:D-inherited}, $3d_i\le\ell_m$ and $B_{m+1}\le2B_m$ bound the fine radius of $Z_i$ by $6(1+B_m)\ell_m$, while the coarse side is at least $100(1+B_m)$. The radial-jet coefficients of $R_i$ through degree~$d$ therefore vanish: this is the format of the incoming remainders, recorded in Lemma~\ref{lem:E2.2a} at level~$m+1$. Taylor's formula along the healthy radial segment gives the factor $\ell_m^{-(d+1)}$, with derivative losses polynomial in $1+d_i$.

(iii) If $\cover_{B_m}(\widehat Z_i)\le B_m$ and $d_i^3\ge\ell_m$, no radial cancellation is used. By~(ii), this class contains every wrapping support with $\cover_{B_m}(\widehat Z_i)\le B_m$; it also contains nonwrapping supports that are not short. The derivative bound without subtraction has losses polynomial in $1+d_i$, and
\begin{equation*}
 1\le\ell_m^{-(d+1)}(1+d_i)^{3(d+1)}
\end{equation*}
gives the same factor $\ell_m^{-(d+1)}$ with a further polynomial factor in $1+d_i$.

In classes (ii) and (iii), the support conversion~\eqref{eq:E-support-conversion} gives $(\sigma_{J_m}+2)\cover_{B_m}(\widehat Z_i)\le\log\mathcal A_m+\tfrac12\sigma_{J_{m+1}}d_i$. The remaining half of the incoming exponent bounds the polynomial in $1+d_i$, since $\sup_{s\ge0}(1+s)^qe^{-\sigma_{J_{m+1}}s/2}<\infty$ for each fixed~$q$. After multiplication by $\ell_m^d$, the contribution of the mark is at most $C\ell_m^{-1}e^{\sigma_{J_{m+1}}d_i}e_i$, with $C$ fixed before the ratio. Thus the smallness of a healthy wrapping mark comes from its unused support weight, not from a Taylor cancellation.

An incoming forced mark retains its own eligibility event, whether or not its support wraps. Its support sum keeps half the original forced suppression: its normalization is $\kappa_{m+1}$, not $\epsilon_{m+1}^{\rm for}$. The remaining suppression absorbs the same derivative and support factors. Thus the bounds before choosing parameters have the form $C_j^{\rm loc}/L+o_{\widehat\beta_m}(1)$, with $C_j^{\rm loc}$ fixed before $L$ and the second term tending to zero at fixed stage and block factor; the ratio and entry choices of item~\ref{it:E-graph}, made in the proof of Lemma~\ref{lem:E1.1}, make it at most $\epsilon^{\rm mk}_m\le L^{-2/5}$. The degree-$d$ Taylor extraction is on healthy charts; the indicators themselves are never differentiated.

Let $\mathbf u$ be the extracted coarse polynomial coordinate, $r$ the incoming anchored remainder norm and $r'$ the resulting coarse remainder norm, including its forced normalization. Summing the contributions gives
\begin{equation}\label{eq:E-atom-estimates}
 |\mathbf u-T_m\mathsf S_L\mathbf x|_1\le\epsilon(|\mathbf x^\circ|_1+r),\qquad
 r'\le C_{\rm loc}|\mathsf S_L\mathbf x^\circ|_1+\epsilon r,
 \qquad r=\|g_{m+1}(\mathbf x)\|_{\rm rem}.
\end{equation}
The regular-remainder terms in~\eqref{eq:E-atom-estimates} use the vanishing Taylor coefficients through degree~$d$ on short supports, and the unused support weight on the other supports. Retained incoming regular terms use the support conversion factor~\eqref{eq:E-support-conversion}. The same sums with the primary mark canceled give the Lipschitz bound $\epsilon\|g-\widetilde g\|_{\rm rem}$ for the coefficient and remainder differences; this uses the entire forced marked family at once.

\emph{Step 2: Linear choice and inverse bound.}
Perform the construction on each homogeneous basis vector and on all its signed and conjugated images. Take their union of supports, with zero entries where necessary, and recombine using the basis coefficients. Average over the finite signed-permutation group and over conjugation. This gives a complex-linear family with the required symmetries and exact density identity. Let $A_d^{\rm sym}\ge1$ bound the sum of absolute recombination coefficients in this finite averaging. After averaging, the error in~\eqref{eq:E-atom-estimates} is at most $\epsilon(1+S)A_d^{\rm sym}|\mathbf x^\circ|_1$, and the remainder is at most $(C_{\rm loc}+\epsilon S)A_d^{\rm sym}|\mathsf S_L\mathbf x^\circ|_1$.

Use~\eqref{eq:E-graph-reserve}. The cross-stage inequality follows from item~\ref{it:E-columns}, fixed before the ratios. Since $\epsilon S\le1\le C_{\rm loc}$, the averaged remainder constant is at most $K$. In the notation $\mathbf y=\mathsf S_L\mathbf x$, the resulting inequalities are
\begin{equation*}
 |\mathbf u-T_m\mathbf y|_1\le \varepsilon_L(|\mathbf x^\circ|_1+r),\qquad
 r'\le K|\mathbf y^\circ|_1+\varepsilon_L r,\qquad r\le S|\mathbf x^\circ|_1.
\end{equation*}
Since $|\mathbf x|_1\le|\mathsf S_L\mathbf x|_1\le L^d|\mathbf x|_1$, the first two conditions in~\eqref{eq:E-graph-reserve} imply
\begin{equation*}
 |\mathbf y|_1\le2K|\mathbf u|_1,\qquad |\mathbf u|_1\le2K|\mathbf y|_1,
 \qquad |\mathbf y^\circ|_1\le2K|\mathbf u^\circ|_1.
\end{equation*}
For example, apply $T_m^{-1}$ and use
$2K|\mathbf u-T_m\mathbf y|_1\le|\mathbf y^\circ|_1$ to absorb the error; the same argument on the nonconstant subspace gives the last bound. Moreover $\varepsilon_L r\le K|\mathbf y^\circ|_1$, so
\begin{equation*}
 r'\le2K|\mathbf y^\circ|_1\le4K^2|\mathbf u^\circ|_1
 \le C^{\rm col}_{J_m}|\mathbf u^\circ|_1.
\end{equation*}
The linear map $\mathbf x\mapsto\mathbf u=\mathsf M_m\mathbf x$ is injective by $|\mathbf x|_1\le S|\mathbf u|_1$, hence invertible in finite dimension. Let $\widetilde g_m(\mathbf x)$ denote the remainders just constructed for the fine coordinate $\mathbf x$. Define the new graph by $g_m(\mathbf u)=\widetilde g_m(\mathsf M_m^{-1}\mathbf u)$. The last display is its required graph bound at level~$m$. The two coefficient estimates of~(i) now follow; property~(ii) is inherited from the symmetric construction. In particular the inverse estimate has no factor $L^d$: $\mathsf S_L^{-1}$ has norm at most one.

For the unit, the primary extraction is exact, $T_m$ fixes it and its nonconstant coordinate is zero. All the preceding remainder bounds are then zero, and the unit is sent to $L^d\mathbf1$ with zero remainders. Finally, when two fine graphs differ by at most $\tau S|\mathbf x|_1$ in the remainder norm, the Lipschitz estimate above and rebasing by the inverse bound of~(i) give a coarse difference at most $\tfrac12\tau C^{\rm col}_{J_m}|\mathbf u|_1$: the two contributions are the changed remainder at fixed input and the changed input coordinate, each bounded by $\tau/4$ using~\eqref{eq:E-graph-reserve}. This is the graph contraction used in Lemma~\ref{lem:E-fixed-level}.
\end{proof}


Choose the reference level $\reflevel{d}$ with $J_{\reflevel{d}}\ge d+3$, $\mathfrak q_{J_m}\ge d+1$ and $\varpi_{J_m}\ge d$ for $m\ge\reflevel{d}$. Apply Lemma~\ref{lem:E-marked-step} along the trajectory, starting from the representation at the cutoff constructed in the proof of Lemma~\ref{lem:E2.1}, and let $\mathsf M_m$ be its coefficient maps. Define, before normalizing the lattice fields,
\begin{equation}\label{eq:E-terminal-map}
 \mathsf Z_{N,p,d}=\mathsf M_{\reflevel{d}}\cdots\mathsf M_{N-1}\quad(N\ge\reflevel{d}).
\end{equation}
An empty product is the identity. For $N<\reflevel{d}$ choose an arbitrary invertible map, as allowed in Definition~\ref{def:D08}.

\paragraph{The curvature square at the terminal map.} We continue the example of Section~\ref{sec:A-example}. The curvature square $c_{F^2}$ is a label of degree $d=4$ with engineering weight four and factor count two, the top weight in~$V_4$. Its column of the leading map is therefore diagonal, $T_mc_{F^2}=(t_m/t_{m+1})^2c_{F^2}$, and $\mathsf S_{\ell_m}$ fixes it. With the leading maps alone, the terminal map would multiply $c_{F^2}$ by the single factor $(t_{\reflevel{4}}/t_N)^2$. The coefficient maps, however, are controlled only up to an error in the full coefficient norm, $|\mathsf M_m\mathbf x-T_m\mathsf S_{\ell_m}\mathbf x|_1\le\varepsilon_{\ell_m}(|\mathbf x^\circ|_1+\|g_{m+1}(\mathbf x)\|_{\rm rem})$ (Step~2 in the proof of Lemma~\ref{lem:E-marked-step}), so at each step the curvature-square input may acquire a component along the unit $\mathbf1$, the only other real $\hypercubic$-invariant label in~$V_4$ (Section~\ref{sec:A-example}), for instance from integrated remainders. The estimates therefore do not force the unit coefficient~$\delta^{F^2}_{N,p}$ of Section~\ref{sec:A-example} to vanish. The reference normalization instead prescribes the whole bare coefficient vector $\mathsf Z_{N,p,4}^{-1}c_{F^2}$; Lemma~\ref{lem:E2.1} bounds it by $|c_{F^2}|_1\prod_{\reflevel{4}\le i<N}C^{\rm col}_{J_{i+1}}=a_N^{-o(1)}|c_{F^2}|_1$, by~\eqref{eq:E-log-cost}. Section~\ref{sec:partII-F} continues the example with the quotient field.

The next lemma carries one insertion from the cutoff to the reference level. Its content is the representation at the cutoff, whose forced part needs the initialization requirement, and the transport bounds obtained by iterating the marked step.
\begin{lemma}[Marked trajectory]\label{lem:E2.1}\label{lem:E2.2}
Let $N\ge\reflevel{d}$ and $\reflevel{d}\le m\le N$. The densities~\eqref{eq:E-marked-evolution} have representations at level~$m$ with coordinates $\mathbf x_m(\mathbf v)$ and remainder families linear in~$\mathbf v$, with fixed supports, such that $\mathbf x_N(\mathbf v)=\mathbf v$, $\mathbf x_m=\mathsf M_m\mathbf x_{m+1}$ for $m<N$, $\mathbf x_{\reflevel{d}}(\mathbf v)=\mathsf Z_{N,p,d}\mathbf v$ with $\mathsf Z_{N,p,d}$ invertible, and
\begin{equation*}
\mathcal E^{\rm mk}_m(\mathbf v)+(\epsilon_m^{\rm for})^{-1/2}\mathcal K^{\rm mk}_m(\mathbf v)
\le C^{\rm col}_{J_m}|\mathbf x_m(\mathbf v)^\circ|_1,
\end{equation*}
also with $\sigma_{J_m}$ replaced by any exponent in $[98,100]$. For $m<N$,
\begin{equation*}
|\mathbf x_{m+1}(\mathbf v)|_1\le C^{\rm col}_{J_{m+1}}|\mathbf x_m(\mathbf v)|_1,
\qquad |\mathbf x_m(\mathbf v)|_1\le C^{\rm col}_{J_{m+1}}\ell_m^d|\mathbf x_{m+1}(\mathbf v)|_1.
\end{equation*}
\end{lemma}
\begin{proof}
At level~$N$, split the unmasked section into its healthy part and its complement $f_{\mathbf v}=\mathbf1_{\Sigma_d\,\mathrm{unhealthy}}P_d[t_N^{-\factorcount}\mathbf v^\circ]$. The healthy part is the polynomial term of~\eqref{eq:I12-marked-representation}; the complement must be carried by the forced family, whose eligibility events (Definition~\ref{def:D06}) cover the flagged configurations. For an eligible family~$\mathcal A$ let $\mathcal A_\Sigma$ be its forced terms whose supports meet~$\Sigma_d$. If $\mathcal A_\Sigma=\varnothing$, exact coverage leaves no flag in~$\Sigma_d$ and $f_{\mathbf v}=0$. Otherwise replace $\mathcal A_\Sigma$ by one new marked forced address with support $\Sigma_d\cup\bigcup_{i\in\mathcal A_\Sigma}Y_i$, inventory $\bigcup_{i\in\mathcal A_\Sigma}I_i$, zero unmarked activity, and $Q=f_{\mathbf v}\prod_{i\in\mathcal A_\Sigma}K_i^{\rm for}$ on the event that the flags in that support are exactly this inventory, $Q=0$ otherwise. An index of the new marked family records the cluster~$\mathcal A_\Sigma$, even when two clusters have the same support and inventory; retain all original forced factors in~$Q$. The remaining supports are disjoint from the merged support, since they avoid~$\Sigma_d$ and were disjoint from every merged support. This is a bijection between eligible families meeting~$\Sigma_d$ and eligible marked families containing one new address, so the forced sum restores $f_{\mathbf v}\density\rho_N$ exactly. The regular remainder is zero and $\mathbf x_N=\mathbf v$.

The support weight of the merged address is at most $e^{\sigma_{J_N}\cover_{B_N}(\Sigma_d)}$ times the product of the old support weights: their connected covers all meet the stencil cover. The activity caps bound the sum of the weighted majorants of addresses meeting~$\Sigma_d$ by $|\Sigma_d|\epsilon_N^{\rm for}$. Dropping disjointness and summing their products over nonempty subsets gives at most $e^{|\Sigma_d|\epsilon_N^{\rm for}}-1$. By~\eqref{eq:A-global-bound}, $|f_{\mathbf v}|\le C_dt_N^{-d}|\mathbf v^\circ|_1$; with~\eqref{eq:I12-activity-caps} and $\cover_{B_N}(\Sigma_d)\le1$,
\begin{equation*}
\begin{aligned}
 (\epsilon_N^{\rm for})^{-1/2}\mathcal K^{\rm mk}_N
 &\le(\epsilon_N^{\rm for})^{-1/2}C_dt_N^{-d}e^{\sigma_{J_N}}
 (e^{|\Sigma_d|\epsilon_N^{\rm for}}-1)|\mathbf v^\circ|_1\\
 &\le|\mathbf v^\circ|_1
\end{aligned}
\end{equation*}
by item~\ref{it:E-birth}. This is the bound of the lemma at $m=N$, since $C^{\rm col}_{J_N}\ge1$.

Iterate the marked step on the same unmarked trajectory. Lemma~\ref{lem:E-marked-step} gives $\mathbf x_m=\mathsf M_m\mathbf x_{m+1}$, the exact equality in~\eqref{eq:E-marked-evolution}, the graph bound and the two transport bounds; the factor $\ell_m^d$ agrees with $a_m^d=\ell_m^da_{m+1}^d$. Each $\mathsf M_m$ is invertible, so their product is invertible and $\mathbf x_{\reflevel{d}}=\mathsf Z_{N,p,d}\mathbf v$. No choice of a different unmarked representation is made. For fixed supports, take the union of the addresses selected for the finitely many basis inputs and fill unused addresses by zero; the basis construction in Lemma~\ref{lem:E-marked-step} then makes the coefficients and both remainders linear in~$\mathbf v$. Finally, for $s\ge0$ and $98\le\tau\le100$ we have $e^{\tau s}\le e^{100s}$, so lowering the support exponent in~\eqref{eq:E-marked-mass} only decreases both sums.
\end{proof}

The normalization retains the symmetries of the local sections. Its unit includes the degree-dependent spacing factor.
\begin{lemma}[Symmetry and the unit]\label{lem:E2.1a}
For every $d$, $N\ge\reflevel{d}$ and $p$, the maps $\mathsf Z_{N,p,d}$ commute with conjugation and~$\hypercubic$. For $\reflevel{d}\le m\le N$,
\begin{equation*}
\mathbf x_m(a_N^d\mathbf1)=a_m^d\mathbf1.
\end{equation*}
Consequently $\mathbf u_d=a_{\reflevel{d}}^d\mathbf1$ has insertion one.
\end{lemma}
\begin{proof}
Complex conjugation and signed permutations commute with integration (Definition~\ref{def:D04}), with Taylor extraction in the symmetric sections of Lemma~\ref{lem:A1}, and with the equivariant coefficient graph of Lemma~\ref{lem:E-marked-step}. Each operation therefore commutes with every $\mathsf M_m$ and their product.

For a unit input, the marked density is $a_N^d\density\rho_N$. Exact integration and multiplication by $\ell_m^d$ give $a_m^d\density\rho_m$ at level~$m$. The selected extraction has zero nonconstant coordinate and zero remainder for this density. Induction proves the displayed identity. In particular
$\mathsf Z_{N,p,d}^{-1}\mathbf u_d=a_N^d\mathbf1$. Substitute into Definition~\ref{def:D08}; factor scaling fixes the polynomial unit, so its insertion equals $a_N^{-d}a_N^d=1$.
\end{proof}

The remainders have the cancellations of the marked format~\eqref{eq:I12-marked-representation} of Section~\ref{sec:renormalization}.
\begin{lemma}[Marked local format]\label{lem:E2.2a}
For $N\ge\reflevel{d}$ and $\reflevel{d}\le m\le N$, the complex regular remainders are Borel, gauge invariant, local and smooth as in Definition~\ref{def:D24}. Their supports contain the anchor. Their values at the identity, and their radial-jet coefficients through degree~$d$ on nonwrapping supports, vanish. The forced remainders satisfy Definition~\ref{def:D25}, and
\begin{equation*}
Q_i\ne0\quad\Longrightarrow\quad\text{the anchor belongs to }Y_i.
\end{equation*}
\end{lemma}
\begin{proof}
The anchored extraction in Lemma~\ref{lem:E-marked-step} subtracts the value at the identity and the radial Taylor polynomial; this gives exactly the two stated regular cancellations. On wrapping supports no radial Taylor cancellation is imposed. In the connected expansion every marked component contains the original insertion anchor. The support is the union of the supports in that component, so it still contains that anchor after enlargement. A component reaches a forced address only through terms that carry a flagged inventory at that level, and then only on the eligibility event of that address (Definition~\ref{def:D25}); an incoming regular mark is never moved to the forced family as a whole (Step~1 in the proof of Lemma~\ref{lem:E-marked-step}). The anchor argument applies to these terms as well. Locality, gauge invariance and Borel measurability are preserved by local integration and by these tests. The regular derivatives are taken on healthy configurations, as in Definition~\ref{def:D24}; no derivative of a relaxed threshold test is used. The fixed basis construction and its support unions retain these properties under iteration.
\end{proof}

\subsubsection{Chessboard seminorms and replicated sources}\label{sec:partII-E.2.2}

\begin{samepage}
A bond tiling has cell boundaries through link midpoints. Its reflected copies place an observable in every cell, with conjugation on odd reflections. The copy $(F)_C$ is the pullback of $F$ under the composition of bond reflections taking the base cell to $C$, conjugated when that composition has odd parity.
\begin{definition}[Chessboard seminorm]\label{def:D27}
For a tiling $\mathcal C^{\rm tile}$ into congruent reflected cells and an observable $F$ in its base cell, set
\begin{equation*}
\|F\|_{\mathcal C^{\rm tile}}=
\left|\left\langle\prod_{C\in\mathcal C^{\rm tile}}(F)_C\right\rangle_{N,p}\right|^{1/|\mathcal C^{\rm tile}|}.
\end{equation*}
Interior means dependence only on links with both endpoints strictly inside the base cell. The centered tiling $\mathcal C^{\rm tile}_m(y)$ has physical cell side $a_m$ and base cell centered at~$y$.
\end{definition}
\end{samepage}
The seminorm measures the size of a repeated observation throughout the box; it converts reflected-product estimates into bounds on one insertion. For interior complex observables these quantities are seminorms and satisfy
\begin{equation}\label{eq:E-chess}
\left|\left\langle\prod_{C\in\mathcal C^{\rm tile}}(F_C)_C\right\rangle_{N,p}\right|
\le\prod_{C\in\mathcal C^{\rm tile}}\|F_C\|_{\mathcal C^{\rm tile}}.
\end{equation}
Apply the abstract chessboard theorem, \cite[Theorem~4.1]{FrohlichIsraelLiebSimon1978}, successively in the four axes. The real vector space is the space of complex cell functions, and the reflection includes conjugation. Translation invariance gives cyclicity; reflection positivity gives the required Schwarz inequality. For the Wilson measure this positivity is \cite[Theorem~2.1, pp.~14--15]{Seiler1982}, with compact matrix group, positive coupling and observables on one side of the bond planes. For $\mathcal C^{\rm tile}_m$, there are $L_pD_m$ cells per axis, an even integer. The cell side $a_m/a_N=D_N/D_m$ is odd, so the centered boundaries pass through link midpoints. Bounded Borel cell functions follow by approximation in the finite-dimensional configuration measure.

The centered seminorm estimate must also control anchors near arbitrary, separated configurations. Write $v=(d,c)$ for the insertion label.
A signed image means the image under~$\hypercubic$; conjugation is also allowed. Two tilings are aligned when their cells have the same axes and every coarse boundary is a fine boundary.
\begin{lemma}[Interior seminorm bound]\label{lem:F1a}
Fix $m\le k\le N$, $N\ge\reflevel{d}$, and~$B_{\rm rt}\ge0$. Suppose, for every anchor~$y'$, every signed image or image under conjugation of~$O_{N,p}(v,y')$ has seminorm at most~$B_{\rm rt}$ in~$\mathcal C^{\rm tile}_k(y')$ and stencil within~$a_k$ of~$y'$. In every translated bond tiling~$\mathcal T$ of side~$a_m$, an anchor at distance at least~$3a_k$ from its cell boundary satisfies
\begin{equation}\label{eq:F-uncentered}
O_{N,p}(v,y)\text{ is interior},\qquad
\|O_{N,p}(v,y)\|_{\mathcal T}\le B_{\rm rt}\,.
\end{equation}
\end{lemma}
\begin{proof}
At an aligned fine-cell center the bound follows from~\eqref{eq:E-chess}; between centers it follows from log-convexity in each axis, using the symmetry of Lemma~\ref{lem:E2.1a}. Since~$a_k/a_N$ is odd, the fine-cell centers are lattice sites. At an aligned fine-cell center, reflect all coarse copies into distinct fine cells. Their signs and conjugations are among the assumed images. Apply~\eqref{eq:E-chess} with~$1$ in every unoccupied fine cell. Taking the coarse seminorm bounds it by~$B_{\rm rt}$.

For intermediate anchors, let~$H(z)$ be the coarse seminorm after~$z$ lattice translations in one axis, with the other coordinates fixed. Translation invariance and reflection give
\begin{equation}\label{eq:F-logconvex}
H(z)^2\le H(z-1)H(z+1)\,.
\end{equation}
Indeed, translate the reflected product by one lattice step. Alternate cells then contain shifts~$z-1$ and~$z+1$; apply~\eqref{eq:E-chess}. On a finite interval a nonnegative sequence satisfying this inequality cannot exceed its endpoint maximum: an interior positive maximum extends to both neighbours and hence to an endpoint. This includes zero endpoints.

Bracket each coordinate of~$y$ by its two neighbouring fine-center coordinates. All anchors in the resulting rectangle have margin at least~$2a_k$, so their stencils are interior. Induct on~$i=0,\ldots,4$: the seminorm is at most~$B_{\rm rt}$ when the first~$i$ coordinates are arbitrary lattice coordinates in these intervals and the remaining coordinates are endpoints. The case~$i=0$ is the aligned-center bound. In the next axis, the endpoint bounds hold uniformly in the coordinates already moved;~\eqref{eq:F-logconvex} supplies the intervening bounds. The case~$i=4$ proves~\eqref{eq:F-uncentered}.
\end{proof}

\paragraph{Replicated sources.}
Several insertions are carried at once by giving each a nilpotent mark. This is the device of the observable field in~\cite[Section~2.3]{BauerschmidtBrydgesSlade2015}, where only terms of order at most one in each external field are retained. For a finite set of slots $I_0$, work in the commutative algebra $\C[z_i:i\in I_0]/(z_i^2:i\in I_0)$. Write $z_I=\prod_{i\in I}z_i$ for the mode~$I$. The exact source density starts at
\begin{equation}\label{eq:E-multiple-density}
 \mathcal H_N(z;U)=\density\rho_N(U)\prod_{i\in I_0}(1+z_iF_i(U)),
 \qquad \mathcal H_m=\blockmap{m}_{\beta_{N,p,m+1}}\mathcal H_{m+1}.
\end{equation}
Thus its normalized $z_I$ integral is the cutoff moment of the corresponding physical insertions. At level~$m$ write its local representation as
\begin{equation}\label{eq:E-source-form}
\begin{aligned}
 &e^{\Phi_m(U)+\sum_i z_if_{i,m}(U)+\sum_X H_{m,X}(z;U)}
 \\
 &\quad\cdot\sum_{\mathcal A\ {\rm eligible}}
 \prod_{Y\in\mathcal A}
 \left(K_{m,Y}(U)+\sum_i z_i k_{i,m,Y}(U)+K^{\ge2}_{m,Y}(z;U)\right).
\end{aligned}
\end{equation}
Here $\Phi_m$ is the unmarked exponent, $f_{i,m}$ and $k_{i,m,Y}$ are the complete one-insertion exponential and forced terms, including their physical spacing factors. The polynomials $H_{m,X}$ and $K^{\ge2}_{m,Y}$ contain only modes $z_I$ with $|I|\ge2$. Their support contains every anchor indexed by a nonzero mode. The regular terms include their health masks, and each forced term retains its eligibility test. For a single mark, $z_1^2=0$ gives $e^{\Phi_m+z_1f_{1,m}}=e^{\Phi_m}(1+z_1f_{1,m})$, and~\eqref{eq:E-source-form} is the representation of Definition~\ref{def:D26} with
\begin{equation*}
 f_{1,m}=a_m^{-d}\Bigl\{x_0+\mathbf1_{\Sigma_d\,\mathrm{healthy}}P_d[t_m^{-\factorcount}\mathbf x_m^\circ]+\sum_i\mathbf1_{Z_i\,\mathrm{healthy}}R_i\Bigr\},
 \qquad k_{1,m,Y_i}=a_m^{-d}Q_i,
\end{equation*}
at $\mathbf v=\mathsf Z^{-1}_{N,p,d}c$. The higher modes record the connected terms that integration creates between different marks. Rescaling mark~$h$ by $\mathsf w_h>0$, its \emph{mark weight}, multiplies the coefficient of $z_I$ by $\prod_{h\in I}\mathsf w_h$; a bound of order one on the rescaled source therefore bounds that coefficient by $\prod_{h\in I}\mathsf w_h^{-1}$.

At the cutoff, the healthy and unhealthy parts of different insertions
must be assembled together. A single forced support can meet more than
one unhealthy stencil, even when their anchors are far apart. The
following construction gives the higher forced modes required in that
case; no activity is divided by.

For a two-mark example, let $K(U)$ be a selected eligible forced
activity, and let $b_1(U),b_2(U)$ be two unhealthy complement factors
whose stencils meet its support. Suppressing the common field~$U$, the
part formed using these factors is
\[
 K(1+z_1b_1)(1+z_2b_2)
 =K+z_1Kb_1+z_2Kb_2+z_1z_2Kb_1b_2.
\]
The coefficient of $z_1z_2$ contains the original activity once. A
product that repeats a mark, such as $z_1(z_1z_2)$, is zero. This exact
algebra does not by itself give the support estimates needed for
propagation.

\begin{lemma}[Joint terminal forced modes]\label{lem:E-joint-terminal}
Fix a finite set of allowed degrees $\mathcal D$ and a finite family of
primary insertions of degrees $d_h\in\mathcal D$, with distinct anchors.
Assume their stencils have the fixed radii of their degrees and their
positive mark weights at $N$ satisfy
$\mathsf w_{h,N}a_N^{-d_h}|\mathbf x_{h,N}|_1\le\zeta_{d_h,N}$.
At a sufficiently late Wilson start of Lemma~\ref{lem:D3}, the exact polynomial
\eqref{eq:E-multiple-density} admits the form~\eqref{eq:E-source-form}:
its regular higher modes vanish, and its forced higher modes are
indexed by connected unions of unhealthy stencils and selected original
forced supports. Each retains the original eligibility tests, the union
of their inventories, and the product of each selected activity once.
For a fixed degree set the weighted forced sums, including the higher
modes, obey the attachment bound of
Lemma~\ref{lem:E-source-attachments} on a sufficiently late terminal
level, uniformly in the number of marks, cutoff and period, provided
there is at most one anchor at each site. The singleton terms agree with
the construction in Lemma~\ref{lem:E2.1}.
\end{lemma}
\begin{proof}
Expand the finite product in~\eqref{eq:E-multiple-density}, and split
each chosen section into its healthy term and its unhealthy complement.
Fix a term and an eligible family of original forced activities. Join
each complement stencil to the selected supports it meets, and join
intersecting complement stencils. Every complement contains a flag, so
exact coverage gives it at least one incident selected support. Collapse
each connected component containing a complement to one new address:
its support and inventory are the unions of those of its constituents,
and its value is the product of the complements and selected activities,
with their original eligibility tests and the exact flag test on the
union. Leave untouched selected supports in the cover. The collapsed
supports are disjoint, and expansion of their constituent indices
recovers the original selection and the assignment of marks uniquely.
Thus this is an identity coefficient by coefficient, including at zeros
of the original activities. Add zero unmarked slots for these indexed
addresses when needed. For one complement the component is exactly the
merged address of Lemma~\ref{lem:E2.1}.

For the bound, use the complete stencil as the read support and the
absolute majorant of each complement. The weighted singleton estimate
used in~\eqref{eq:E-attachment-allowance} bounds its size by
$A_{j,d_h}\widehat\beta_N^{q^{\rm ht}_{d_h}}\zeta_{d_h,N}$, with the
fixed stencil and support factors included. It is at most
$A_{j,d_h}\widehat\beta_N^{-4}$. At most a fixed number of stencils of
each admitted degree can meet a site, because anchor occupancy is one;
the finite maximum of their radii and constants gives one site bound,
independent of the number of marks. The original Wilson forced
majorants have, at every fixed larger support exponent, a site sum
bounded by a polynomial in $B_N$ times $\exp(-c\beta_Nt_N^2)$, by the
counting estimate in the proof of Lemma~\ref{lem:D3}. The extra
exponent absorbs the union support. Every new forced component contains
at least one original forced support. Charge half of each such support's
Wilson suppression before applying the rooted tree bound of
Lemma~\ref{lem:D-rooted}; root at a constituent containing the site
being tested, which may be an original forced support rather than a
stencil. Summing the possible site roots and dividing by
$\epsilon_N^{\rm for}$ gives, for fixed stage constants $C_j,c_j>0$,
\[
 (2B_N+1)^4\sup_x\sum_{a:x\in Z_{a,N}}v^{\rm new,for}_{a,N}
 \le C_j(1+B_N)^{C_j}\widehat\beta_N^{C_j}
 \frac{e^{-c_j\beta_Nt_N^2}}{\epsilon_N^{\rm for}}
 \le\frac{\eta_N^{\rm att}}4.
\]
The forced normalization is $\kappa_N=(\epsilon_N^{\rm for})^{1/2}$;
dividing by the smaller $\epsilon_N^{\rm for}\le1$ gives the stronger
upper bound displayed here.
The remaining halves of the suppressions make the forced-support tree
sum convergent; the stencil site bound includes all degrees together.
Since $\beta_Nt_N^2\ge(\cutoffsize_N)^2/2$ and
$\epsilon_N^{\rm for}=e^{-(\cutoffsize_N)^{1/5}}$, the prefactor tends
to zero on the selected late terminal levels. Increasing the entry makes
this contribution and the regular singleton contribution each fit one
quarter of the attachment threshold; the latter uses the strengthened
entry bound~\eqref{eq:E-attachment-allowance}. Their sum fits the
combined site bound in Lemma~\ref{lem:E-source-attachments}. All
estimates precede summation over the finite marker family and are
uniform in the period. Absolute convergence permits the countable
forced-address sums and their regrouping.
\end{proof}

Connected terms with at least two marks need large supports, because distinct marked anchors are far apart. We record this once, since the stopping estimate uses it again.
\begin{lemma}[Marked covering estimate]\label{lem:E-cover}
Let $B\ge8$. If a support $X$ contains $n\ge2$ anchors at pairwise periodic max-distance at least $B^2$, then
\begin{equation}\label{eq:E-marked-cover}
 Bn\le64\cover_{B}(X).
\end{equation}
\end{lemma}
\begin{proof}
Take a minimal cover of $X$ by $M=\cover_B(X)$ radius-$B$ balls with connected union. Join their centers by a tree with edges of length at most $2B+1$. Its doubled traversal has length at most $6B(M-1)$. Assign each anchor to one covering center. Distinct anchors have distinct centers, and the distance between any two assigned centers is at least $B^2-2B$. Visit these centers in their cyclic order in the traversal. The triangle inequality gives $n(B^2-2B)\le6B(M-1)$, so $n(B-2)\le6M$. For $B\ge8$ we have $B-2\ge\tfrac34B$, hence $Bn\le8M\le64M$. The argument uses the periodic max-distance and needs no nonwrapping assumption.
\end{proof}
The estimate permits each source weight to be enlarged by $e^{\delta B/64}$ per mark when $\delta$ units of the support exponent are removed.

The rooted expansion needs a bound on all the other marked factors. We verify it for the primary insertions of Definition~\ref{def:D08}, with their selected coefficient graphs, before taking any replicated integral. For mark~$h$ of degree $d_h$ at level~$i$ let $\mathbf x_{h,i}$ be its dimensionless singleton coefficient, let $b_{h,i}=a_i^{-d_h}|\mathbf x_{h,i}|_1$, and let $v_{a,i}^{\rm src}$ be the weighted integrated majorants of the other marked factors on their complete read supports $Z_{a,i}$, with derivative and support weights as in Lemma~\ref{lem:D-connected}.
\begin{lemma}[Marked attachments]\label{lem:E-source-attachments}
Fix a finite set of degrees $\mathcal D$ and let $\max_{d\in\mathcal D}\reflevel{d}\le k\le N$, with $d\le J_i$ for every $d\in\mathcal D$ and $i\ge k$. Consider finitely many primary insertions of degrees $d_h\in\mathcal D$ on one selected trajectory, whose anchors have level-$i$ ancestors pairwise at periodic distance at least $B_i^2$ for $k\le i\le N$. Give mark~$h$ weights with $\mathsf w_{h,i+1}=\widehat\beta_{i+1}^{-1}\mathsf w_{h,i}$ and $\mathsf w_{h,k}b_{h,k}\le\zeta_{d_h,k}$. Then, for $k\le i\le N$,
\begin{equation*}
 (2B_i+1)^4\sup_x\sum_{a:x\in Z_{a,i}}v_{a,i}^{\rm src}\le\eta^{\rm att}_i/2,
\end{equation*}
uniformly in the cutoff, the period and the number of marks.
\end{lemma}
\begin{proof}
The proof is an induction down the levels: the rooted clause of Lemma~\ref{lem:D-connected} bounds each step, the covering estimate controls the changing mark weights. Closing the bound also requires the higher-mode estimate identified below. The singleton coefficients are controlled by Lemma~\ref{lem:E-marked-step} and their format by Lemma~\ref{lem:E2.2a}.

For the local expansion, give a mode $z_I$ the weight $\prod_{i\in I}w_i$, with $w_i>0$. The resulting coefficient norm is submultiplicative: multiplication sums over disjoint splittings of~$I$, exactly the splittings in the product of the two weighted norms. The coefficientwise rooted clause of Lemma~\ref{lem:D-connected} applies in this algebra, with the attachment bound at the preceding stage of the construction. Derivative and support weights are included in the separate majorants $v_{a,i}$; no submultiplicativity is asserted for them. Connected terms with at least two marks obey~\eqref{eq:E-marked-cover} with $B=B_k$, because distinct anchors are separated by $B_k^2$ and $B_k\ge8$.

The weight change from $\mathsf w_{h,i+1}$ to $\mathsf w_{h,i}$ introduces the factor $\widehat\beta_{i+1}^{|I|}$ on a mode~$I$. For $|I|\ge2$, Lemma~\ref{lem:E-cover} at level~$i$ gives
\begin{equation*}
 \widehat\beta_{i+1}^{|I|}
 \le\exp\!\left\{\frac{64\log\widehat\beta_{i+1}}{B_i}
                        \cover_{B_i}(X)\right\}.
\end{equation*}
The coefficient of the covering number tends to zero at fixed stage, since $B_i\asymp(\log\widehat\beta_i)^2$ and $\widehat\beta_{i+1}/\widehat\beta_i\to1$. The source entry bounds therefore absorb this increase into the extra support exponent at each step. The singleton terms use the coefficient and graph bounds established below. This controls the varying weights on all modes with at least two marks.

The weights $\mathsf w_{h,i}$ estimate the intermediate source polynomials; the weight in the final integral remains $\mathsf w_{h,k}$. The inverse estimate of Lemma~\ref{lem:E-marked-step} gives $b_{h,i+1}\le\ell_i^{d_h} C^{\rm col}_{J_{i+1}}b_{h,i}$. Applying~\eqref{eq:E-source-allowance-transport} to each $d_h$ gives
\begin{equation*}
 \mathsf w_{h,i}b_{h,i}\le\zeta_{d_h,i}\qquad(k\le i\le N).
\end{equation*}
For each $h$ the graph norm has the same bound with its column constant included in $A_{j,d_h}$. The normalized forced terms retain their original eligibility events. Their remaining large-field suppression absorbs the fixed derivative factors before the forced normalization is divided out.

Here is the conversion from a marked-anchor bound to the site sum required by Lemma~\ref{lem:D-connected}. Each support contains its marked anchor. If it also contains a site~$x$ at distance $R$ from that anchor, its connected covering number is at least $R/(2B_i+1)$. One reserved unit of support exponent thus gives a factor at most $e^{-R/(2B_i+1)}$. Distinct anchors have occupancy one, so summing this factor over them is bounded by
\begin{equation*}
 \sum_{z\in\Z^4}e^{-|z|_\infty/(2B_i+1)}\le C(1+B_i)^4.
\end{equation*}
This estimate also bounds the periodic sum, by lifting the distinct anchors to nearest representatives. For a mark of degree $d_h$, the finite derivative products and polynomial evaluations are bounded by $A_{j,d_h}\widehat\beta_i^{q^{\rm ht}_{d_h}}$. Rooted attachment sums use the other reserved support unit to absorb $e^{\eta^{\rm att}_in_\star}$; summing their convergent tree majorants enlarges the bound by at most two, after decreasing $\eta^{\rm att}_i$ if necessary. For the singleton weighted majorants $v_{a,i}^{(1)}$ this gives
\begin{equation*}
 (2B_i+1)^4\sup_x\sum_{a:x\in Z_{a,i}}v_{a,i}^{(1)}
 \le2A_{j,*}(2B_i+1)^4(1+B_i)^4\widehat\beta_i^{-4}\le\eta^{\rm att}_i/4,
\qquad A_{j,*}=\max_{d\in\mathcal D}A_{j,d}.
\end{equation*}
The last inequality is the maximum of the finitely many corresponding
inequalities~\eqref{eq:E-attachment-allowance}; occupancy is used once
for the combined family, not once per degree. The joint terminal
representation of Lemma~\ref{lem:E-joint-terminal} also supplies a
forced higher-mode site bound at the initial level. To complete the
induction, one must propagate a stronger anchored majorant for the
already-formed higher modes and the new joins, at a fine support
exponent that retains the output reserve after reweighting. The
singleton estimate above does not by itself establish that propagation.
The application of the complete attachment bound at later levels is
therefore conditional on this remaining estimate.
\end{proof}

The next estimate is the source of all volume-uniform bounds. Place one copy of the insertion in every cell of the level-$m$ tiling, give each copy a nilpotent mark, and integrate down to a level~$k$ at which every cell contains $(a_m/a_k)^4$ sites. The marks stay at distance at least $B_i^2$ at every intermediate level. The attachment condition for their complete source depends on the higher-mode propagation still required in Lemma~\ref{lem:E-source-attachments}; the following bound is conditional on that input. At level~$k$ each mark contributes the factor $e^6/\mathsf w_k$, where $\mathsf w_k^{-1}=a_k^{-d}e^{o(s_k)}$ is the size of the transported insertion, and each site contributes the unmarked normalization $e^{q_k}$ of Lemma~\ref{lem:E-normalization-cost}. Taking the $|\mathcal C^{\rm tile}_m|$-th root gives the per-insertion bound $\mathsf w_k^{-1}e^{6+(a_m/a_k)^4q_k}$.

We also use the estimate for a residual observable~$E$ with an exact source representation. Call $\mathsf wE$ controlled at level~$k$ when its weighted singleton exponential bound is at most~$3$, its forced singleton bound is at most~$1$, and the connected regular and forced bounds contribute at most one each per mark, with the original eligibility tests. The source evolution from its initial level down to~$k$ must retain these bounds and satisfy the attachment bounds~\eqref{eq:E-attachment-allowance} at every intervening level. This condition concerns the complete source, including its forced terms.
Fix $d$ and a bounded coefficient ball. A stopping pair is a pair of levels $\reflevel{d}\le m\le k$ with $B_k^2\le a_m/a_k$.
\begin{lemma}[Replicated source bound]\label{lem:E-replication}
There are mark weights $\mathsf w_k>0$, independent of~$p$, such that, for each stopping pair, some $N_0(d,\text{ball},m,k)$ gives, for all $N\ge N_0$ and all~$p$,
\begin{equation*}
 \mathsf w_k^{|\mathcal C^{\rm tile}_m|}
 \left|\left\langle\prod_{C\in\mathcal C^{\rm tile}_m}(O(c))_C\right\rangle_{N,p}\right|
 \le e^{|\Lambda_{n_{p,k}}|q_k+6|\mathcal C^{\rm tile}_m|}.
\end{equation*}
For every $\varepsilon'>0$, there is $k_0(\varepsilon')$ with $6-\log\mathsf w_k\le(d+\varepsilon')s_k$ for all $k\ge k_0(\varepsilon')$.
The same integral bound holds with $O(c),\mathsf w_k$ replaced by $E,\mathsf w$ whenever $\mathsf wE$ is controlled at~$k$.
\end{lemma}
The proof bounds the complete source before taking its replicated integral.
\begin{proof}
The marks are attached through Lemma~\ref{lem:E-source-attachments}, the expansion of Lemma~\ref{lem:D-connected} bounds the higher modes, and Lemma~\ref{lem:E-normalization-cost} bounds the unmarked ratio; Lemmas~\ref{lem:E-marked-step} and~\ref{lem:E2.2a} control the singleton coefficients and their format. Enlarge the fixed coefficient radius to include its signed and conjugated images. Put
\begin{equation*}
 E_k=R' C^{\rm col}_{J_k}\prod_{i<k}C^{\rm col}_{J_{i+1}},\qquad
 \mathsf w_k=a_k^d\widehat\beta_k^{-w_d}/E_k,
\end{equation*}
where $R'>0$ depends on the coefficient ball and $w_d$ is a fixed nonnegative integer large enough for its finite polynomial evaluations and forced counterparts. The transport inequality of Lemma~\ref{lem:E2.1}, iterated from $\reflevel{d}$, bounds the level-$k$ coefficients by the product in $E_k$. Threshold powers introduce only a fixed power of $\widehat\beta_k$; enlarging $w_d$ includes them. Increase $w_d$ to at least $q^{\rm ht}_d+4$. Multiplying every mark by $\mathsf w_k$ then bounds its full singleton size by $\widehat\beta_k^{-w_d}\le\zeta_{d,k}$.

The reflected anchors at level~$k$ have spacing $a_m/a_k\ge B_k^2$, and their parents agree under each block map. Since $B_{i+1}^2\le2B_i^2$ and $\ell_i\ge3$, their spacing is at least $B_i^2$ at every finer level. Apply Lemma~\ref{lem:E-source-attachments} with initial weight $\mathsf w_k$. Its full-interval bound verifies the hypothesis of the rooted clause of Lemma~\ref{lem:D-connected}. Use the mode-weighted coefficient norm of that proof; derivative and support weights remain in the separate majorants.

Singleton terms are bounded directly by their one-mark norms. Taylor's product rule introduces at most the factorial of the derivative order; the stage column constant was chosen above this factorial. With $j=J_{k+1}$, the rooted clause of Lemma~\ref{lem:D-connected} supplies the derivative bound at support exponent $\sigma_{J_k}+3$ and the forced bound $v_\star e^{\eta^{\rm att}_kn_\star}e^{-r_j\beta'_u t_{k+1}^2}$ at exponent $\sigma_{J_k}+2$, where $v_\star$ is the integrated majorant of the distinguished factor and $n_\star$ its covering number. Multiplying by a mark weight changes only the distinguished factor in its tree sum. Item~\ref{it:E-columns} bounds the derivative factorial. The weighted sums therefore retain those bounds; decreasing the final exponent from $100$ to any value in $[98,100]$ only decreases the sums. The forced normalization uses half the suppression, leaving the other half for the marked derivative and attachment losses. Their spatial anchor is always one of the marks; a nonzero mode is assigned only to a support containing all its marked anchors. This proves the support assertion in~\eqref{eq:E-source-form}, as well as its absolute convergence. Restricting the set of marks commutes with the construction, since every operation takes disjoint unions of their modes.

At the stopping level, the weighted supremum of each complete singleton exponential term is at most $3$, and that of its forced profile is at most~$1$. Indeed, its polynomial part is at most its coefficient bound, its marked regular sum is at most that bound, and the normalization of the forced part is at most one; Lemma~\ref{lem:E2.1} and the definition of $E_k$ bound the sum by~$3$. The forced profile follows from its $\kappa_k$ normalization. For~\eqref{eq:E-source-form}, the local absolute sums give an additional two per mark: exponentiation of the connected regular terms and the eligible forced sum each contribute at most one. Integrating the remaining scalar-free unmarked exponent and dividing by the unmarked partition function gives at most $e^{|\Lambda_{n_{p,k}}|q_k}$, by Lemma~\ref{lem:E-normalization-cost}. Thus coefficient extraction in the weighted source norm yields
\begin{equation*}
 \mathsf w_k^{|I_0|}\left|\left\langle\prod_{i\in I_0}F_i\right\rangle_{N,p}\right|
 \le\exp\{|\Lambda_{n_{p,k}}|q_k+(3+1+2)|I_0|\}.
\end{equation*}
Exact integration preserves these normalized coefficients, proving the first assertion. The ratio of site and cell counts is $(a_m/a_k)^4$. For a residual source controlled at weight~$\mathsf w$, the same coefficient extraction uses $\mathsf w$ in place of~$\mathsf w_k$; only the bounds $3+1+2$ and the unmarked ratio enter this calculation. Taking the cell root gives
\begin{equation}\label{eq:E-residual-root}
 \|E\|_{\mathcal C^{\rm tile}_m}
 \le \mathsf w^{-1}\exp\{(a_m/a_k)^4q_k+6\}.
\end{equation}
In particular $\mathsf w=c\mathsf w_0$ gives a bound $C_{m,k,\mathsf w_0}/c$, uniformly in the cutoff and period. This proves the residual clause, for which $E$ need not be a primary insertion.

Finally~\eqref{eq:E-log-cost}, summed as in~\eqref{eq:E-running}, gives
\begin{equation*}
 \sum_{i<k}\log C^{\rm col}_{J_{i+1}}=o(s_k),\qquad
 \log\widehat\beta_k=o(s_k).
\end{equation*}
The separate final column is included in this sum for $k>0$. Hence, for every $\varepsilon'>0$, there is $k_0(\varepsilon')$ with $6-\log\mathsf w_k\le(d+\varepsilon')s_k$ for all $k\ge k_0(\varepsilon')$. All choices precede the period, and all the estimates are on full cutoff tails.
\end{proof}


The centered estimate fixes its constant before the physical cell size. A sufficiently fine cutoff then places the fixed insertion stencil strictly inside that cell.
\begin{lemma}[Centered bound]\label{lem:E2.3}
For every $d$ and $R,\varepsilon>0$, there is $A(d,R,\varepsilon)\ge0$ such that, for every $m\ge\reflevel{d}$, some $N_0(d,R,\varepsilon,m)\ge m$ gives, for $N\ge N_0$, every $p,y$ and $\|c\|\le R$,
\begin{equation*}
O_{N,p}^{(d)}(c,y)\text{ is interior},\qquad
\|O_{N,p}^{(d)}(c,y)\|_{\mathcal C^{\rm tile}_m(y)}\le Aa_m^{-d-\varepsilon}.
\end{equation*}
\end{lemma}
\begin{proof}
The replicated source bound of Lemma~\ref{lem:E-replication} gives the seminorm at the stopping level $k=k(m)$ of Lemma~\ref{lem:E1.1}; the choice of $k(m)$ makes the per-cell site cost negligible against $s_k$. For $0<\varepsilon'<1$,
\begin{align}\label{eq:E-centered-input}
 \|O(c,0)\|_{\mathcal C^{\rm tile}_m(0)}
 &\le\exp\{(a_m/a_k)^4q_k+6-\log\mathsf w_k\},\notag\\
 6-\log\mathsf w_k&\le(d+\varepsilon')s_k\quad(k\ge k_0(\varepsilon')),\notag\\
 (a_m/a_{k(m)})^4q_{k(m)}&\le\varepsilon' s_{k(m)}\quad(m\ge m_0(\varepsilon')).
\end{align}
The cutoff threshold is uniform in the coefficient ball and in~$p$.

Put $\tau=s_{k(m)}$, $\upsilon=s_m$ and $W=(a_m/a_k)^4$. Since $q_k\ge1$,
\begin{equation*}
\tau-\upsilon=\log(a_m/a_k)\le W\le Wq_k\le\varepsilon'\tau.
\end{equation*}
The exponent in the first line of~\eqref{eq:E-centered-input} is therefore at most $(d+2\varepsilon')\tau\le(d+2\varepsilon')\upsilon/(1-\varepsilon')$. Choose $\varepsilon'$ so that $(d+2\varepsilon')/(1-\varepsilon')\le d+\varepsilon$. This proves the required estimate with constant one for all sufficiently large~$m$. For the finitely many remaining $m\ge\reflevel{d}$, use the first line of~\eqref{eq:E-centered-input} at the selected fixed $k(m)$ and take the maximum of its ratio to $a_m^{-d-\varepsilon}$ and~one. This fixes~$A$ before~$m$.

Finally the stencil of Definition~\ref{def:D08} has fixed lattice radius. For fixed~$m$ its physical radius tends to zero as $N\to\infty$, uniformly in~$p$. Translation covariance moves the centered calculation to any~$y$.
\end{proof}

\subsubsection{Negligible errors and local matching}\label{sec:partII-E.2.3}

The following estimates concern the insertions of Definition~\ref{def:D08}, with the terminal maps~$\mathsf Z_{N,p,d}$ of Lemma~\ref{lem:E2.1}. They are used to identify the quotient fields in Section~\ref{sec:partII-F}. A statement holds on full cutoff tails if it holds for all $N\ge N_0$ and every~$p$. A nonnegative family $b_{N,p}$ has subpower growth when, for every $\varepsilon>0$, $b_{N,p}\le a_N^{-\varepsilon}$ on such a tail. A coefficient family is tame if
\begin{equation}\label{eq:F-subpower}
 \forall\varepsilon>0\ \exists N_0\ \forall N\ge N_0\ \forall p:
 \|\mathsf Z_{N,p,d}^{-1}c_{N,p}\|_\infty\le a_N^{-\varepsilon}.
\end{equation}
For later volume extraction we also use strictly increasing subsequences $N=\mathsf n(j)$, $p=P(k)$ and the ordered limit
\begin{equation}\label{eq:E2-tails}
 \operatorname*{olim}X_{N,p}=x
 \quad\Longleftrightarrow\quad
 \forall\varepsilon>0\ \exists k_0\ \forall k\ge k_0\ \exists j_0(k)\ \forall j\ge j_0(k):
 |X_{\mathsf n(j),P(k)}-x|<\varepsilon.
\end{equation}
The ordered versions of tameness and negligibility replace full cutoff tails by these quantifiers. We prove the full-tail versions of the bounds below; the same constructions work for ordered-tame families, with their ordered thresholds, and this is the form used in Section~\ref{sec:partII-F}.

An anchor is the lattice site at which an insertion is placed; translates of a bond tiling are by lattice vectors. In a base cell of side~$a_k$, let~$z^a$ be the anchor's integer coordinate measured from the first site in direction~$a$, and put
\begin{equation}\label{eq:E2-margin}
M_{n_*,k,N}=\left\lfloor\frac{a_k}{8(n_*+1)a_N}\right\rfloor,
\qquad M_{n_*,k,N}\le z^a,\quad z^a+M_{n_*,k,N}+1\le a_k/a_N\,.
\end{equation}
These inequalities specify a relative interior margin.

\begin{definition}[Negligible errors]\phantomsection\label{def:E2-small}
An observable family~$E_{N,p}(y)$ is negligible if there are $\eta_{N,p}\ge0$ with $\sup_p\eta_{N,p}\to0$ as $N\to\infty$ such that, for every arity~$n_*$ and level~$m$, there are $k\ge m$ and $N_0\ge k$ with the following property: for all $N\ge N_0$, all~$p$, every translated bond tiling~$\mathcal C^{\rm tile}$ of side~$a_k$ and every base-cell anchor~$y$ satisfying~\eqref{eq:E2-margin},
\begin{equation*}
E_{N,p}(y)\text{ is interior},\qquad
\|E_{N,p}(y)\|_{\mathcal C^{\rm tile}}\le\eta_{N,p}\,.
\end{equation*}
\end{definition}
The tolerance is chosen before $n_*$ and~$m$; the level and the cutoff tail may depend on these two requests. Here $n_*$ is the number of insertions in a later product test, and the lower level~$m$ requires the chosen cell scale to be arbitrarily fine. All residuals below are translation covariant, and their signed and conjugated images obey the same estimates. To pass from centered to arbitrary interior seminorms, choose $k'\ge k$ with $24(n_*+1)a_{k'}\le a_k$ and apply the finite reflection argument of Lemma~\ref{lem:F1a}. It uses only fixed stencils and the chessboard inequality, so it also applies to these residuals on full cutoff tails.

Local matching converts a local observable of controlled size into an insertion of the same degree. It makes no assumption on the size of the physical coefficient it produces; it bounds its inverse terminal coordinate, which is the quantity in~\eqref{eq:F-subpower}. The derivative bound below controls Taylor extraction, and the supremum bound controls the forced complement.
\begin{lemma}[Local comparison]\phantomsection\label{lem:E-local-comparison}
\begin{enumerate}
\item[(i)] \emph{Local matching.} Let $F_{N,p}(y)$ be local observables such that $a_N^{d'}F_{N,p}$ is continuous on a fixed stencil and regular on healthy configurations in the sense of Definition~\ref{def:D05a}, with derivative bounds and a supremum bound over all configurations of subpower growth. Then there are tame $b_{N,p}\in V_{d'}$ with $O^{(d')}(b)-F$ negligible.
\item[(ii)] \emph{Auxiliary equation insertions.} For the auxiliary section of Lemma~\ref{lem:A-R4}, there are invertible $P_{N,p}$, with $\|P_{N,p}^{-1}\|_1\le2$ eventually uniformly in~$p$, such that, with $B_{d,N,p}=\mathsf Z_{N,p,d}P_{N,p}t_N^{\factorcount}\mathsf R_d$, the comparison~\eqref{eq:E-equation-comparison} holds for every equation label $\pi_{N,p}\in\mathcal I_d$ with tame $B_{d,N,p}\pi_{N,p}$.
\end{enumerate}
\end{lemma}
Here $\mathcal I_d$ is the equation ideal described in the proof of Lemma~\ref{lem:E2.6}, and the comparison is
\begin{equation}\label{eq:E-equation-comparison}
 O^{(d)}_{N,p}(B_{d,N,p}\pi_{N,p},y)
 -a_N^{-d}V_d^{\rm aux}[\mathsf R_d\pi_{N,p}](y)
 \quad\text{is negligible}.
\end{equation}
Both parts use the same backward coefficient correction.
\begin{proof}
We represent the observable as a one-mark source at the cutoff, carry the difference with the selected primary insertion backwards with the marked step (Lemma~\ref{lem:E-marked-step}), and bound its replicated seminorm by Lemma~\ref{lem:E-replication}. Write $h_N=a_N^{d'}F_{N,p}(0)$. Enlarge its stencil by the section stencil. The radial Taylor coefficients through degree~$d'$ (in the sections of Lemma~\ref{lem:A1}) give a coefficient $v_N$ and a remainder. They are bounded by a fixed multiple of the derivative bound: it is a linear map on a finite jet space. Subtract the section with that coefficient on the healthy part. Its regular remainder has zero value and zero radial jet through degree~$d'$. Assign the complementary observable to the forced supports meeting the enlarged stencil. This gives an exact one-mark representation of $\density\rho_N h_N$, with
\begin{equation}\label{eq:E-local-start}
 |v_N|_1+r_N\le E_{N,p},\qquad
 \text{forced majorant}\le E'_{N,p},
\end{equation}
where $E,E'$ have subpower growth. The support weight of this fixed enlarged stencil is bounded by $e^{81\sigma_{J_N}}$ eventually; its $B_N$ covering number is at most~$81$. Threshold powers, the Taylor-coefficient bound and this support weight are included in $E,E'$. The entry conditions make each of them subpower. No division by the forced normalization is made on the uncontrolled complement: it is retained with its original eligibility tests and the supremum bound $E'$.

To compare with the selected graph, carry this representation backwards and at the reference level choose the primary coefficient with the same polynomial coordinate. Subtract the corresponding selected inserted density. Its polynomial coordinate at the reference level is zero. The following calculation carries this matching condition towards the fine scale.

For one step, let $u,v$ be the coarse and fine coordinates, $R,T$ their remainder bounds, $Y=|\mathsf S_Lv|_1$, $K$ the norm bound for the leading map and its inverse, $\varepsilon_L=L^{-3/8}$ and $\vartheta_L=L^{1/8}$. The marked-step inequalities give
\begin{equation}\label{eq:E-match-cone}
 Y\le K\bigl(|u|_1+\varepsilon_L(Y+T)\bigr),\qquad R\le KY+\varepsilon_L T.
\end{equation}
Suppose $a|u|_1\le R$, with $a=32(1+C^{\rm col}_{J_i})^2\ge4K^2$. Then $4K|u|_1\le Y+\varepsilon_L T$. Substitution in the first inequality, using $4K\varepsilon_L\le1$, gives $Y\le4K\varepsilon_L T$. Item~\ref{it:E-match} in Lemma~\ref{lem:E1.1} gives
\begin{equation*}
 4K\,32(1+C^{\rm col}_{J_{i+1}})^2\varepsilon_L\le1,\qquad
 (4K^2+1)\vartheta_L\varepsilon_L\le1.
\end{equation*}
Consequently $32(1+C^{\rm col}_{J_{i+1}})^2|v|_1\le T$ and $\vartheta_L R\le T$. The matching condition starts this induction with $u=0$ at the reference level. Multiplying the second inequality through the interval gives the inverse-scale gain, while the finest-level bound is closed by absorbing the coefficient contribution into its remainder bound. In particular, with subpower $E''$ absorbing $E,E'$,
\begin{equation}\label{eq:E-backward-match}
 |\Delta_N|_1\le\frac{2(1+C^{\rm col}_{J_N})E''_{N,p}}{\ell_{N-1}^{1/8}},\qquad
 R_k\le2(1+C^{\rm col}_{J_N})E''_{N,p}e^{-(s_N-s_k)/8}.
\end{equation}
Here $\Delta_N$ is the difference from the selected fine polynomial coordinate. The same bounds hold when a local source is first represented at any level $n\le N$, with its polynomial coordinate matched to the selected graph at a lower reference level. If $E''_{n;N,p}$ bounds its starting data after the coefficient contribution is absorbed, then
\begin{equation}\label{eq:E-backward-match-level}
 |\Delta_n|_1\le\frac{2(1+C^{\rm col}_{J_n})E''_{n;N,p}}{\ell_{n-1}^{1/8}},\qquad
 R_k\le2(1+C^{\rm col}_{J_n})E''_{n;N,p}e^{-(s_n-s_k)/8}
 \quad(\reflevel{d'}\le k<n).
\end{equation}
Indeed the cone starts at the same zero matched coordinate, and its induction uses only the steps through~$n$. The boundary absorption uses the starting bound there. The product $\prod_{i=k}^{n-1}\ell_i^{-1/8}$ is exactly $e^{-(s_n-s_k)/8}$. Taking $n=N$ recovers~\eqref{eq:E-backward-match}. The factor $e^{-(s_n-s_k)/8}$ is the inverse-scale gain: each matched step contracts the remainder by $\ell^{-1/8}$.

To estimate the replicated residual, for a source of size $E$ started at $N$, use at level~$k$ the mark weight
\begin{equation*}
 \mathsf w=\frac{\zeta_{d',k}/4}{a_k^{-d'}E\exp\{-(s_N-s_k)/8\}},
\end{equation*}
where $\zeta_{d',k}=\widehat\beta_k^{-(q^{\rm ht}_{d'}+4)}$ is the source threshold fixed before~$N,p$ in~\eqref{eq:E-attachment-allowance}. Its reciprocal is a fixed power of the reference coupling. At every finer step multiply this weight by the source contraction factor of that step. Its product with the dimensionless singleton size is at most one quarter of that source threshold; the other three quarters bound the graph remainder, the forced complement and their sum. This follows inductively from the two inequalities in~\eqref{eq:E-atom-estimates}. For the first step the forced complement is bounded by $E'$ in~\eqref{eq:E-local-start}; the large-field suppression absorbs its fixed support and derivative factors before division by the new forced normalization. The radius transport~\eqref{eq:E-source-allowance-transport} and the singleton site-sum calculation in Lemma~\ref{lem:E-source-attachments} control the singleton attachments. The complete residual is controlled at level~$k$ in the sense of Lemma~\ref{lem:E-replication} only after the higher-mode propagation identified there has been established. The singleton polynomial coordinate at the reference level is zero, so the representation constructed this way is exactly the difference with the selected primary coefficient.

Here $m$ labels the tiling and $k$ the finer stopping level at which its replicated sources are integrated. Replicate the difference at a stopping pair $m\le k$, and use the weighted source representation~\eqref{eq:E-source-form}. Coefficient extraction contributes $\mathsf w^{-1}$ per mark; the complete singleton bounds and the normalized unmarked ratio~\eqref{eq:E-vacuum-ratio} contribute a constant depending on $m,k$. Consequently, for every fixed such pair,
\begin{equation}\label{eq:E-local-error}
 \|O^{(d')}(b)-F\|_{\mathcal C^{\rm tile}_m(0)}
 \le C_{m,k}\mathfrak e_{N,p}a_N^{1/8},\qquad
 \|\mathsf Z_{N,p,d'}^{-1}b_{N,p}\|_\infty\le2E'''_{N,p},
\end{equation}
where $\mathfrak e,E'''$ have subpower growth. The factor $e^{s_k/8}$ and the fixed weight factor $\zeta_{d',k}^{-1}$ belong to $C_{m,k}$. The same argument applies after each signed permutation and conjugation, with translation covariance and the union of their fixed stencils. Absorbing two factors $a_N^{-1/32}$ bounds the first expression by $a_N^{1/16}$ on every fixed pair. The interior-seminorm comparison after Definition~\ref{def:E2-small} gives the common tolerance $a_N^{1/32}$, chosen before all cell requests. This proves negligibility and tameness on full cutoff tails, and part~(i).

For part~(ii), perform the construction on all its basis polynomials together, with the original Taylor normalization. Let $u_j$ be the primary bare coefficient matching its $j$th basis vector, and set $Px=\sum_jx_ju_j$, after the finite symmetry averaging. For this fixed section the regular derivative bound $H_{j,d}$ at stage~$j$ is fixed before its ratio, independently of the later coupling. The sharper coefficient estimate in~\eqref{eq:E-backward-match} therefore has numerator bounded by $2(1+C^{\rm col}_j)H_{j,d}$ on each basis input. The ratio condition, including the finite basis-recombination factor, gives
\begin{equation*}
 \left|x-Px\right|_1\le\tfrac14|x|_1.
\end{equation*}
Indeed, before summing over the fixed basis the right side is the last-step factor $\ell_{N-1}^{-1/8}$ times the stage coefficient bound; the stage ratio was chosen after that bound. The matrix is therefore invertible and $|P^{-1}x|_1\le2|x|_1$. The construction is common to all later cell requests. Multiplication by $t_N^{\factorcount}\mathsf R_d$ changes a polynomial label into its auxiliary bare insertion coordinate; multiplication by $\mathsf Z$ changes that coordinate to the reference normalization. These are exactly the maps $B$ in the statement. The inverse bound preserves tameness, and~\eqref{eq:E-local-error} compares the two sections on its transported equation ideal. Finite combinations use the same subpower factors. The argument only replaces eventual bounds on coefficients, so ordered-tame families obey the corresponding ordered version as well.
\end{proof}


\subsubsection{Equation identities and coefficient corrections}\label{sec:partII-E.2.4}

The equation projections separate the coefficient directions controlled by exact cutoff identities. The free equation space consists of polynomial labels that vanish on curvature jets satisfying the Maxwell equations. Its transported image is the kernel of the projections chosen below, and the free quotient has dimension~$\dim_\C\maxwellspace{d}$. The projections and the equation observables in Lemmas~\ref{lem:E2.6}--\ref{lem:E2.7} are chosen together, before the spacing and volume subsequences.
\begin{lemma}[Equation projections]\label{lem:E2.6}
For every~$d$, there are complex linear idempotents $\Pi_{d,N,p}:V_d\to V_d$ and $C^\Pi_d<\infty$ such that
\begin{equation*}
\|\Pi_{d,N,p}\|\le C^\Pi_d,\qquad
\dim\operatorname{ran}\Pi_{d,N,p}\le\dim_\C\maxwellspace{d},\qquad
\ker\Pi_{d,N,p}=B_{d,N,p}\mathcal I_d,
\end{equation*}
with the maps $B$ of Lemma~\ref{lem:E-local-comparison}(ii).
\end{lemma}
\begin{proof}
The kernel is prescribed by the transported equation ideal, and the bound comes from choosing orthogonal projections in a fixed norm. Let $\mathcal I_d\subset V_d$ be the polynomials of Definition~\ref{def:D07} belonging to the ideal generated by the linear Maxwell equations in radial jet coordinates. Radial coordinates already satisfy the linear Bianchi identities. Taking invariant polynomials commutes with the on-shell quotient: average any representative over the compact adjoint group to obtain an invariant representative. Thus the on-shell map is onto $\maxwellspace{d}$ (Definition~\ref{def:D11}) and
\begin{equation}\label{eq:E-equation-count}
\dim V_d-\dim\mathcal I_d=\dim_\C\maxwellspace{d}.
\end{equation}

Lemma~\ref{lem:E-local-comparison}(ii) constructs invertible coefficient maps $B_{d,N,p}:V_d\to V_d$. They transport $\mathcal I_d$ to the coefficient equations for auxiliary local representatives, including their change to the selected insertion coordinates. Set $U_{d,N,p}=B_{d,N,p}\mathcal I_d$. In the Euclidean norm on the fixed basis take the orthogonal projection $Q_{d,N,p}$ onto $U_{d,N,p}$ and set $\Pi_{d,N,p}=1-Q_{d,N,p}$. Then $\Pi^2=\Pi$ and $\ker\Pi=U_{d,N,p}$. Equivalence with the fixed coefficient norm bounds $\|\Pi\|$ independently of $N,p$, even when $B$ is poorly conditioned. Since $B$ is invertible, $\dim\ker\Pi=\dim\mathcal I_d$, and~\eqref{eq:E-equation-count} gives the dimension bound.

The same local comparison maps will be used in Lemma~\ref{lem:E2.7}. They and these projections are chosen for all cutoffs before any subsequences.
\end{proof}

The next estimate converts a coefficient equation into an observable with exact separated moment identities.
\begin{lemma}[Equation observables]\label{lem:E2.7}
For these same projections, every tame family~$c_{N,p}\in\ker\Pi_{d,N,p}$ admits a measurable observable family~$\mathcal Q_{N,p}(y)$, bounded at each anchor and cutoff on a full cutoff tail, for which $O^{(d)}_{N,p}(c_{N,p},y)-\mathcal Q_{N,p}(y)$ is negligible. For every fixed spectator list and~$\delta>0$, on a full cutoff tail,
\begin{equation*}
\int \mathcal Q_{N,p}(y_i)\prod_{h\ne i}O_{N,p}(v_h,y_h)\,d\mu_{N,p}=0
\qquad(\operatorname{sep}_{\rm tor}(y)\ge\delta)\,.
\end{equation*}
\end{lemma}
\begin{proof}
The observable is the auxiliary equation section of Lemma~\ref{lem:A-R4}, whose separated moments vanish exactly; Lemma~\ref{lem:E-local-comparison}(ii) compares it with the insertion. Lemma~\ref{lem:A-R4} supplies $\widetilde P_d=V_d^{\rm aux}$ and $\mathsf R_d$. Using the kernel identity of Lemma~\ref{lem:E2.6}, write $c_{N,p}=B_{d,N,p}\pi_{N,p}$ with $\pi_{N,p}\in\mathcal I_d$, and set
\begin{equation}\label{eq:E-equation-section}
 \mathcal Q_{N,p}(y)=a_N^{-d}V_d^{\rm aux}[\mathsf R_d\pi_{N,p}](y).
\end{equation}
Lemma~\ref{lem:E-local-comparison}, specifically~\eqref{eq:E-equation-comparison}, proves that $O^{(d)}(c)-\mathcal Q$ is negligible. Its inverse change has subpower growth for tame~$c$. The common tolerance and interior-seminorm comparison were proved there and apply to this family.

The exact zero identity is~\eqref{eq:A-aux-zero} in Lemma~\ref{lem:A-R4}: every auxiliary equation term is a Wilson action derivative whose multiplier and transport avoid the varied link, so Haar invariance makes its exterior moment zero. The auxiliary section has one finite equation radius. At prescribed positive physical separation, every fixed spectator stencil is exterior to it for all sufficiently fine cutoffs, uniformly in~$p$. Its product is bounded and measurable, as required by that lemma. Apply the identity with that product as spectator. Finally $\widetilde P_d$ is continuous on a finite product of compact groups. At each fixed cutoff its evaluation with the chosen finite coefficients is bounded.
\end{proof}

Changing the degree changes the insertion normalization. The next estimate corrects that change before taking limits.
\begin{lemma}[Corrected inclusion]\label{lem:E2.8}
For every $d\le d'$ and $c\in V_d$, there is a tame family~$b_{N,p}\in V_{d'}$ such that
\begin{equation*}
O^{(d')}_{N,p}(b_{N,p},y)-O^{(d)}_{N,p}(c,y)
\quad\text{is negligible}\,.
\end{equation*}
\end{lemma}
\begin{proof}
The target observable satisfies the hypotheses of local matching (Lemma~\ref{lem:E-local-comparison}(i)) in degree~$d'$. After multiplication by $a_N^{d'}$, it is
\begin{equation*}
a_N^{d'-d}P_d[t_N^{-\factorcount}\mathsf Z_{N,p,d}^{-1}c](y).
\end{equation*}
It is a local continuous function on the fixed degree-$d$ stencil. Its ordered derivatives in the allowed link directions and its supremum have subpower bounds. Indeed, the inverse terminal coordinate has this growth by the first transport inequality of Lemma~\ref{lem:E2.1} and~\eqref{eq:E-log-cost}; threshold powers have subpower growth because $\log\widehat\beta_N=o(s_N)$. The factor $a_N^{d'-d}\le1$ can only decrease these bounds.

Local matching extracts a degree-$d'$ coefficient $b_{N,p}$ whose inverse terminal coordinate still has subpower growth. Its residual seminorm is bounded by $C_{m,k}\mathfrak e_{N,p}a_N^{1/8}$ at every fixed stopping pair of levels, with $\mathfrak e$ of subpower growth. The absorption and interior-seminorm argument in the proof of Lemma~\ref{lem:E-local-comparison} gives one tolerance $a_N^{1/32}$ for all the later cell requests. The residual is therefore negligible. For $d=d'$ one may simply take $b_{N,p}=c$, with zero residual.
\end{proof}

The physical difference quotient requires a correction in the next degree.
\begin{lemma}[Corrected difference]\label{lem:E2.9}
For every degree~$d$, $c\in V_d$ and axis~$a$, there is a tame family~$b_{N,p}\in V_{d+1}$ such that
\begin{equation*}
O^{(d+1)}_{N,p}(b_{N,p},y)
-\frac{O^{(d)}_{N,p}(c,y+e_a)-O^{(d)}_{N,p}(c,y)}{a_N}
\quad\text{is negligible}\,.
\end{equation*}
\end{lemma}
\begin{proof}
Multiplying the target difference by $a_N^{d+1}$ removes the inverse spacing and gives exactly
\begin{equation*}
P_d[t_N^{-\factorcount}\mathsf Z_{N,p,d}^{-1}c](y+e_a)
-P_d[t_N^{-\factorcount}\mathsf Z_{N,p,d}^{-1}c](y),
\end{equation*}
to which local matching (Lemma~\ref{lem:E-local-comparison}(i)) applies in degree $d+1$. Its stencil is the union of one fixed stencil and its unit translate. It is continuous and local. Moving an ordered derivative anchor by one lattice site multiplies its spatial weight by at most a stage-dependent factor. The factor is at most $2^{\varpi_{J_N}}$ per derivative. The logarithmic condition $\varpi_{J_N}=o(s_N)$ in Lemma~\ref{lem:E1.1} makes each fixed product of these factors subpower. The other bounds are those in the preceding proof. Local matching in degree $d+1$ now supplies a tame coefficient family and the same negligible residual. The inverse spacing has been removed before matching, so it causes no further loss in that residual. The four axes use translates of this one computation and their fixed spatial-weight bounds.
\end{proof}

\subsubsection{Comparison with free fields}\label{sec:partII-E.2.5}

Fix~$d$ and a real basis $\mathfrak m_1=1,\mathfrak m_2,\ldots,\mathfrak m_{\dim_\C\maxwellspace{d}}$ of~$\maxwellspace{d}$. Choose finitely many one- or two-point tests. In test~$i$, the first label is the variable~$\mathfrak m_j$; the optional second label is a fixed basis element. Let $y^i\in\Gen{n_i}$, $n_i\in\{1,2\}$, be the test's configuration. The free test matrix is
\begin{equation}\label{eq:E2-free-matrix}
\mathcal G_{ij}=\maxwellcorr_{n_i}(\mathfrak m_j,\text{ test }i\text{ spectator};y^i)\,.
\end{equation}
The next lemma supplies a finite family that distinguishes every free field of degree at most~$d$.
\begin{lemma}[Free tests]\label{lem:E2-free-tests}
For every~$d$ and real basis~$\mathfrak m_1=1,\mathfrak m_2,\ldots,\mathfrak m_{\dim_\C\maxwellspace{d}}$, there are finitely many one- or two-point tests at generic configurations, with every second label a basis element, for which the matrix~\eqref{eq:E2-free-matrix} has rank~$\dim_\C\maxwellspace{d}$.
\end{lemma}
The proof uses only the Maxwell correlations of Definition~\ref{def:D11} and the free positivity in Lemma~\ref{lem:G1.1}, independently of Section~\ref{sec:partII-F}.
\begin{proof}
Let~$c\ne0$ in~$\maxwellspace{d}$. Lemma~\ref{lem:G1.1} gives a strictly positive reflected two-point value for~$c$ and the reflection of~$\bar c$. Expand that second label in the basis~$(\mathfrak m_j)$. At least one of the resulting two-point functionals is nonzero on~$c$. Each free Wick correlation is a finite sum of products of derivatives of $\Gamma(x)=(4\pi^2|x|^2)^{-1}$ between distinct sites, hence is smooth there. This permits a perturbation of the pair to~$\Gen{2}$ preserving this nonzero value. Thus the common kernel of all permitted test functionals is zero. Starting with~$\maxwellspace{d}$, choose a test nonzero on a nonzero vector in the current common kernel. Each choice strictly lowers its dimension. After finitely many choices the common kernel is zero, so~$\mathcal G$ has rank~$\dim_\C\maxwellspace{d}$.
\end{proof}
Fix such a family for the rank comparison below.

For an index~$\nu$ of the comparison sequence, a common coefficient list~$b^\nu_{j,N,p}\in V_d$ and lattice-compatible configuration $x^{\nu,i}=r_\nu y^{\nu,i}$ define the test insertions $F_h^{\rm tst}=O^{(d)}_{N,p}(b^\nu_{j_h,N,p},x_h^{\nu,i}/a_N)$, with~$j_1=j$ and the spectator index fixed by test~$i$. Physical points are read modulo the lattice period. Let~$\mathcal G_J^\nu$ be the Maxwell moment of the labels~$(\mathfrak m_{j_h})_{h\in J}$ at~$(y_h^{\nu,i})_{h\in J}$, with $\mathcal G_\varnothing^\nu=1$. Define the relative errors by
\begin{equation}\label{eq:F-relative}
R_\varnothing=1,\qquad
R_I=\prod_{h\in I}F_h^{\rm tst}-
\sum_{\varnothing\ne J\subseteq I}\mathcal G_J^\nu R_{I\setminus J}\,.
\end{equation}
The rank argument uses only the singleton and pair cases. Suppressing the superscript~$\nu$, these are
\begin{equation*}
\begin{aligned}
R_{\{1\}}&=F_1^{\rm tst}-\mathcal G_{\{1\}},\\
R_{\{1,2\}}&=(F_1^{\rm tst}-\mathcal G_{\{1\}})(F_2^{\rm tst}-\mathcal G_{\{2\}})
-(\mathcal G_{\{1,2\}}-\mathcal G_{\{1\}}\mathcal G_{\{2\}}).
\end{aligned}
\end{equation*}
Thus the pair error subtracts the connected free two-point value from the product of the centered insertions. The lifts~$b^\nu$ include the scale normalization, so~$F_h^{\rm tst}$ carries no additional power of~$r_\nu$. A lift is the reference coefficient whose Wick-ordered leading polynomial at scale~$r_\nu$ is the specified basis element~$\mathfrak m_j$. The construction uses $r_\nu=a_{t_\nu}(\cutoffsize_{t_\nu})^b$ for one fixed $b>0$ chosen from the degree and the finite tests, before~$\nu$ varies.

The source estimates below also apply to tests with any fixed arity bound~$n\ge2$. For a slot set $I_0$ with $|I_0|\le n$, define the free generating polynomial and its logarithm, the free cumulant generating function,
\begin{equation*}
 \mathsf G^{\rm Max}(z)=\sum_{J\subseteq I_0}\mathcal G_Jz_J,
 \qquad \mathsf C^{\rm Max}(z)=\log\mathsf G^{\rm Max}(z),
 \qquad z_i^2=0\quad(i\in I_0).
\end{equation*}
The logarithm is a finite polynomial because the constant term of $\mathsf G^{\rm Max}$ is one and its nonconstant part is nilpotent. Its coefficient on a nonempty~$I$ is
\begin{equation*}
 [z_I]\mathsf C^{\rm Max}
 =\sum_{\mathcal P\text{ partition of }I}
 (-1)^{|\mathcal P|-1}(|\mathcal P|-1)!
 \prod_{J\in\mathcal P}\mathcal G_J.
\end{equation*}
For two slots these are the two singleton means and the connected pair moment; the rank comparison uses just this case. In the source estimates, the test positions may range over any fixed compact set $K\subset\Conf{|I_0|}$ with $\sepinf\ge\delta>0$. Constants may depend on $K,\delta,n$ and the finite degree and derivative orders; the lifts are common to these positions. The coefficients to be estimated are
\begin{equation}\label{eq:E-corrected-coefficients}
 \mathfrak c_I=[z_I]\left(e^{-\mathsf C^{\rm Max}(z)}\prod_h(1+z_hF_h^{\rm tst})\right).
\end{equation}

For the source comparison, use the intervals~\eqref{eq:E-window}. Fix $b>0$ large enough for the degree and the finite tests, put $l=t-h_t(b)$, $k=l-v_t$ and $r=a_t(\cutoffsize_t)^b$. The interval from $t$ to~$l$ separates the reading points from the last fluctuation scale; the interval from $l$ to~$k$ lets the mark weights of the corrected sources increase while their remainders decrease. With $L=\ell_{J_t}^{\rm st}$, these intervals satisfy
\begin{equation}\label{eq:E-window-geometry}
 L^{h_t(b)}\ge(\cutoffsize_t)^{b+2},\qquad
 L^{v_t}\ge B_t^{10},\qquad
 \widehat\beta_i\in[\widehat\beta_t/2,\widehat\beta_t]\quad(k\le i\le t)
\end{equation}
on sufficiently late intervals. The last assertion follows from $h_t(b)+v_t=o(t)$ at fixed stage and the entry choices. All rounded test points lie strictly inside the level-$k$ cell: their distance from the origin is $O(r)$, while $a_l/r\ge(\cutoffsize_t)^2$. Rounding changes their rescaled coordinates by $O((\cutoffsize_t)^{-b})$.

The comparison may use a different sequence of positive couplings below
the reading level~$t$. The next lemma records what is preserved when
the \emph{entire} marked density is pushed along that sequence.
\begin{lemma}[Changing the comparison course]\label{lem:E-course-mass}
Fix a cutoff, a period and a finite set of squarefree marks. At level~$i$
let $\lambda_i$ be product Haar measure, and let
$K_i^b(U,dV)=k_i^b(U,V)\lambda_i(dV)$ be the block probability kernel
from a fine field~$U$ at level~$i+1$ to a coarse field~$V$ at level~$i$,
using the positive coupling $b_{i+1}$. Suppose every coefficient of
the terminal marked density $H_N(z;U)$ is integrable with respect to
$\lambda_N$. Its empty coefficient is $\rho_N\ge0$, with
$Z=\int\rho_N\,d\lambda_N>0$. Define $H_m^b$ by pushing the whole
polynomial coefficientwise from~$N$ to~$m$ through these kernels.

For every mode~$I$, the integral $\int[z_I]H_m^b\,d\lambda_m$ equals
$\int[z_I]H_N\,d\lambda_N$, independently of the positive course~$b$.
If two courses use the same kernels from~$N$ down to~$t$, their marked
densities at~$t$ agree $\lambda_t$-almost everywhere. Multiplication by
a field-independent squarefree polynomial $Q(z)$ commutes with every
push; if $Q(0)=1$, it also leaves the empty coefficient and its integral
$Z$ unchanged.

If $[z_I]H_N=\rho_NF_I$, define the joint path law by drawing $U_N$ from
$Z^{-1}\rho_N\lambda_N$ and then drawing $U_i$ from
$K_i^b(U_{i+1},\cdot)$ in descending order. For
$F_I\in L^1(\rho_N\lambda_N/Z)$, the conditional expectation of
$F_I(U_N)$ given $U_m=V$ is
\[
 \frac{[z_I]H_m^b(V)}{[z_\varnothing]H_m^b(V)}
 \qquad\text{for }\lambda_m\text{-almost every }V\text{ with }[z_\varnothing]H_m^b(V)>0.
\]
\end{lemma}
\begin{proof}
For an integrable coefficient $h$, the push from $i+1$ to~$i$ is
$V\mapsto\int h(U)k_i^b(U,V)\,d\lambda_{i+1}(U)$. Fubini's theorem and
$\int k_i^b(U,V)\,d\lambda_i(V)=1$ show that this push is integrable
and preserves its integral. Iteration proves the first assertion;
iteration with identical kernels proves the equality at~$t$. Multiplication
by a fixed polynomial $Q$ is a fixed linear map on the finite coefficient
space, so it commutes with the push. Its unit empty coefficient preserves
the unmarked density.

For the conditional statement, the law of~$U_m$ has density
$[z_\varnothing]H_m^b/Z$. For every bounded measurable test~$\varphi$,
Fubini gives
\[
 \mathbb E\bigl[\varphi(U_m)F_I(U_N)\bigr]
 =Z^{-1}\int\varphi(V)[z_I]H_m^b(V)\,d\lambda_m(V).
\]
The pushed signed or complex measure in this identity is absolutely
continuous with respect to the pushed unmarked measure. Thus the ratio
above is its Radon--Nikodym derivative; its value on the null set where
the denominator vanishes is immaterial.
\end{proof}
The lemma preserves the original normalized Wilson integrals, including
when $Q=e^{-\mathsf C^{\rm Max}}$ is the finite squarefree correction.
It does not identify the lower-level densities or conditional laws of
two different courses. A lawful representation tied to the chosen
course, and the guarded source estimates for that representation, are
separate inputs below.

Fix the arity bound $n\ge2$ and the compact separated position set specified above. Put $\mathcal M=\cutoffsize_t$. Let $q_{\rm der}$ bound the required derivative orders, and let $c,\delta_{\rm G}>0$ be the localization and healthy-chart exponents in the estimates below. Define
\begin{equation}\label{eq:E-source-initial}
 \mathcal E_t^{\rm G}=\mathcal M^{-b}
 +C^{\rm G}_j\mathcal M^{q_{\rm der}}
   (\widehat\beta_t^{-\delta_{\rm G}}+e^{-cB_t}).
\end{equation}

\paragraph{Comparison windows.} A comparison window is a triple of levels $k<l<t$; put $\mathcal M=\cutoffsize_t$. The proofs of Lemmas~\ref{lem:E-source-polynomial}--\ref{lem:E-source-integral} use the window only through the following conditions:
\begin{enumerate}
\item[(Win1)] $\log(a_l/a_t)\ge(b+2)\log\mathcal M$, $\log(a_k/a_l)\ge10\log B_t$, and $\widehat\beta_i\in[\widehat\beta_t/2,\widehat\beta_t]$ for $k\le i\le t$;
\item[(Win2)] $\mathcal C_{J_t}\prod_{k\le i\le t}\mathcal C_{J_i}\le\mathcal M^{1/64}$;
\item[(Win3)] every level of $[k,t]$ is past the entries of item~\ref{it:E-entries} for the finitely many requests, and $Q^{\rm cmp}_{J_k}\le\log\mathcal M/(16n)$.
\end{enumerate}
Late windows~\eqref{eq:E-window} satisfy (Win1)--(Win3) by~\eqref{eq:E-window-geometry}, \eqref{eq:E-source-loss} and Lemma~\ref{lem:E1.1}. A window need not lie in one stage: under (Win1)--(Win3) the three lemmas hold with $\mathcal C_j^{h_t(b)+2}$ replaced by $\mathcal C_{J_t}\prod_{l\le i\le t}\mathcal C_{J_i}$.

The polynomial comparison supplies the initial error for the subsequent source iteration.
\begin{lemma}[Middle-level source comparison]\label{lem:E-source-polynomial}
For each fixed finite label family, sufficiently late intervals admit common real lifts, fixing the unit and bounded in $N,p$ for that interval. For every comparison window satisfying (Win1)--(Win3) and uniformly over the prescribed compact position set, the coefficients of $z_I$, $I\ne\varnothing$, of $e^{-\mathsf C^{\rm Max}}\prod_h(1+z_hF^{\rm tst}_h)$, represented at level~$l$, have total regular and normalized forced bound
\begin{equation*}
 \mathcal C_j^{h_t(b)+2}\mathcal E_t^{\rm G}.
\end{equation*}
The assertion holds on a full cutoff tail, uniformly in~$p$.
\end{lemma}
The proof gives the Wick correction and the inverse-separation gain. Its
application to the actual pushed source still needs a joint-law estimate:
the sampled physical reads, forecast, guard, determinant and residual
loading must be matched to the same sector law for every mode, including
the empty mode. The resulting value and derivative estimates must give
the stated regular and normalized forced bounds with the original tests.
In addition, a comparison based on a restarted coupling course below~$t$
requires an exact density and coupling tie to that course; an estimate
on the selected running history would instead need its own proof. Until
one of these suppliers is established, this lemma and its applications
below are conditional on those inputs.
\begin{proof}
The free covariance is compared with the lattice covariance of Lemmas~\ref{lem:B2} and~\ref{lem:B-windows} at the reading scale, the internal contractions are removed by Wick ordering, and the connected expansion of Lemma~\ref{lem:D-connected} carries the corrected sources from $t$ to~$l$, with the marked step of Lemma~\ref{lem:E-marked-step} supplying the singleton coefficients. The sources have at most $n$ marks in each cell; the number of cells is unrestricted. Use the coefficientwise rooted clause of Lemma~\ref{lem:D-connected} in their weighted nilpotent algebra. Only point separation is needed for the covariance comparison; coincident coordinates of distinct points are allowed.


On a fixed healthy chart, subtract the radial Taylor polynomial of each local insertion. The coefficient norm of the remainder is bounded by its next ordered derivative times the corresponding power of the coordinate radius. All derivatives needed here are among the finitely many stage orders. Express each polynomial in curvature and differentiated-curvature variables, carrying along their powers of~$r$. In the finite jet coordinates $\xi_1,\ldots,\xi_{n_\xi}$ write the covariance as $C=(C_{ab})$ and put $\Delta_C=\sum_{a,b}C_{ab}\,\partial^2/\partial\xi_a\partial\xi_b$. Gaussian contraction and Wick ordering are the mutually inverse operators
\begin{equation*}
 e^{\Delta_C/2}\pi=\sum_{j\ge0}\frac1{2^jj!}\Delta_C^j\pi,
 \qquad {:}\pi{:}_C=e^{-\Delta_C/2}\pi.
\end{equation*}
The series terminate on the chosen polynomial space, since each contraction removes two factors. We Wick-order each insertion with respect to its internal contractions; contractions between distinct insertions remain. The leading coefficient change $T_m$ of Lemma~\ref{lem:E-marked-step} and Wick ordering, in the jet normalization of Definition~\ref{def:D07}, are invertible and preserve the unit and real form. The error in replacing the selected coefficient map by this leading map has norm less than $1/4$ after its inverse, by the marked estimate~\eqref{eq:E-atom-estimates} and the entry condition. Solving this finite linear system gives one lift for each free basis polynomial, common to every row. Transport only between the fixed levels $\reflevel{d}$ and~$t$ bounds these reference coefficients independently of~$N,p$ for that interval. All matrices are real on the real form, and the unit equation has the exact solution~$\mathbf u_d$.

The covariance calculation involves finitely many derivatives of $\Gamma(x)=(4\pi^2|x|^2)^{-1}$. For each pair of coloured plaquette stencils the rescaled Fourier chart $K(\xi)$ is supported in $|\xi|\le2$. Its leading symbol is the Maxwell curvature symbol $H(\xi)$. For every fixed required order~$q_{\rm der}$, the quadratic-form calculation of Lemma~\ref{lem:B2}, differentiated in frequency, gives
\begin{equation*}
 |\partial^jK(\xi)|\le C|\xi|^{-j},\qquad
 |\partial^j(K-H)(\xi)|\le C|\xi|^{1-j}
 \quad(0<|\xi|\le1,\ j\le q_{\rm der}).
\end{equation*}
These bounds follow by differentiating the inverse quadratic symbol: its curl factors remove the inverse-square singularity; replacing each lattice difference by its linear term leaves one extra power of~$|\xi|$. The regulator and its derivatives obey the same fixed bounds; different colours have zero cross covariance. To convert the second estimate to position space, decompose a cutoff of $K-H$ into annuli $|\xi|\asymp2^{-v}$. A derivative of order~$s$ in position inserts $s$ powers of~$\xi$. After $q_{\rm der}>5+s$ integrations by parts on each annulus its transform at~$Rx$, with $x$ in a compact separated set, is bounded by
\begin{equation*}
 C2^{-v(5+s)}\min\{1,(R2^{-v})^{-q_{\rm der}}\}.
\end{equation*}
Splitting the sum at $2^v\asymp R$ gives $CR^{-5-s}$. The smooth part away from zero decreases faster than this. The leading homogeneous symbol is the corresponding curvature derivative of~$\Gamma$ away from the origin, with the Fourier normalization of Part~I. Including the curvature scaling powers therefore gives absolute error $O(R^{-1})$ in the rescaled coordinates, including where a leading tensor component vanishes.

For the finite jet samples, terms of the matching homogeneity give this leading derivative exactly. Every other term has at least one extra inverse power of~$R$; summing their finitely many absolute coefficients preserves the $O(R^{-1})$ bound. Here $R=(\cutoffsize_t)^b$, and rounding in the rescaled variables gives an additional $O(R^{-1})$. Localizing the sector covariance introduces $C(\cutoffsize_t)^{q_{\rm der}} e^{-cB_t}$ by the inverse-covariance bound in Lemma~\ref{lem:B-windows}, with smaller $c$, since $\widehat\beta_i\ge\widehat\beta_t/2$ throughout the interval. Conditional on the joint-law estimate above, its nonlinear healthy-chart error is $C(\cutoffsize_t)^{q_{\rm der}}\widehat\beta_t^{-\delta_{\rm G}}$ for some fixed $\delta_{\rm G}>0$. The derivative orders and powers~$q_{\rm der}$ are fixed by the degree and the finite row family; their stage constants are fixed before entry. Since $\log\cutoffsize_t=o(\log\widehat\beta_t)$ and $B_t\asymp(\log\widehat\beta_t)^2$, both errors are smaller than any prescribed fixed negative power of~$\cutoffsize_t$ after entry. Choose $b$ so that the inverse-separation error is smaller than the finite power required by the mark weight below. The chart and inspected supports lie inside the nonwrapping comparison interval in~\eqref{eq:E-window-geometry}, so these local estimates hold on every period, with the same constants.


The finite-polynomial estimate combining these errors is as follows. Let $U$ bound the Gaussian comparison polynomials, $T$ the one-factor error, and $Z$ the two-factor error in their coefficient norm. For the connected two-point term, subtract the product of the one-point terms. Expanding that product gives the bound
\begin{equation}\label{eq:E-Gaussian-polynomial}
 T+Z+(U+T)T+TU.
\end{equation}
For example, if $a_i$ and $a_i^0$ are the two one-point terms, then
$a_1a_2-a_1^0a_2^0=(a_1-a_1^0)a_2+a_1^0(a_2-a_2^0)$, whose norm is at most $(U+T)T+UT$. The terms $T,Z$ are the errors in the remaining one- and two-mark polynomial terms. Gaussian contraction expands into pairings, so the bound for any fixed polynomial is a finite polynomial in the bounds for its contractions. After $b$ has been fixed sufficiently large, the preceding covariance and Taylor estimates make~\eqref{eq:E-Gaussian-polynomial} smaller than every source bound required at~$l$.

For more slots, let $\mathcal G'_J$ be the Gaussian polynomial moment before replacing its covariance by the Maxwell covariance. On the compact separated position set, all these moments are bounded by a fixed~$U\ge1$, and their discrepancies are at most $\varepsilon$, a constant times the covariance and Taylor error above. For a partition with $q$ blocks, telescoping its product bounds the discrepancy by $q\varepsilon(U+\varepsilon)^{q-1}$. The partition formula for $\mathsf C^{\rm Max}$ therefore bounds the sum of all nonempty corrected coefficients by
\begin{equation*}
 C_n\varepsilon(1+U+\varepsilon)^{n-1}.
\end{equation*}
There are finitely many subsets and partitions for fixed~$n$. Their constants are included in~$\mathcal C_j$. Subtracting $\mathsf C^{\rm Max}$ cancels every leading free coefficient in the logarithm; for a pair this is precisely the mean-product subtraction in~\eqref{eq:E-Gaussian-polynomial}. A support carries a mode only when it contains all its anchors. The support-independent scalar correction is retained separately per cell. Thus no local term is assigned to a support lacking one of its marks.

For terms meeting a flagged region, keep that region's complete source support and its original eligibility test. At a step $i+1\to i$ in stage~$j$, let $v_\star$ be the integrated majorant of the distinguished marked factor and $n_\star$ its covering number, as in Lemma~\ref{lem:D-connected}. Its rooted forced bound is $v_\star e^{\eta^{\rm att}_in_\star}e^{-r_j\beta'_u t_{i+1}^2}$ before division by the forced source normalization. Conditional on the flagged-source estimate stated above, half this suppression absorbs the derivative and support factors of the finite marks; the other half bounds their normalized forced remainder by the error above after entry. The finite polynomial contractions, Taylor subtractions and connected attachments contribute at most $\mathcal C_j$ at each of the $h_t(b)$ steps. Their product is the factor in the statement. At each step the mode-weighted coefficient norm is submultiplicative; the derivative and support weights remain in the separate integrated majorants. The rooted sums are anchored at marks, so their bounds have no factor counting the replicated cells.
\end{proof}

The next estimate increases the mark weight while retaining the full regular and forced representation.
\begin{lemma}[Corrected Gaussian sources]\label{lem:E-Gaussian-sources}
Uniformly on that set, arbitrarily late intervals admit common real lifts of the free basis, bounded in $N,p$ for each fixed interval and fixing the unit, such that the corrected source $e^{-\mathsf C^{\rm Max}(z)}\prod_h(1+z_hF_h^{\rm tst})$, carried to level~$k$, has a representation~\eqref{eq:E-source-form} there with mark weight
\begin{equation*}
 \mathsf w_t=\exp(\log\cutoffsize_t/(8n)).
\end{equation*}
At this weight its regular and forced source sums are each at most one, with scalar source bound at most one per cell.
For each such interval there is $N_0$ such that the assertion holds for all $N\ge N_0$ and all~$p$.
\end{lemma}
Its proof iterates the middle-level error over the second part of the interval.
\begin{proof}
Starting from the error of Lemma~\ref{lem:E-source-polynomial} at~$l$, a quadratic recurrence from the connected expansion of Lemma~\ref{lem:D-connected} carries the bound to~$k$; the extra length $v_t$ of the window absorbs the per-step losses. Work in the submultiplicative mode-weighted coefficient norm defined in the proof of Lemma~\ref{lem:E-source-attachments}. Keep derivative and support weights in separate integrated majorants and retain the forced eligibility tests.

We now bound the loss over all steps; their number grows with~$t$. Use~\eqref{eq:E-source-initial}. Increasing the fixed power $q_{\rm der}$ includes the finite Taylor and rounding bounds above. The mark weight contributes at most $\mathsf w_t^n=\mathcal M^{1/8}$ to any mode in one cell. Enlarge $\mathcal C_j$ to include all per-step finite polynomial and support constants. From~\eqref{eq:E-log-cost} and the two ceilings in~\eqref{eq:E-window}, on sufficiently late intervals with fixed~$b$,
\begin{equation}\label{eq:E-source-loss}
 \mathcal C_j^{h_t(b)+v_t+2}
 \le\mathcal M^{(b+6)/j^2}B_t^{10/j^2}
 \le\mathcal M^{1/64}.
\end{equation}
Indeed $\log B_t=2\log\log\widehat\beta_t+o(1)$ and $\log\mathcal M=(\log\log\widehat\beta_t)^2$, so $B_t=\mathcal M^{o(1)}$. The two ceilings and the $+2$ contribute at most $L^{4/j^2}\le\mathcal M^{4/j^2}$, since $\ell^{\rm st}_{J_t}\le\cutoffsize_t$ by item~6 of Table~\ref{tab:constant-choices}; the final inequality holds at all sufficiently large stages for each fixed~$b$.

We use the following one-step source estimate. For the fixed lifts of Lemma~\ref{lem:E-source-polynomial}, let $S_i$ majorize the corrected singleton terms, and $P_i$ the corrected terms with at least two marks, in the weighted regular norm at mark weight~$\mathsf w_t$, including their scalar terms per cell. Use the same quantities for the forced norm after division by its positive source normalization, and take the larger values. For $k\le i<l$, the connected product expansion gives
\begin{align}\label{eq:E-source-step}
 S_i&\le K_j^{\rm src} S_{i+1}+\mathcal M^{1/2}S_{i+1}^2,\notag\\
 P_i&\le K_j^{\rm src} P_{i+1}
          +\mathcal M^{1/2}(S_{i+1}+P_{i+1})^2.
\end{align}
Here $K_j^{\rm src}\ge1$ is fixed at the stage, and $\mathcal C_j\ge K_j^{\rm src}+2$. To obtain these bounds, separate connected terms linear in the incoming corrected sources from those with at least two source factors. The former use the one-mark integration bound. In the latter, the product rule leaves a product of two source norms; summing the additional attachments uses the rooted clause of Lemma~\ref{lem:D-connected}. Its finite derivative and localization losses are bounded by $\mathcal M^{1/2}$ after entry. Every higher product is absorbed in the same quadratic bound while the source norm is at most $\mathcal M^{-1/2}$. For the forced terms the retained eligibility event supplies the forced suppression before its normalization is divided out. The entry requirement in item~\ref{it:E-entries} includes this normalized forced bound. Disjoint cell modes use the multiplicative weights of Lemma~\ref{lem:E-replication}; their connecting supports keep the same estimate. No factor counts the number of cells.

Set $\mathfrak a_i=S_i+P_i$. Equations~\eqref{eq:E-source-step} imply $\mathfrak a_i\le K_j^{\rm src}\mathfrak a_{i+1}+2\mathcal M^{1/2}\mathfrak a_{i+1}^2$. The first part of the interval gives $\mathfrak a_l\le\mathsf w_t^n\mathcal C_j^{h_t(b)+2}\mathcal E_t^{\rm G}$. By~\eqref{eq:E-source-loss}, for $b\ge1$ and sufficiently late entry,
\begin{equation}\label{eq:E-source-iteration}
 \mathsf w_t^n\mathcal C_j^{h_t(b)+v_t+2}\mathcal E_t^{\rm G}
 \le\mathcal M^{-3/4},\qquad
 \mathfrak a_i\le\mathcal C_j^{l-i}\mathfrak a_l\le\mathcal M^{-3/4}<\mathcal M^{-1/2}.
\end{equation}
The second inequality follows by backward induction: below $\mathcal M^{-1/2}$ the quadratic term in the recurrence is at most $2\mathfrak a_{i+1}$. For the first, the covariance contribution is at most $\mathcal M^{1/8+1/64-b}$, whose exponent is at most $-55/64$. The other contributions tend to zero faster than every fixed negative power of~$\mathcal M$, since $\log\mathcal M=o(\log\widehat\beta_t)$ and $B_t\asymp(\log\widehat\beta_t)^2$. Thus the initial-error condition is compatible with the whole interval, although its length grows with~$t$.

The extra length $v_t$ also provides the support margin: a level-$l$ neighborhood of radius $B_t^2$ has level-$k$ radius at most $B_t^2/L^{v_t}\le B_t^{-8}$. It therefore fits inside the protected cell around its rounded reading points. The successive kernel neighborhoods add a bounded geometric sum in those units; their fixed margin was included in the choice of the comparison cell. Thus the inspection and attachment tests used in the iteration stay inside the stated interval. The unit has no remainder by Lemma~\ref{lem:E-marked-step}.

The support margin in~\eqref{eq:E-window-geometry} makes every replicated inspected stencil fit its own cell. Terms meeting flagged regions keep their original forced eligibility and the suppression used in Lemma~\ref{lem:D-connected}; their source amplitude replaces only the amplitude of the distinguished factor. Thus the local source sums at the enlarged mark weight are at most one in both regular and forced norms. The constant source contribution is at most one per cell. These estimates prove the lemma on arbitrarily late intervals, with all thresholds independent of~$p$.
\end{proof}

The remaining normalization uses the complete positive unmarked density. In particular it cannot discard the forced sum or replace it by its healthy part.
\begin{samepage}
\begin{lemma}[Normalized source integral]\label{lem:E-source-integral}
For the corrected sources of Lemma~\ref{lem:E-Gaussian-sources}, put $u_t=\log\cutoffsize_t/(16n)$. For each such interval there is $N_0$ such that, for all $N\ge N_0$ and all~$p$, for every nonempty mode~$I$,
\begin{equation*}
 \left|\left\langle\prod_{C\in\mathcal C^{\rm tile}_k}(\mathfrak c_I)_C\right\rangle_{N,p}\right|
 \le e^{|\mathcal C^{\rm tile}_k|q_t^{\rm cmp}}
       \mathsf w_t^{-|I||\mathcal C^{\rm tile}_k|},
 \qquad q_t^{\rm cmp}\le u_t.
\end{equation*}
\end{lemma}
\end{samepage}
\begin{proof}
The bound compares the source integral with the complete positive unmarked density of Definition~\ref{def:D06}, whose positivity follows from the density identity of Theorem~\ref{prop:RG}, the positive Wilson start and the probability kernels of Lemma~\ref{lem:B1}; the source sizes come from Lemma~\ref{lem:E-Gaussian-sources}. Let $\rho$ be the unmarked representation at the width level~$k$, $L=n_{p,k}$ its torus side, $V=|\Lambda_L|=|\mathcal C^{\rm tile}_k|$, $g=\beta(\rho)^{-1/2}$, $H_{\rm path}=A_{J_k}^{\rm path}$, and let $\delta,\Lambda^{\rm loop}$ bound its regular and loop terms. The admissibility inequalities give $0<g\le t_k$, $\delta\le1$, $\Lambda^{\rm loop}\le\Lambda_k^{\rm loop}$ and $\beta(\rho)\le4\widehat\beta_k$.

Choose constants $s_H\in(0,1]$, $b_H\ge1$ such that the chord ball of radius $r/4$ has Haar measure at least $s_Hr^{\dimG}$ for small $r$, and the clipped one-link integrals are at most $b_H(g/\sigma_*)^{\dimG}$, where $\sigma_*$ is the smaller of one and the linear clip coefficient. These follow from the nonsingular exponential chart and compactness of the group, with constants chosen before all levels. Write
\begin{equation*}
 \mathfrak h=\log b_H-\dimG\log\sigma_*,\qquad
 \mathcal P(\rho,X,g)=2E_{\rm ex}+6+
 \beta(\rho)E_{\rm base}+X+3\mathfrak h-4\log s_H.
\end{equation*}
Here $E_{\rm ex}$ bounds the loop and regular exponent per site, including the exponent discrepancy $\delta$ in the clipped comparison, and on links of size at most $g/4$,
$\beta(\rho)E_{\rm base}=3(c_{{\rm dir},J_k}+H_{\rm path}^2)$. For any integrable complex~$F$ with
$z|F(U)|\le e^{\Phi_\rho(U)+VX}$, $z\ge0$, the ratio estimate is
\begin{equation}\label{eq:E-positive-ratio}
 \frac{z|\int F\,dU|}{\int\density\rho\,dU}
 \le e^{V\mathcal P(\rho,X,g)-4\dimG\log g}.
\end{equation}
The gauge fixing cancels the small-coupling powers in this ratio. The denominator has $3V+1$ free links and the numerator will integrate $3V-3$ selected plaquette factors, so their Haar-ball powers differ only by $g^{\dimG(3V-3)}/g^{\dimG(3V+1)}=g^{-4\dimG}$, a cost that is not extensive in~$V$; all remaining constants are exponential in~$V$. For the denominator, fix an axial spanning tree with $V-1$ links. Gauge invariance removes those links, leaving $3V+1$ free links. Restrict each free link to the ball of chord $g/4$. All plaquettes are then healthy, so the forced covering sum equals one. The lower action bound and the Haar-ball bound give
\begin{equation*}
 \int\density\rho\,dU\ge
 e^{v_{{\rm sc}}V-V(E_{\rm ex}+\beta(\rho)E_{\rm base})}
 (s_Hg^{\dimG})^{3V+1}.
\end{equation*}
For the numerator, the clipped plaquette inequality is
\begin{equation*}
 \beta(\rho)E_{\rm plaq}(U)\ge
 \frac{\chord(U_\square)}{g/\sigma_*}-1.
\end{equation*}
It follows in the quadratic branch from $(a/g-1)^2\ge0$ and $\sigma_*\le\min(1,c^{\rm dir})$, and in the linear branch from $\sigma_*\le c^{\rm lin}$ and $g$ below the threshold. Summing over the six plaquettes per site retains the product of $e^{-\chord(U_\square)/(g/\sigma_*)}$ in the upper integrand.

% The explicit axial elimination argument adapts Humanise YM, rank-root-normalization.tex; see the Area D migration handoff for CC BY 4.0 provenance.
Keep the $3V-3$ plaquette factors selected below and discard the others, whose values lie in $[0,1]$. Represent coordinates by $0,\ldots,L-1$, and let $p(x)$ be the largest index of a nonzero coordinate, with $p(0)=0$. The axial tree joins each $x\ne0$ to $x-e_{p(x)}$. Leave the four free wrap links $(-e_i,i)$ untreated. For every other free link $e=(x,i)$ put $t=\max\{i,p(x)\}$ and order links by decreasing lexicographic priority
\[
 \bigl(t,\mathbf 1_{\{i=t\}},\textstyle\sum_a x_a\bigr).
\]
If $i<t$, choose the $(i,t)$ plaquette based at $x-e_t$. Its two $t$-links are in the tree, and its other $i$-link has lower priority if free. If $i=t$, freeness implies $x_i=L-1$ and all higher coordinates vanish. Since $e$ is not one of the four retained links, $w=x+e_i\ne0$. With $j=p(w)<i$, choose the $(j,i)$ plaquette based at $x-e_j$. One $j$-link is in the tree; the other free links, if any, have lower priority. In both cases the chosen plaquette contains $e$ once. Two links cannot choose the same plaquette, since each would have lower priority than the other. Consequently the highest-priority remaining link occurs in no other remaining selected factor, and Haar invariance integrates its own factor to the one-link integral. Repeat in descending order. The four untreated links account for $(3V+1)-4=3V-3$; the argument includes the periodic boundary. Thus
\begin{equation*}
 z|\int F\,dU|\le
 e^{v_{{\rm sc}}V+V(E_{\rm ex}+6+X)}
 \bigl(b_H(g/\sigma_*)^{\dimG}\bigr)^{3V-3}.
\end{equation*}
Divide these inequalities. The powers of $g$ differ by $-4\dimG$. Using $3V-3\le3V$, $3V+1\le4V$, $b_H\ge1$ and $s_H\le1$ gives~\eqref{eq:E-positive-ratio}.

One may use the clipped comparison exponent in this calculation even if it is not the exponent used for the source bound. Their difference has absolute value at most $V\delta$. Multiplying the upper source bound by $e^{V\delta}$ and bounding the two positive partition integrals in the opposite direction by the same factor accounts for this change. These two factors are included in $2E_{\rm ex}$ because each occurrence of $E_{\rm ex}$ includes $\delta$. Thus~\eqref{eq:E-positive-ratio} applies to the actual source integral on the original exact history.

The enlarged-weight bounds of Lemma~\ref{lem:E-Gaussian-sources} give a pointwise bound with $X\le2+2n$. Use the weighted $\ell^1$ norm of the squarefree coefficients. The regular exponential is bounded by the exponential of the sum of its scalar and regular source norms. Write $K_Y^{\rm src}$ for the sum of the nonempty marked coefficients at forced address~$Y$, keeping its index, support and inventory even when supports coincide. For the complete forced sum, dropping eligibility and disjointness only in the upper bound gives
\[
 \sum_{\mathcal A}\prod_{Y\in\mathcal A}
 \bigl(|K_Y|+\|K_Y^{\rm src}\|_{\mathsf w_t}\bigr)
 \le\prod_Y\bigl(1+|K_Y|+\|K_Y^{\rm src}\|_{\mathsf w_t}\bigr)
 \le e^{\sum_Y|K_Y|+\sum_Y\|K_Y^{\rm src}\|_{\mathsf w_t}}.
\]
The unmarked forced sum is at most $V\epsilon_k^{\rm for}\le V$ by Definition~\ref{def:D05b}; the scalar source is at most one per cell, and the regular and forced source sums are each at most one per mark. Thus the total exponent is at most $V(2+2n)$. No individual forced activity need be positive. Before integration the coefficient of the disjoint union of the chosen modes across all cells includes the factor $z=\mathsf w_t^{|I|V}$. The entry inequalities imply
\begin{equation*}
 E_{\rm ex}\le8H_{\rm path}+2,\qquad
 -4\dimG\log g/V\le1.
\end{equation*}
For the second, $-4\dimG\log g=2\dimG\log\beta(\rho)$ is $O(\log\widehat\beta_k)$, whereas $V\ge D_k=e^{s_k}$. The first follows from
$\widehat\beta_kt_k^3\le1$, $\widehat\beta_kt_k^4\le1$ and $\Lambda_k^{\rm loop}t_k^2\le1$ on sufficiently late levels. Hence choose a stage constant
\begin{equation}\label{eq:E-comparison-cost}
 Q_j^{\rm cmp}\ge2(8A_j^{\rm path}+2)+6+3(c_{{\rm dir},j}+(A_j^{\rm path})^2)+(2+2n)
 +3\mathfrak h_j-4\log s_H+1
\end{equation}
for each arity used at stage~$j$. It bounds the exponent per cell in~\eqref{eq:E-positive-ratio}. All its quantities are fixed before the entry to stage~$j$. Item~\ref{it:E-entries} in Lemma~\ref{lem:E1.1} requires $Q_j^{\rm cmp}\le\log\cutoffsize_t/(16n)$ on the chosen intervals. This proves $q_t^{\rm cmp}\le u_t$ and the asserted coefficient estimate. In particular the volume appears only through the number of cells, so the cutoff threshold is independent of~$p$.
\end{proof}


The next estimate uses the corrected source comparison in Lemma~\ref{lem:E-Gaussian-sources}. Each relative error is interior whenever its insertions are interior, since~\eqref{eq:F-relative} is a polynomial in those insertions.
\begin{lemma}[Relative rank comparison]\label{lem:E2.10}
For every such finite test family, there are $r_\nu\downarrow0$, $y^{\nu,i}\to y^i$, levels $k_\nu\le t_\nu$ tending to infinity, $\gamma_\nu\ge0$ tending to zero, and real lifts~$b^\nu_{j,N,p}$ of the free basis~$(\mathfrak m_j)$, bounded in~$N,p$ for fixed~$\nu$ and common to all tests. They satisfy $b^\nu_{1,N,p}=\mathbf u_d$. Each~$x^{\nu,i}$ is generic on the spacing-$a_{t_\nu}$ grid. Eventually uniformly in~$p$, all test insertions are interior to~$\mathcal C^{\rm tile}_{k_\nu}(0)$, and, for every nonempty~$I$,
\begin{equation*}
\|R_I\|_{\mathcal C^{\rm tile}_{k_\nu}(0)}\le\gamma_\nu\,.
\end{equation*}
\end{lemma}

\begin{proof}[Proof of Lemma~\ref{lem:E2.10}]
The relative errors are the corrected coefficients~\eqref{eq:E-corrected-coefficients}, and Lemmas~\ref{lem:E-Gaussian-sources}--\ref{lem:E-source-integral}, with the free correlations of Definition~\ref{def:D11}, bound their seminorms by a negative power of the cutoff size. Apply those lemmas with arity bound $n=2$, so $\mathsf w_t=(\cutoffsize_t)^{1/16}=e^{2u_t}$ and $u_t=\log\cutoffsize_t/32$. The normalized source lemma gives
\begin{equation}\label{eq:E-rank-input}
 \left|\left\langle\prod_{C\in\mathcal C^{\rm tile}}(\mathfrak c_I)_C\right\rangle_{N,p}\right|
 \le\exp\{|\mathcal C^{\rm tile}|q_t^{\rm cmp}\}
 \mathsf w_t^{-|I||\mathcal C^{\rm tile}|},\qquad q_t^{\rm cmp}\le u_t.
\end{equation}
The cutoff threshold is independent of~$p$.

Taking the $|\mathcal C^{\rm tile}|$-th root in~\eqref{eq:E-rank-input} gives
\begin{equation*}
\|\mathfrak c_I\|_{\mathcal C^{\rm tile}}\le e^{u_t-2|I|u_t}\le e^{-u_t}
=(\cutoffsize_t)^{-1/32}\longrightarrow0.
\end{equation*}
Since $e^{\mathsf C^{\rm Max}}$ has coefficient $\mathcal G_J$ at $z_J$, multiplication by $e^{\mathsf C^{\rm Max}}$ shows that $\mathfrak c_I$ is exactly $R_I$ in~\eqref{eq:F-relative}. Choose successive comparison windows with both levels increasing, reading radii strictly decreasing to zero, and $e^{-u_t}$ decreasing to zero. The windows of Lemma~\ref{lem:E1.1} allow these finitely many requests at each step. Call the resulting levels $k_\nu\le t_\nu$, radii $r_\nu$ and errors $\gamma_\nu=e^{-u_{t_\nu}}$. Rounding on the finer grid preserves genericity for each fixed~$\nu$ once that grid is fine enough. Interior insertion supports imply interiority of every $R_I$, by its polynomial formula.
\end{proof}

\subsubsection{Bounds for lattice differences}\label{sec:partII-E.2.6}

Positional differences are a separate problem. A difference quotient multiplies the insertion by the inverse spacing,
\begin{equation}\label{eq:E-inverse-spacing}
\frac{O^{(d)}_{N,p}(c,y+e_a)-O^{(d)}_{N,p}(c,y)}{a_N}
=a_N^{-(d+1)}\Bigl(P_d[t_N^{-\factorcount}\mathsf Z_{N,p,d}^{-1}c](y+e_a)
-P_d[t_N^{-\factorcount}\mathsf Z_{N,p,d}^{-1}c](y)\Bigr),
\end{equation}
so Lemma~\ref{lem:E2.3}, applied to the two terms on the left, bounds its seminorm only by $2Aa_m^{-d-\varepsilon}/a_N$, which diverges as $N\to\infty$. The remedy has three steps. Local matching represents the difference quotient as a degree-$(d+1)$ insertion up to a negligible error (Lemma~\ref{lem:E2.9}); the factor $a_N^{-(d+1)}$ becomes the normalization of the new insertion, but its coefficient is only tame, not bounded. The equation projection removes the directions of that coefficient whose insertions vanish in separated correlations (Lemmas~\ref{lem:E2.6}--\ref{lem:E2.7}). The remaining coefficient lies in a space of dimension at most $\dim_\C\maxwellspace{d+1}$ and is bounded by testing it against the finite free matrix (Lemmas~\ref{lem:E2-free-tests} and~\ref{lem:E2.10}).

The tests of the difference must be bounded without any bound on differentiated insertions. Averaging the difference quotient over $Q_N\approx h_{\rm av}/a_N$ consecutive translates telescopes to a difference over the fixed physical distance $Q_Na_N$, which is bounded by undifferentiated insertions alone; this is why the argument is not circular. Comparing averaged tests at nearby positions uses the following displacement estimate, which also involves only undifferentiated insertions. For a grid configuration~$x$ define its minimum coordinate separation on the physical torus by
\begin{equation*}
 g_p(x)=\min_{i\ne j}\min_{a=1,\ldots,4}\min_{z\in L_p\Z}
 |x_i^a-x_j^a+z|,
\end{equation*}
with value $+\infty$ for one slot. This separation supplies the empty slabs for reflection positivity. A coordinate hybrid of $x,y$ chooses each coordinate from either configuration. Assume the insertions and their single-axis reflections are interior at a centered level~$D\le N$, with all their seminorms at most~$A\ge0$.
\begin{lemma}[Displacement of correlation tests]\label{lem:E-displacement}
For every $r>0$ with $a_D<r$, a one-coordinate grid move $x\mapsto y$ of size at most~$r$ from $g_p(x)\ge4r$ satisfies
\begin{equation*}
 |M_{N,p}(v;x/a_N)-M_{N,p}(v;y/a_N)|
 \le A^n|x-y|/r.
\end{equation*}
If $|x-y|\le r$ and every coordinate hybrid has coordinate separation at least~$4r$, then
\begin{equation*}
 |M_{N,p}(v;x/a_N)-M_{N,p}(v;y/a_N)|
 \le4nA^n|x-y|/r.
\end{equation*}
\end{lemma}
\begin{proof}
A move changes the lengths of the two empty slabs next to the moved insertion, and the spectral theorem for the transfer operator bounds the effect of that change; the reflection inequality~\eqref{eq:E-chess} bounds what remains. Lemma~\ref{lem:E2.3} supplies the assumed seminorms for bounded coefficient lists on full cutoff tails. The separation hypothesis gives every directed coordinate distance at least~$4r$, at any period. No fixed test neighbourhood is used.

Fix an axis and all other coordinates. Reflection positivity and translation invariance give a positive transfer operator, rescaled to $0\le T\le1$. On the periodic axis the expectation is a normalized trace. If $X$ denotes the insertion being moved, its two adjacent empty slabs have lengths $n_1^{\rm sl},n_2^{\rm sl}$ in fine sites. A move by $l\ge0$ changes their factor by
\begin{align}\label{eq:E-trace-displacement}
 T^{n_1^{\rm sl}+l}XT^{n_2^{\rm sl}-l}-T^{n_1^{\rm sl}}XT^{n_2^{\rm sl}}
 &=T^{n_1^{\rm sl}}(T^lX-XT^l)T^{n_2^{\rm sl}-l}\notag\\
 &=-T^{n_1^{\rm sl}}(1-T^l)XT^{n_2^{\rm sl}-l}
   +T^{n_1^{\rm sl}}X(1-T^l)T^{n_2^{\rm sl}-l}.
\end{align}
Both expressions are inserted in the same remaining cyclic product and divided by the same partition trace. The stencil is contained in a centered cell of lattice side $d_{\rm cell}=a_D/a_N$. Since $a_D<r$, the directed separations exceed $2d_{\rm cell}+2r/a_N$, and $la_N\le r$ by assumption. Splitting the two surrounding empty slabs at their reflection planes leaves $n_{\rm sl}+1\ge2r/a_N$ fine sites next to each occurrence of $1-T^l$. The spectral theorem gives
\begin{equation*}
 \|(1-T^l)T^{n_{\rm sl}}\|
 \le\sup_{0\le t\le1}l(1-t)t^{n_{\rm sl}}
 \le l/(n_{\rm sl}+1)\le la_N/(2r).
\end{equation*}
For each normalized trace in~\eqref{eq:E-trace-displacement}, reflection Schwarz at the slab ends and the chessboard inequality bound the remaining factors by the product of their centered seminorms, at most~$A^n$. One applies this inequality to the complete cyclic product, so no supremum norm of~$X$ is required. The two terms together are bounded by $A^nla_N/r$. An arity-one moment is translation invariant and has change zero.

For a negative move reverse the chosen axis. Coordinate separation is unchanged and the assumed reflected seminorms have the same bound, so the same slab estimate applies. Replace the $4n$ coordinates one at a time. Each starting configuration is a coordinate hybrid and hence has the required separation. Each move is at most $|x-y|\le r$. Summing proves the second bound. The constants depend only on the seminorms, arity and slab margin, uniformly in~$N,p$.
\end{proof}

All the preceding comparison lemmas hold on full cutoff tails, uniformly in~$p$. Local matching also applies to any fixed iterated difference: after multiplying by $a_N^{d+h}$, a difference of order~$h$ is a finite sum of translates of the degree-$d$ section, with subpower derivative bounds. Matching in degree $d+h$ leaves a seminorm error tending to zero at every fixed cell size, uniformly in~$p$. This finite-cutoff argument supplies the difference bounds required in Part~I.
\begin{lemma}[Difference estimate]\label{lem:E2.4}
For every arity $n$, insertion-label list $v$, difference list $\boldsymbol l$ and $r>0$, the constants in Proposition~\ref{prop:UB} can be chosen independently of~$p$:
\begin{equation*}
\sup_p\sup_{\operatorname{sep}_{\rm tor}(y)\ge r}
|\Delta_{\boldsymbol l}M_{N,p}(v;y)|\le B(n,v,\boldsymbol l,r)
\qquad(N\ge N_0(n,v,\boldsymbol l,r)).
\end{equation*}
\end{lemma}
\begin{proof}
We first bound the chessboard seminorm of one differentiated insertion at a fixed fine cell size, uniformly in $N,p$ (Steps~1--4), and then apply~\eqref{eq:E-chess} (Step~5). Let $F=O^{(d)}(c)$, let $\Delta_{\boldsymbol a}F$ be a fixed iterated normalized difference in its anchor, and put $d'=d+|\boldsymbol a|$.

\emph{Step 1: Matching and the equation part.} Local matching gives tame coefficients $c^{\boldsymbol a}$ in degree $d'$ with $O^{(d')}(c^{\boldsymbol a})-\Delta_{\boldsymbol a}F$ of vanishing seminorm. Write $c^{\rm pr}=\Pi_{d',N,p}c^{\boldsymbol a}$. For a tame coefficient in $\ker\Pi_{d'}$, the auxiliary equation section has seminorm zero. To see this, take all the other reflected copies as the spectator in the exact Haar identity in Lemma~\ref{lem:A-R4}. For a fixed physical cell, its side in lattice units eventually exceeds twice the equation radius plus two; all those copies are then exterior. The section-comparison residual of Lemma~\ref{lem:E2.7} tends to zero in seminorm. In particular it applies to $z=(c^{\boldsymbol a}-c^{\rm pr})/(1+\|c^{\rm pr}\|)$. Indeed, $c^{\boldsymbol a}/(1+\|c^{\rm pr}\|)$ is tame, and $c^{\rm pr}/(1+\|c^{\rm pr}\|)$ has norm at most one. Every bounded coefficient family is tame, because Lemma~\ref{lem:E2.1} gives
\begin{equation*}
 \|\mathsf Z_{N,p,d'}^{-1}\|\le C_{d'}\prod_{i=\reflevel{d'}}^{N-1}C^{\rm col}_{J_{i+1}}=a_N^{-o(1)}
\end{equation*}
by~\eqref{eq:E-log-cost}. Subtracting these two tame families proves tameness of~$z$ before~$c^{\rm pr}$ is bounded.

\emph{Step 2: The averaged tests are bounded.} It remains to bound~$c^{\rm pr}$. Use a fixed sufficiently short-distance element of the comparison sequence of Lemma~\ref{lem:E2.10} for degree~$d'$. Let $\delta_0>0$ be the minimum coordinate separation of its configuration on the physical torus. Fix a slab margin $0<b<\delta_0/8$ and a centered level~$D$ with $a_D<b$. Average the variable slot of every test over $Q_N=\lceil h_{\rm av}/a_N\rceil$ sites in each direction occurring in~$\boldsymbol a$, where $0<h_{\rm av}<b/[2(|\boldsymbol a|+1)]$. Eventually $Q_Na_N\le2h_{\rm av}$, so the total shift in each coordinate is less than~$b$ and every coordinate hybrid used in the averaging has coordinate separation at least~$4b$. Denote this linear map on coefficients by $\mathcal P_{N,p}$. Telescoping in one direction gives
\begin{equation*}
\frac1{Q_N}\sum_{j=0}^{Q_N-1}
\frac{f(z+(j+1)e_a)-f(z+je_a)}{a_N}
=\frac{f(z+Q_Ne_a)-f(z)}{Q_Na_N}.
\end{equation*}
Since $Q_Na_N\ge h_{\rm av}$, iterating bounds the tests of~$\Delta_{\boldsymbol a}F$ by $(2/h_{\rm av})^{|\boldsymbol a|}$ times the supremum of tests of~$F$. The latter is bounded by Lemma~\ref{lem:E2.3}, the interior-seminorm comparison and~\eqref{eq:E-chess}. Matching errors tend to zero in these tests: put the error and the bounded spectator insertions in distinct cells with fixed margins, use Lemma~\ref{lem:F1a} for their interior seminorms, and apply~\eqref{eq:E-chess}. All these bounds hold on full cutoff tails and use no positional difference estimate. Thus $\|\mathcal P_{N,p}c^{\boldsymbol a}\|\le C_{\rm test}$ eventually, uniformly in~$p$.

\emph{Step 3: Invertibility of the tests.} The finite-dimensional lower bound is obtained before any cutoff limit. The free matrix has a positive lower bound on the unit sphere. The relative-error estimate and the equation-seminorm bound put the tests of the projected lifts within $O(\gamma_\nu)+o_{N\to\infty}(1)$ of the free matrix, uniformly in~$p$. Choose the fixed~$\nu$ with this error smaller than its matrix lower bound. Averaging over sufficiently small~$h_{\rm av}$ preserves the lower bound: the undifferentiated correlation tests have uniform positional continuity by Lemma~\ref{lem:E-displacement}, with margin~$b$, the centered level~$D$ and the coordinate-hybrid separations just chosen. The level and its seminorm bounds are fixed before reducing~$h_{\rm av}$, so the displacement bound tends to zero with~$h_{\rm av}$. The projected lifts are therefore linearly independent. Their number equals $\dim_\C\maxwellspace{d'}$, whereas Lemma~\ref{lem:E2.6} bounds $\dim\operatorname{ran}\Pi_{d'}$ by that number. They consequently span the range. Their uniform coefficient bounds give
\begin{equation}\label{eq:E-detection}
\|c^{\rm pr}\|\le C_{\rm det}\|\mathcal P_{N,p}c^{\rm pr}\|
\end{equation}
with $C_{\rm det}$ independent of~$N,p$ on the retained tails. Choose an invertible square minor $H_0$ of the free matrix and the corresponding interacting minor~$H$. For $\|H-H_0\|\le(2\|H_0^{-1}\|)^{-1}$, the inverse bound $\|H^{-1}\|\le2\|H_0^{-1}\|$ follows.

\emph{Step 4: The bound on $c^{\rm pr}$.} The equation coefficient $z=(c^{\boldsymbol a}-c^{\rm pr})/(1+\|c^{\rm pr}\|)$ is tame. Its tests tend to zero, so eventually $C_{\rm det}\|\mathcal P_{N,p}z\|\le1/2$. Equation~\eqref{eq:E-detection} yields
\begin{equation*}
\|c^{\rm pr}\|\le C_{\rm det}C_{\rm test}+\tfrac12(1+\|c^{\rm pr}\|),\qquad
\|c^{\rm pr}\|\le2C_{\rm det}C_{\rm test}+1.
\end{equation*}
Lemma~\ref{lem:E2.3} now bounds the seminorms of~$O^{(d')}(c^{\rm pr})$ at each fixed sufficiently fine scale. Now $c^{\boldsymbol a}-c^{\rm pr}$ is tame, so its equation residual has vanishing seminorm. Together with the matching residual, the triangle inequality for seminorms gives the same type of bound for~$\Delta_{\boldsymbol a}F$. All thresholds and constants in this argument are independent of~$p$.

\emph{Step 5: Products at separated points.} Distribute the differences in $\boldsymbol l$ among their indicated slots. A difference in slot~$i$ acts only on that insertion, so the resulting moment is exactly the expectation of the product of these differentiated insertions. Choose a fixed tiling side smaller than~$r$ and translate its boundaries away from all anchors. In each axis the excluded boundary strips occupy less than one full cell length if their widths are chosen below $a_m/[4(n+1)]$. For sufficiently large~$N$ a lattice translation with these margins exists. Distinct anchors occupy distinct cells, since their torus separation exceeds the cell side. The interior-seminorm comparison, with a finer centered scale, bounds all the factors uniformly. Apply~\eqref{eq:E-chess} and put one in every unused cell. Their product gives the asserted~$B$. Empty difference lists and arities at most one obey the same argument. This proves the lemma and Proposition~\ref{prop:UB}.
\end{proof}

\subsubsection{Bounded approximations}\label{sec:partII-E.2.7}

The later phase argument needs bounded approximations with a bound independent of the cutoff. Compactness at each cutoff gives no such bound. We obtain it by stopping at one fixed coarse lattice and pulling bounded observables back through the probability block kernels.

In the next two lemmas, a statement holds eventually in the ordered limits when it holds for all sufficiently large indices of the chosen volume sequence and, at each such volume, all sufficiently fine cutoffs of the chosen spacing sequence. No common cutoff threshold in other subsequently fixed parameters is asserted. The chessboard seminorms used below are defined in Definition~\ref{def:D27}.
\begin{lemma}[Bounded coarse approximations]\label{lem:E-coarse-bounded}
For every $d$, $k\ge\reflevel{d}$, real $c\in V_d$ and $\delta>0$, there are measurable real observables $\mathcal B_{N,p}(c,y)$ and $M_{\rm bnd}<\infty$, independent of $N,p,y$, with $|\mathcal B_{N,p}(c,y)|\le M_{\rm bnd}$ and
\begin{equation*}
 \|\mathcal B_{N,p}(c,y)-O^{(d)}_{N,p}(c,y)\|_{\mathcal C^{\rm tile}_k(y)}\le\delta
\end{equation*}
eventually in the ordered limits. Both observables are interior to that cell, and the same holds for their signed and reflected images.
\end{lemma}
\begin{proof}
We stop at a fixed coarse level~$h$, replace the insertion by a bounded coarse polynomial pulled back through the probability block kernels (Definition~\ref{def:D04}), and bound the discrepancy by replicating it in every cell (Lemma~\ref{lem:E-replication}); the marked step of Lemma~\ref{lem:E-marked-step} matches the coefficients. The construction is independent of the phase decomposition. Choose a fixed $m\ge k$ with $B_m^2\le a_k/a_m$, using Lemma~\ref{lem:E1.1}. At a finer level $h\ge m$, the bounded coarse polynomial with coordinate~$v$ is
\begin{equation*}
 \mathsf V_h(v;U)=a_h^{-d}\left(v_0+
 \mathbf1_{\Sigma_d\,\mathrm{healthy}}P_d[t_h^{-\factorcount}v^\circ](U)\right).
\end{equation*}
Its supremum for each basis input is bounded by a constant depending on~$d,h$, uniformly in~$N,p$: the healthy polynomial uses a fixed compact coordinate range, and the unit term is constant. Its inspected stencil has a fixed radius, including the links needed for the relaxed threshold test.

Let $\mathcal K_{N,h}(U,dV)$ be the probability kernel obtained by composing the block kernels from $N$ to~$h$, and set
\begin{equation*}
 (\mathcal P_{N,h}F)(U)=\int F(V)\,\mathcal K_{N,h}(U,dV).
\end{equation*}
It preserves real values and does not increase the supremum norm. The locality of the block kernels confines its dependence to the fine preimage of the inspected stencil. For reflected cells these ancestor cones are disjoint once $h$ has the stated margin. At each intermediate level the block kernels are products in the coarse-link variables given the preceding field; integrating disjoint ancestor cones therefore gives conditional independence, and hence $\prod_C\mathcal P_{N,h}F_C=\mathcal P_{N,h}(\prod_CF_C)$. Choose $h$ sufficiently large that this preimage lies strictly inside the centered level-$k$ cell, including its boundary-link margin.

To compare with the actual insertion, choose auxiliary fine coefficients $u_j$ whose selected level-$h$ polynomial coordinates differ by at most $1/4$ from the basis vectors. They are obtained by carrying the pure coarse polynomial back through the inverse graph construction in Lemma~\ref{lem:E-marked-step}. Apply~\eqref{eq:E-match-cone} to the difference with this primary coarse polynomial. Write $S_h=C^{\rm col}_{J_h}$, and let $\mathcal J_h$ be the matrix of the selected level-$h$ coordinates of these trials. For a basis input, the cone and the finest-level boundary estimate give
\begin{equation}\label{eq:E-coarse-inversion}
 32(1+S_h)^2|e_j-\mathcal J_he_j|_1\le R_h,
 \qquad R_h\le2(1+S_h),\qquad
 |e_j-\mathcal J_he_j|_1\le\frac1{16(1+S_h)}\le\frac14.
\end{equation}
Taking the maximum over columns gives $\|\mathcal J_h-1\|_1\le1/4$, which implies $\|\mathcal J_h^{-1}\|_1\le2$. If $c$ is the required reference coefficient, linear recombination therefore gives
\begin{equation*}
 O(c)=\sum_j t_j O(\mathsf Z_{N,p,d}u_j),\qquad
 \sum_j|t_j|\le2e^{\sum_{i<h}\log C^{\rm col}_{J_{i+1}}}|c|_1.
\end{equation*}
Replace each term on the right by $\mathcal P_{N,h}\mathsf V_h(e_j)$. This defines a bounded complex observable $\mathcal B'_{N,p}$; its bound depends on $h,d,c$, and hence may depend on $\delta$, but is independent of~$N,p$.

To control the discrepancy, take one reflected copy in each cell of the tiling. For a fine field~$U$, let $O_C(U)$ be the copy of $O(\mathsf Z u_j)$ in cell~$C$; for a coarse field~$V$, let $\mathsf V_{h,C}(V)$ be the corresponding copy of the coarse polynomial. Given~$U$, the ancestor cones are disjoint, so the conditional factorization above gives the exact identity
\begin{equation*}
 \prod_C\bigl((\mathcal P_{N,h}\mathsf V_{h,C})(U)-O_C(U)\bigr)
 =\int\prod_C\bigl(\mathsf V_{h,C}(V)-O_C(U)\bigr)
       \mathcal K_{N,h}(U,dV).
\end{equation*}
Formula~\eqref{eq:E-multiple-density} applies to the fine insertions; at level~$h$ multiply the source density by the corresponding coarse factors. The probability kernel supplies no additional supremum factor. Conditional on the higher-mode propagation input identified in Lemma~\ref{lem:E-source-attachments}, weighted coefficient extraction, Lemma~\ref{lem:E-replication} and~\eqref{eq:E-vacuum-ratio} give, after recombination,
\begin{equation}\label{eq:E-bounded-error}
 \|\mathcal B'_{N,p}-O(c)\|_{\mathcal C^{\rm tile}_k}
 \le C_{k,m,d,c}\,
 e^{\sum_{i<h}\log C^{\rm col}_{J_{i+1}}}(1+C^{\rm col}_{J_h})
 e^{-(s_h-s_m)/8}.
\end{equation}
The prefactor has subpower growth in~$a_h^{-1}$ by~\eqref{eq:E-log-cost}; the last factor tends to zero. Choose $h$ so that this expression is at most~$\delta$, before choosing the cutoff. This also meets the stencil-fit condition.

For real~$c$, the original insertion is real by Lemmas~\ref{lem:A1} and~\ref{lem:E2.1a}. Set $\mathcal B=\operatorname{Re}\mathcal B'$. In the centered cell,
\begin{equation*}
 \mathcal B-O(c)=\tfrac12(\mathcal B'-O(c))+\tfrac12\overline{\mathcal B'-O(c)}.
\end{equation*}
The chessboard seminorm and its invariance under conjugation show that this does not increase~\eqref{eq:E-bounded-error}; the supremum bound also does not increase. Translation, signed permutation and conjugation use the same finite collection of stencils and bounds. This proves the lemma on the stated ordered tails.
\end{proof}

Fix a degree bound~$d$, a lower level~$m$ and an ordered extraction. A real grid word is a list $w=((d_i,c_i),x_i)_{i=1}^n$ with $d_i\le d$, $\bar c_i=c_i$ and distinct spacing-$a_k$ positions. We assume its configuration is generic. Let $\vartheta$ denote linear time reflection on positions and labels, without complex conjugation of scalar coefficients. For an ordered word $w$, write $w^*$ for conjugation and reversal of order, and $\theta w=\vartheta w^*$ for the antilinear reflection of Osterwalder and Schrader. Put
\begin{equation*}
 F_{N,p}(w)=\prod_iO^{(d_i)}_{N,p}(c_i,x_i/a_N),\qquad
 z_j(T)=Tj a_k(1,1,1,1),
\end{equation*}
\begin{equation*}
 f_{j,N,p,T}=F_{N,p}(\tau_{z_j(T)}w)
             -F_{N,p}(\tau_{z_j(T)}\vartheta w).
\end{equation*}
These are the real word differences used in the off-diagonal variance estimate in Section~\ref{sec:partII-H}.
\begin{lemma}[Bounded real word approximations]\label{lem:E2-bounded-real}
There is $k\ge m$, chosen before~$w$, such that every such word has, for each $\varepsilon>0$, measurable real approximations $g_{j,N,p,T}$ and $R_\varepsilon<\infty$ with $|g_{j,N,p,T}|\le R_\varepsilon$ for all $N,p,T,j$. For each $i<j$ and every sufficiently large integer~$T$, eventually in the ordered limits,
\begin{equation*}
 \left|\langle g_{i,N,p,T}g_{j,N,p,T}\rangle_{N,p}
 -\operatorname{Re}\langle f_{i,N,p,T}f_{j,N,p,T}\rangle_{N,p}\right|\le\varepsilon.
\end{equation*}
\end{lemma}
\begin{proof}
Each insertion of the word is replaced by its bounded approximation from Lemma~\ref{lem:E-coarse-bounded}; the chessboard inequality~\eqref{eq:E-chess}, with the seminorm bounds of Lemmas~\ref{lem:E2.3} and~\ref{lem:F1a}, controls the replacement error by seminorms alone. Choose
$k\ge\max(m,\max_{e\le d}\reflevel{e})$. It depends only on~$d,m$, before the length, labels and positions of~$w$. All its finer comparison levels remain available by Lemma~\ref{lem:E1.1}.

Include every label in~$w$ and its signed and conjugated images in a finite list. Lemma~\ref{lem:E2.3} gives a bound $C$ for their centered seminorms at any one fixed sufficiently fine level. Choose that level so that its cells and the margins required by Lemma~\ref{lem:F1a} fit between distinct spacing-$a_k$ points. Apply Lemma~\ref{lem:E-coarse-bounded} with error~$\delta$ at a still finer centered level, preserving these margins. Denote its real bounded approximations by~$\mathcal B_i$. The same family, translated to any anchor, is used for every copy of the word. Its seminorms are at most $C+\delta$, independently of its possibly much larger supremum bound.

For a product of $n_{\rm word}$ such insertions at distinct grid sites, telescope the products, replacing one factor at a time. Put all $n_{\rm word}$ factors into distinct cells with the chosen margins. The interior-seminorm comparison and~\eqref{eq:E-chess} give
\begin{equation}\label{eq:E-word-error}
 \left|\left\langle\prod_{a=1}^{n_{\rm word}} \mathcal B_a-\prod_{a=1}^{n_{\rm word}} O_a\right\rangle_{N,p}\right|
 \le n_{\rm word}\delta(C+\delta)^{n_{\rm word}-1}.
\end{equation}
Only the seminorms, not the supremum bounds on~$\mathcal B_a$, occur on the right. For each fixed finite configuration, sufficiently large periods avoid wrapping identifications; its stencils then lie in the prescribed cells on the inner cutoff tail. This explains the order of the two eventualities.

Let $G_{j,N,p,T}$ be the product of the $\mathcal B_i$ for the translated word, and let $G^\theta_{j,N,p,T}$ be the corresponding product for its reflected word. Set $g_j=G_j-G^\theta_j$. Choose $M_i^{\rm bnd}$ to bound both the original and reflected approximations in slot~$i$. Then
\begin{equation*}
 |g_j|\le 2\prod_iM_i^{\rm bnd}=:R_\varepsilon,
\end{equation*}
independently of~$N,p,T,j$. Every~$g_j$ is real and measurable. Set $g_j=0$ outside the ordered tails on which the coarse approximations are defined, so this supremum bound holds for all $N,p,T,j$. For fixed $i<j$, sufficiently large $T$ separates the two translated finite lists, for each of the four combinations of a word and a reflected word. Indeed their first coordinates differ by $T(j-i)a_k$ plus a fixed finite difference. Each list is itself distinct because the original word is generic and reflection is injective.

Expand $g_ig_j$ and $f_if_j$ into their four product terms and apply~\eqref{eq:E-word-error} with $n_{\rm word}=2n$. Choose $0<\delta\le1$ so that $4(2n)\delta(C+1)^{2n-1}\le\varepsilon$. Taking real parts cannot increase the error. Intersecting the four eventual sets proves the assertion. The empty word has zero difference and may be assigned $g_j=0$. There is no common cutoff threshold in $T,i,j$; a finite list of pairs is handled by intersecting its eventual sets. The choice gives such a grid level above every prescribed lower level~$m$, which is the cofinality required in Section~\ref{sec:partII-H}.
\end{proof}


\subsection{Full ultraviolet convergence}\label{sec:partII-E.3}

We prove convergence as the ultraviolet cutoff increases, at fixed volume. The argument is a triangle inequality through a common fixed level~$m$. At a separated configuration~$x$ on a fixed common grid,
\begin{equation}\label{eq:E-three-term}
|M_{N,p}-M_{N',p}|\le
\underbrace{|M_{N,p}-C_{N,p,m}|}_{\text{stopping error}}
+\underbrace{|C_{N,p,m}-C_{N',p,m}|}_{\text{fixed-level difference}}
+\underbrace{|C_{N',p,m}-M_{N',p}|}_{\text{stopping error}},
\end{equation}
where $M_{N,p}=M_{N,p}(v;x/a_N)$ is the cutoff moment and $C_{N,p,m}=C_{N,p,m}(v;x)$ is its approximation computed at level~$m$ from the unmarked representation and the one-mark families~\eqref{eq:E-coarse-integral}. The intended bounds are $2F_0e^{-B_m/64}$ for each stopping error, with $F_0$ independent of~$m$, and convergence of the fixed-level difference as $N,N'\to\infty$ at fixed~$m$. The proofs of Lemmas~\ref{lem:E-stopping} and~\ref{lem:E-fixed-level} identify estimates still needed for these bounds. Once those estimates are established, choose $m$ to make both stopping errors small, hold it fixed, and then send both cutoffs to infinity.

The middle term compares two cutoff histories at the same level~$m$. This differs from the contraction of the marked step, which compares two coefficient graphs carried by one unmarked history. Table~\ref{tab:E-fixed-level} lists the discrepancies in the middle term and the estimate that controls each.

\Needspace{18\baselineskip}
\begin{table}[!h]
\centering\small
\begin{tabular}{@{}p{0.4\textwidth}p{0.54\textwidth}@{}}
\toprule
\raggedright Discrepancy between cutoffs $N,N'$ & \raggedright Controlled by\tabularnewline
\midrule
\raggedright unmarked couplings, activities and paths at a fixed level & \raggedright convergence along the full cutoff sequence (Lemma~\ref{lem:D-fixed-level})\tabularnewline
\raggedright two marked graphs on one unmarked history & \raggedright one marked step halves their distance (Lemma~\ref{lem:E-marked-step}(iii))\tabularnewline
\raggedright marked steps on the two finite-cutoff histories & \raggedright a common graph class and fixed-input convergence on a dense core are needed; see Lemmas~\ref{lem:E-common-core}--\ref{lem:E-fixed-level}\tabularnewline
\raggedright scalar normalization $e^{v_{{\rm sc},m}|\Lambda_{n_{p,m}}|}$ & \raggedright not compared: it cancels in the normalized ratios~\eqref{eq:E-normalized-density} and~\eqref{eq:E-coarse-scalar-free}\tabularnewline
\raggedright scalar-free partition integral & \raggedright Cauchy by unmarked convergence; bounded below on a cutoff tail by a positive product Haar ball on all links\tabularnewline
\bottomrule
\end{tabular}
\caption{The discrepancies in the fixed-level difference $C_{N,p,m}-C_{N',p,m}$.}\label{tab:E-fixed-level}
\end{table}

The third row compares the actual finite-cutoff steps. Lemma~\ref{lem:D-fixed-level} gives convergence of their unmarked coordinates in the stated weighted spaces; that convergence alone does not construct the common marked graph class or prove convergence on its core. The required fixed-input bounds include the original eligibility tests and output tails, as described in Lemma~\ref{lem:E-fixed-level}. No marked step on a limiting completed history is assumed. Fix~$p$. In the normalized density at a fixed level, the scalar normalization appears as the same field-independent factor in numerator and denominator:
\begin{equation}\label{eq:E-normalized-density}
\frac{\density\rho_{N,p,m}}{\int\density\rho_{N,p,m}}
=\frac{e^{v_{{\rm sc},m}|\Lambda_{n_{p,m}}|}\,e^{-\mathcal V^\flat_{N,p,m}}\mathcal F_{N,p,m}}
{e^{v_{{\rm sc},m}|\Lambda_{n_{p,m}}|}\int e^{-\mathcal V^\flat_{N,p,m}}\mathcal F_{N,p,m}}
=\frac{e^{-\mathcal V^\flat_{N,p,m}}\mathcal F_{N,p,m}}
{\int e^{-\mathcal V^\flat_{N,p,m}}\mathcal F_{N,p,m}},
\end{equation}
where $\mathcal V^\flat_{N,p,m}=\pathaction_{N,m}+\sum_i\mathbf1_{X_i\,\mathrm{healthy}}E_i^{\rm reg}=\mathcal V_{\rho_{N,p,m}}+v_{{\rm sc},m}|\Lambda_{n_{p,m}}|$ is the scalar-free part of the potential, and $\mathcal F_{N,p,m}$ is the forced sum, both as in Definition~\ref{def:D06}. The factor $e^{v_{{\rm sc},m}|\Lambda_{n_{p,m}}|}$ depends on the cutoff, but it has canceled, so no convergence of that factor is proved or used.

For a fixed finite insertion-label list and a separated common-grid configuration~$x$, let $o_{i,m}=x_i/a_m$ be its sites at all sufficiently fine levels. They are read modulo the period; their parents agree at consecutive levels. In~\eqref{eq:E-source-form}, form $\mathcal H^0_{N,p,m}$ by setting $H_{m,X}=K^{\ge2}_{m,Y}=0$, keeping the complete singleton functions $f_{i,m}$ and~$k_{i,m,Y}$. Define
\begin{equation}\label{eq:E-coarse-integral}
 C_{N,p,m}(v;x)=
 \frac{\int[z_{\{1,\ldots,n\}}]\mathcal H^0_{N,p,m}(z;U)\,dU}
      {\int\density\rho_{N,p,m}\,dU}.
\end{equation}
The numerator is specified by the unmarked representation and the one-mark families alone. Both integrands contain the factor $e^{v_{{\rm sc},m}|\Lambda_{n_{p,m}}|}$, the numerator through~$\Phi_m$, and it cancels as in~\eqref{eq:E-normalized-density}:
\begin{equation}\label{eq:E-coarse-scalar-free}
\begin{aligned}
 C_{N,p,m}(v;x)&=
 \frac{e^{v_{{\rm sc},m}|\Lambda_{n_{p,m}}|}\int[z_{\{1,\ldots,n\}}]e^{-v_{{\rm sc},m}|\Lambda_{n_{p,m}}|}\mathcal H^0_{N,p,m}\,dU}
      {e^{v_{{\rm sc},m}|\Lambda_{n_{p,m}}|}\int e^{-v_{{\rm sc},m}|\Lambda_{n_{p,m}}|}\density\rho_{N,p,m}\,dU}
 \\
 &=\frac{\int[z_{\{1,\ldots,n\}}]e^{-v_{{\rm sc},m}|\Lambda_{n_{p,m}}|}\mathcal H^0_{N,p,m}\,dU}
      {\int e^{-v_{{\rm sc},m}|\Lambda_{n_{p,m}}|}\density\rho_{N,p,m}\,dU}.
\end{aligned}
\end{equation}
After the cancellation both integrands are scalar-free; they are formed from the scalar-free density and the singleton functions, whose convergence is proved below. The difference from the exact source is the stopping error, estimated in Lemma~\ref{lem:E-stopping}.

The first lemma treats one fixed coarse level. Its proof uses a comparison of the actual marked steps on a common class of graphs. The elementary observation below shows which convergence of those steps is sufficient.
\begin{lemma}[Contracting steps on a common core]\label{lem:E-common-core}
Let $X_i$ be nonempty complete metric spaces of diameter at most $D$, and let $T_{N,i}:X_{i+1}\to X_i$, defined for all sufficiently large~$N$, have a common Lipschitz constant $0\le q<1$. Suppose that, for each fixed~$i$, the sequence $T_{N,i}g$ is Cauchy for every $g$ in a dense subset of $X_{i+1}$. Suppose also that, for every fixed~$h$ and all sufficiently large~$N$, the selected trajectory is defined through level~$h$, lies in the respective spaces and obeys $g_{N,i}=T_{N,i}g_{N,i+1}$ at every $i<h$. Then $g_{N,m}$ is Cauchy for each fixed~$m$.
\end{lemma}
\begin{proof}
If $g_N\to g$ in $X_{i+1}$ and $a$ belongs to the dense subset, then
\[
 d_i(T_{N,i}g_N,T_{N',i}g_{N'})
 \le q\,d_{i+1}(g_N,a)+d_i(T_{N,i}a,T_{N',i}a)
       +q\,d_{i+1}(a,g_{N'}).
\]
First let $N,N'\to\infty$, then choose $a$ arbitrarily close to~$g$. Completeness gives a common limit. Fix $h>m$ and start each cutoff recursion at the same point of~$X_h$. Applying this argument at the finitely many levels from $h$ to~$m$ gives Cauchy auxiliary outputs. Each differs from the selected output by at most $Dq^{h-m}$. Given a tolerance, choose $h$ first and then choose the cutoffs, proving the claim.
\end{proof}

For a fixed graph, finite address truncations give a candidate core: the polynomial coordinate space has a finite $\ell^1$ basis, and the weighted remainder norm is an anchored sum over addresses. The tail of each basis column tends to zero; a single finite union controls all columns. Closing that union under the finite symmetry and conjugation orbits preserves those symmetries. The retained functions keep their local cancellations and original eligibility tests. This is approximation of each fixed graph, not a uniform tail estimate over the graph ball. Its use below requires a common address dictionary and admissibility of the truncated graphs for the actual steps.

The next lemma treats one fixed coarse level.
\begin{lemma}[Fixed-level marked convergence]\label{lem:E-fixed-level}
For every fixed $p,v,x$ as above and every sufficiently fine fixed level~$m$, the scalar-free unmarked density and the complete singleton functions in~\eqref{eq:E-source-form} are Cauchy in their common function spaces as $N\to\infty$. Their normalized integrals satisfy
\begin{equation}\label{eq:E-coarse-cauchy}
 C_{N,p,m}(v;x)\text{ is Cauchy as }N\to\infty.
\end{equation}
The scalar-free partition integral is bounded below by a positive constant depending on~$p,m$.
\end{lemma}
\begin{proof}
The proof describes the comparisons in Table~\ref{tab:E-fixed-level}, subject to the remaining verification identified below. Lemma~\ref{lem:D-fixed-level} gives full-sequence convergence of the scalar couplings and scalar-free activities in fixed weighted spaces at each fixed~$p,m$, and includes the prescribed paths with the strict reserve. Thus the cutoff-dependent path arrays converge in the same charts. The scalar normalization is excluded, as in~\eqref{eq:E-normalized-density}.

Use $\kappa_i=(\epsilon_i^{\rm for})^{1/2}$ in the remainder norm~\eqref{eq:E-marked-mass}. For coefficient graphs $g,\widetilde g$ at level~$i$ put
\begin{equation*}
 \operatorname{dist}^{\rm gr}_i(g,\widetilde g)=\sup_{\mathbf x\ne0}\frac{\|g(\mathbf x)-\widetilde g(\mathbf x)\|_{\rm rem}}{C^{\rm col}_{J_i}|\mathbf x|_1},
\end{equation*}
and let $\operatorname{dist}^{\rm gr}_i(N,N')$ be this distance between the graphs selected at cutoffs $N,N'$, after the fixed-level identification of their unmarked coordinates. The graph bound in the definition of a coefficient graph gives $\operatorname{dist}^{\rm gr}_i\le2$; by Lemma~\ref{lem:E-marked-step}(iii), one step on a common unmarked representation halves it. For different cutoffs the comparison is between actual finite-cutoff steps on common addresses. Lemma~\ref{lem:E-common-core} reduces the needed comparison to fixed inputs in a dense common core, provided the hypotheses of that lemma hold for these steps. No marked step on the limiting completed unmarked coordinates is needed.

\emph{Remaining verification.} Construct a common complete class on a fixed dictionary, with the original support and eligibility tests, on which all the actual steps and their half-contraction are defined. A truncated graph generally represents a different marked density from the selected one; its admissibility must be proved by applying the same marked construction to its own perturbation. In particular, adding an address with zero unmarked amplitude preserves the unmarked density, but a nonzero mark in that slot still contributes to its first variation. Its original eligibility and the corresponding bounds must be retained.

For each fixed finite-address input, establish convergence of the actual step in the full regular ordered-derivative and forced supremum norms, including the output tails and forward coefficient readout. The latter also needs a uniform same-cutoff Lipschitz bound in the input graph, as in the coefficient estimate of Lemma~\ref{lem:E-marked-step}; together with fixed-core convergence, it transfers convergence to the selected, cutoff-dependent graphs. A weighted kernel and derivative-score comparison in fixed physical links is one possible route. For the finite coefficient changes use $M^{-1}-\widetilde M^{-1}=M^{-1}(\widetilde M-M)\widetilde M^{-1}$ with common inverse bounds. These are statements about the actual finite-cutoff maps, not consequences of contraction alone. Subject to their verification, Lemma~\ref{lem:E-common-core} gives the Cauchy property at level~$m$. The finitely many coefficient maps between the reference level and~$m$ converge as well. Solving their fixed-dimensional linear systems then gives Cauchy polynomial coordinates for each fixed reference label. Substituting in the graph gives the complete singleton functions.

On the fixed finite torus, weighted convergence bounds the sums of the regular terms uniformly. For the forced gas, the product telescope in Lemma~\ref{lem:D6} bounds the difference while retaining each original eligibility test. Equivalently, incorporate those tests into the entries before grouping identical supports and inventories; there are only finitely many such supports and inventories. The resulting scalar-free density and the singleton functions are uniformly Cauchy on their fixed health domains. The health threshold at level~$m$ is fixed; no moving indicator is being differentiated or replaced. Formula~\eqref{eq:E-source-form}, with its higher modes set to zero, is a finite polynomial in the grouped functions and their exponentials. Its scalar-free coefficients are consequently uniformly Cauchy and integrable against Haar probability measure.

The denominator has a common positive lower bound. Choose $g>0$ below every admitted $\beta_{N,p,m}^{-1/2}$ on the retained fixed-level tail. Restrict every actual link to a fixed chord neighbourhood of the identity of radius below $\min\{t_m/4,g/4\}$, shrinking it once if necessary so that every plaquette has chord below $t_m$ and no flag is present. This set has positive product Haar measure. The forced sum there is one, and the scalar-free path action and activities are uniformly bounded above on the cutoff tail. The same positive lower bound for the integral therefore applies to every cutoff on that tail. Dividing the Cauchy numerator integrals by these positive Cauchy denominators gives~\eqref{eq:E-coarse-cauchy}, subject to the fixed-level convergence input above.
\end{proof}

The source density transfers exactly, but its singleton part does not: connected components with two or more marks remain. Two marks show what the stopping removes and what it keeps. Write $H_{m,X}=z_1z_2h_X$ and $K^{\ge2}_{m,Y}=z_1z_2\kappa_Y$, put $K^{\mathcal B}=\prod_{Y\in\mathcal B}K_{m,Y}$, and suppress $m$ and~$U$ on the singleton functions. Since $z_1^2=z_2^2=0$, expanding the exponential and the product in~\eqref{eq:E-source-form} gives
\begin{equation}\label{eq:E-two-mark}
\begin{aligned}
 [z_1z_2]\mathcal H_m={}&e^{\Phi_m}\sum_{\mathcal A}\Bigl[f_1f_2K^{\mathcal A}
 +\sum_{Y\in\mathcal A}(f_1k_{2,Y}+f_2k_{1,Y})K^{\mathcal A\setminus\{Y\}}
 +\sum_{Y\ne Y'}k_{1,Y}k_{2,Y'}K^{\mathcal A\setminus\{Y,Y'\}}\Bigr]
 \\
 &+e^{\Phi_m}\sum_{\mathcal A}\Bigl[\Bigl(\sum_Xh_X\Bigr)K^{\mathcal A}
 +\sum_{Y\in\mathcal A}\kappa_YK^{\mathcal A\setminus\{Y\}}\Bigr],
\end{aligned}
\end{equation}
with sums over eligible families and ordered pairs of distinct $Y,Y'\in\mathcal A$. The contribution in the first bracket is $[z_1z_2]\mathcal H^0_{N,p,m}$. It keeps the marked forced terms $k_{i,Y}$, which enter through the same eligible families as the unmarked activities. The contribution in the second bracket is removed, and its normalized integral is the stopping error for two marks. On a configuration without flags only the empty family is eligible, and the coefficient reduces to $e^{\Phi_m}(f_1f_2+\sum_Xh_X)$.

The next estimate bounds the effect of omitting all components with two or more marks. Its quantifiers are those needed in~\eqref{eq:E-three-term}, where $m$ is chosen before the cutoffs; so the constant~$F_0$ below must not depend on $m$ or~$N$. Convergence at each fixed level does not supply such a bound, since it concerns $C_{N,p,m}$ and not its distance from the cutoff moment.
\begin{lemma}[Stopping estimate]\label{lem:E-stopping}
For fixed $p,v,x$, there are $F_0<\infty$ and $m_0$ such that, for each $m\ge m_0$ and every sufficiently large $N$,
\begin{equation}\label{eq:E-stop-error}
 |M_{N,p}(v;x/a_N)-C_{N,p,m}(v;x)|
 \le2F_0e^{-B_m/64}.
\end{equation}
The constant $F_0$ is independent of $m,N$; the cutoff threshold may depend on~$m$.
\end{lemma}
Every mode with at least two marks has support covering number at least $B_m/64$, by Lemma~\ref{lem:E-cover}. Multiply all such modes, and only those, by a complex~$z$. The complete singleton functions, the original forced eligibility tests and the unmarked representation are unchanged. Lowering the support exponent by two absorbs $|z|<r_m=e^{B_m/64}$, so the normalized integral $F_m(z)$ is holomorphic on this disk, with $F_m(1)$ the cutoff moment and $F_m(0)=C_{N,p,m}$. The argument needs a bound by a constant $F_0$ at one base level~$k_0$ chosen before~$m$; the required propagation is identified in the proof. Given that bound, Cauchy's estimate on circles of radius $r_m/2$ centered on $[0,1]$ gives
\begin{equation*}
 |F_m(1)-F_m(0)|\le\frac{F_0}{r_m/2}=2F_0e^{-B_m/64}.
\end{equation*}
The error is the bound on the disk divided by half its radius; it is small because the radius grows with~$m$ while the bound does not.
\begin{proof}
The expansion is that of~\eqref{eq:E-source-form} for the exact source evolution~\eqref{eq:E-multiple-density}; the proposed base-level bound uses Lemma~\ref{lem:E-replication} and the finite one-mark coefficient and graph bounds below. Choose a base level~$k_0$ above the reference levels and the grid containing~$x$. At all sufficiently fine levels, distinct ancestral sites have separation at least $B_i^2$: their separation is a fixed positive physical number divided by~$a_i$, while $B_i^2$ grows only as a power of $\log\widehat\beta_i$. Their occupancy is one. Enlarge $k_0$ so these assertions and the entry conditions for the marked source estimates hold at every subsequent level.

Choose mark weights $\mathsf w_h>0$ at this base level so the weighted complete singleton sizes and coefficient norms lie below their prescribed bounds. The reference coefficient of each primary insertion is prescribed in Definition~\ref{def:D08}. Iterating the inverse coefficient estimate of Lemma~\ref{lem:E-marked-step}(i) over the finitely many steps from that mark's reference level to $k_0$ bounds $b_{h,k_0}=a_{k_0}^{-d_h}|\mathbf x_{h,k_0}|_1$ uniformly in~$N$. The graph bound then controls the full singleton norms. Choose the fixed weights small enough that $\mathsf w_hb_{h,k_0}\le\zeta_{d_h,k_0}$ as well as the stated singleton bounds. This choice does not use cross-cutoff graph convergence. Let $s_0^{\rm reg},s_0^{\rm for}$ bound the weighted regular singleton supremum and the sum of forced singleton suprema at~$k_0$. Set
\begin{equation}\label{eq:E-stopping-constant}
 F_0=\frac{\exp\{|\Lambda_{n_{p,k_0}}|q_{k_0}
                      +n(s_0^{\rm reg}+s_0^{\rm for}+2)\}}{\prod_h\mathsf w_h}.
\end{equation}
The numerator is the normalized-integral bound of the several-mark expansion, now with these fixed base weights. It depends on the base level, not on the later stopping level.

At level~$m>k_0$, every nonzero mode of $H_{m,X}$ or $K^{\ge2}_{m,Y}$ contains at least two marks and its support contains their anchors. The covering estimate of Lemma~\ref{lem:E-cover} gives
\begin{equation*}
 \cover_{B_m}(X)\ge B_m/64.
\end{equation*}
Multiply every one of these higher-mode functions by a common complex variable~$z$. For $|z|\le e^{B_m/64}$, its amplitude increase is absorbed by decreasing the support exponent by two:
\begin{equation*}
 |z|e^{(\sigma-2)\cover_{B_m}(X)}
 \le e^{\sigma\cover_{B_m}(X)}.
\end{equation*}
Singleton functions are not multiplied by~$z$. The source representation therefore remains integrable on $|z|<e^{B_m/64}$, with its prescribed singleton functions, original forced eligibility tests and two fewer units of support exponent. Here $\sigma=\sigma_{J_m}=100$. For $s\ge0$ and $98\le\tau\le100$, $e^{\tau s}\le e^{100s}$, so every regular or forced weighted sum decreases when its exponent decreases. The coefficientwise rooted clause of Lemma~\ref{lem:D-connected} bounds the connected regular terms at exponent $\sigma+3$ and the forced terms at $\sigma+2$. The entry chosen in item~\ref{it:E-entries} controls the marked product factors. Applying these bounds in the weighted source norm of~\eqref{eq:E-source-form} leaves the two units used here, with final exponent $98$.

To carry this modified density back to the same base level~$k_0$, one needs the coefficientwise marked attachment induction restarted at level $m$, with the actual modified higher modes admitted at exponent $98$ and with weights satisfying the base graph and coefficient hypotheses. Its anchored higher-mode bound and this restart are not yet established in this draft. Multiplication by $z$ preserves the original eligibility tests and cannot lower marker degree; it does not by itself prove the required bound. If that induction supplies the stated base estimates, the singleton exponential and forced factors contribute at most $n(s_0^{\rm reg}+s_0^{\rm for})$ to the logarithmic majorant, and the remaining connected regular and forced sums contribute at most $2n$. Equation~\ref{eq:E-vacuum-ratio} contributes $|\Lambda_{n_{p,k_0}}|q_{k_0}$. Coefficient extraction divides by $\prod_i\mathsf w_i$, giving~\eqref{eq:E-stopping-constant}. The resulting bound on the normalized full-mode integral $F_m(z)$ is
\begin{equation*}
 |F_m(z)|\le F_0.
\end{equation*}
This avoids the level-$m$ volume factor that would arise from estimating the normalized integral only where it was stopped.

For $e^{B_m/64}>2$, the Cauchy derivative estimate on circles of radius $e^{B_m/64}/2$ centered on $[0,1]$ gives
\begin{equation*}
 |F_m(1)-F_m(0)|\le2F_0e^{-B_m/64}.
\end{equation*}
At $z=1$ this is the exact multiple-source integral. Its equality to the original cutoff moment follows from Fubini and the probability normalization of every block kernel in~\eqref{eq:E-multiple-density}. At $z=0$ it is exactly~\eqref{eq:E-coarse-integral}. Substitution proves the estimate. For each fixed~$m$, the singleton bounds use the finite one-mark estimates and the propagated weight relation. The cutoff threshold may depend on~$m$, while $F_0$ stays fixed.
\end{proof}


The fixed-level, stopping and positional estimates would give the full limit once the remaining inputs to those estimates are established.
\begin{lemma}[Lattice-point convergence]\label{lem:E3.1}
For every $p,n,v$, there is $f_{p,v}\in C^\infty(\NoHit{L_p}{n})$ such that, for every $\boldsymbol l$, compact $K\subset\NoHit{L_p}{n}$ and $\varepsilon>0$, some $N_0(p,n,v,\boldsymbol l,K,\varepsilon)$ gives
\begin{equation*}
|\Delta_{\boldsymbol l}M_{N,p}(v;x/a_N)-\partial_{\boldsymbol l}f_{p,v}(x)|<\varepsilon
\end{equation*}
at every spacing-$a_N$ point of $K$ in the periodic lift, for every $N\ge N_0$.
\end{lemma}
\begin{proof}
Pointwise convergence on the common grids is~\eqref{eq:E-three-term}; the difference bounds of Lemma~\ref{lem:E2.4} extend it to uniform convergence of all derivatives. By Lemmas~\ref{lem:E-fixed-level}--\ref{lem:E-stopping}, the right side of~\eqref{eq:E-three-term} is at most $4F_0e^{-B_m/64}+|C_{N,p,m}-C_{N',p,m}|$; taking first $m$ and then $N,N'$ large proves the Cauchy property at every separated configuration on a fixed common grid. Since every ratio is an integer, each coarser grid is contained in every finer grid. The union of these grids is dense in the periodic noncollision domain.

We next extend this limit and identify its derivatives. Choose a compact neighbourhood $K'$ of~$K$ contained in that domain, and a positive lower bound for its torus separation. Lemma~\ref{lem:E2.4} bounds the lattice functions and all their fixed differences there. One more difference gives, by telescoping along coordinate lattice paths,
\begin{equation*}
|G_N(z)-G_N(z')|\le B\,|z-z'|_1
\end{equation*}
for nearby physical grid points, where $G_N=\Delta_{\boldsymbol l}M_{N,p}$. For order zero, a finite dense grid in $K'$ and this estimate extend the pointwise Cauchy property to uniform convergence on~$K$. The multilinear interpolants obey the same bound, with an additional $O(a_N)$ displacement to their vertices. The resulting limit~$f_{p,v}$ is continuous and periodic.

Suppose $G_N$ already converges locally uniformly to~$g$. Bounds for two more differences compare its first forward difference with a difference over a fixed physical interval. Put $h_N=Q_Na_N\to h>0$, with $Q_N\in\N$. Summing successive first differences and bounding their variation by the second difference gives
\begin{equation}\label{eq:E-discrete-taylor}
\left|\frac{G_N(z+h_Ne)-G_N(z)}{h_N}-\Delta_eG_N(z)\right|
\le\tfrac12B h_N.
\end{equation}
Choose $h$ small and then let $N\to\infty$. The endpoint values converge uniformly, so~\eqref{eq:E-discrete-taylor} proves that $\Delta_eG_N$ is uniformly Cauchy on~$K$. Call its continuous limit~$g_e$. The exact telescoping identity on grid segments, followed by Riemann summation of the uniformly convergent first differences, gives
\begin{equation*}
g(z+te)-g(z)=\int_0^t g_e(z+ue)\,du.
\end{equation*}
Thus $\partial_eg=g_e$. Induct on the length of the difference list. Every required additional order is available from Lemma~\ref{lem:E2.4}, so all derivative limits exist and $f_{p,v}$ is smooth. Finite covers by such neighbourhoods give the asserted uniform statement on~$K$.
\end{proof}

The same limit applies when observation points move across periodic representative boundaries. Smearing then follows from a compact-support Riemann sum.
\begin{lemma}[Interpolation and smearing]\label{lem:E3.2}
For fixed $p,n,v,\boldsymbol l$ and $x_N\to x\in\NoHit{L_p}{n}$, the multilinear interpolants of $\Delta_{\boldsymbol l}M_{N,p}$ at~$x_N$ converge to $\partial_{\boldsymbol l}f_{p,v}(x)$. For every $p$-admissible separated test~$t$,
\begin{equation*}
t\partial_{\boldsymbol l}f_{p,v}\in L^1,\qquad
 a_N^{4n}\sum_y t(\hat y)\Delta_{\boldsymbol l}M_{N,p}(v;y)
\longrightarrow\int t\partial_{\boldsymbol l}f_{p,v}.
\end{equation*}
\end{lemma}
\begin{proof}
Both assertions follow from the locally uniform convergence of Lemma~\ref{lem:E3.1}, with the placement and moment conventions of Definition~\ref{def:D03} and admissibility from Section~\ref{sec:continuum}. Choose a compact neighbourhood of~$x$ inside the periodic noncollision domain. Every vertex used by the interpolant at~$x_N$ lies there eventually and tends to~$x$. The interpolant is a convex combination of the vertex values. Uniform convergence at these vertices and continuity of the limiting derivative give the moving-point assertion. Work on the periodic lift throughout; no continuity of the choice of centered representatives is needed.

The closed support of~$t$ is a compact subset of that domain, strictly inside the chosen period box. The smooth limiting derivative is bounded there, so its product with~$t$ is integrable. On a fixed compact neighbourhood, Lemma~\ref{lem:E3.1} makes the supremum of the difference between the lattice values and the limiting derivative tend to zero. Its contribution to the sum is at most that supremum times $\|t\|_\infty$ times the total volume of the grid cubes meeting the support. This volume is bounded independently of~$N$. The remaining sum is the Riemann sum for the continuous compactly supported function $t\partial_{\boldsymbol l}f_{p,v}$ and converges to its integral.
\end{proof}

\begin{proof}[Proof of Theorem~\ref{prop:UV}]
Set $M^{\rm UV}_{p,v}=f_{p,v}$. Lemma~\ref{lem:E3.1} gives its smoothness on the full periodic noncollision domain, and Lemma~\ref{lem:E3.2} gives the asserted interpolation and smearing limits for every fixed derivative list. Thus every inner limit in Definition~\ref{def:D09} exists at every admissible fixed volume along the full cutoff sequence.
\end{proof}
